% !TEX program = pdflatex
\documentclass[12pt,a4paper]{article}
\usepackage{ctex}
\usepackage{amsmath,amssymb,bm}
\usepackage{booktabs,longtable,array,multirow}
\usepackage[a4paper,margin=2.5cm]{geometry}
\emergencystretch=3em % 允许长串（arXiv 编号等）换行时拉伸，避免中英混排溢出
\usepackage{xcolor}
\usepackage{fancyhdr}
\usepackage{titlesec}
\usepackage{tcolorbox}
\tcbuselibrary{skins,breakable}
\usepackage{enumitem}
\usepackage[colorlinks=true,linkcolor=blue,citecolor=blue]{hyperref}
\usepackage{listings}
\usepackage{tocloft}
\usepackage{float}
\usepackage{tikz}
\usetikzlibrary{arrows.meta,positioning,shapes.geometric,fit,calc,backgrounds,decorations.pathreplacing}
\usepackage{pgfplots}
\pgfplotsset{compat=1.17}
% ---------- TikZ 统一样式 ----------
\tikzset{
  blk/.style={draw=cblue, fill=cblue!8, rounded corners=2pt, minimum height=8mm, align=center, font=\small, inner sep=4pt},
  grn/.style={draw=cgreen!70!black, fill=cgreen!8, rounded corners=2pt, minimum height=8mm, align=center, font=\small, inner sep=4pt},
  gry/.style={draw=cgray, fill=cgray!8, rounded corners=2pt, minimum height=8mm, align=center, font=\small, inner sep=4pt},
  orn/.style={draw=orange!80!black, fill=orange!12, rounded corners=2pt, minimum height=8mm, align=center, font=\small, inner sep=4pt},
  dgt/.style={draw=cblue!60, fill=cblue!4, rounded corners=2pt, minimum height=8mm, align=center, font=\small, inner sep=4pt},
  arr/.style={-{Stealth[length=2.6mm]}, thick, draw=cgray},
  brr/.style={-{Stealth[length=2.6mm]}, thick, draw=cblue},
  rrr/.style={-{Stealth[length=2.6mm]}, thick, draw=orange!80!black, dashed},
  lab/.style={font=\scriptsize\itshape, text=cgray},
  lay/.style={draw=cblue!70, fill=cblue!6, rounded corners=3pt, minimum height=11mm, align=center, font=\small, inner sep=5pt},
}
\definecolor{cblue}{RGB}{41,98,164}
\definecolor{cgreen}{RGB}{46,125,50}
\definecolor{cgray}{RGB}{97,97,97}
\definecolor{codebg}{RGB}{245,246,248}
\titleformat{\section}{\Large\bfseries\color{cblue}}{\thesection}{1em}{}
\titleformat{\subsection}{\large\bfseries\color{cgreen}}{\thesubsection}{1em}{}
\titleformat{\subsubsection}{\normalsize\bfseries\color{cgray}}{\thesubsubsection}{1em}{}
\pagestyle{fancy}
\fancyhf{}
\fancyhead[L]{\small\itshape 深度学习与大语言模型：前沿技术全景报告}
\fancyhead[R]{\small 2026年9月}
\fancyfoot[C]{\thepage}
\renewcommand{\headrulewidth}{0.4pt}
\setlength{\headheight}{14pt}
\newtcolorbox{infobox}[1]{breakable,colback=cblue!4,colframe=cblue,fonttitle=\bfseries,title={#1},boxrule=0.8pt,left=8pt,right=8pt,top=6pt,bottom=6pt}
\newtcolorbox{keybox}[1]{breakable,colback=cgreen!4,colframe=cgreen,fonttitle=\bfseries,title={#1},boxrule=0.8pt,left=8pt,right=8pt,top=6pt,bottom=6pt}
\newtcolorbox{warnbox}[1]{breakable,colback=orange!8,colframe=orange!70!black,fonttitle=\bfseries,title={#1},boxrule=0.8pt,left=8pt,right=8pt,top=6pt,bottom=6pt}
\lstset{basicstyle=\ttfamily\small,backgroundcolor=\color{codebg},frame=single,numbers=left,numberstyle=\tiny\color{gray},breaklines=true,tabsize=2}
\setlist{topsep=4pt,itemsep=2pt}
% ---- 浮动体自适应参数：图表可漂浮补页，消除 [H] 造成的页尾大片留白 ----
\renewcommand{\topfraction}{0.9}       % 页顶浮动体最多占版心比例
\renewcommand{\bottomfraction}{0.6}    % 页底浮动体最多占版心比例
\renewcommand{\textfraction}{0.07}     % 每页至少保留的正文比例
\renewcommand{\floatpagefraction}{0.8} % 浮动页至少需填满的比例
\setcounter{topnumber}{3}
\setcounter{bottomnumber}{2}
\setcounter{totalnumber}{5}
\begin{document}
% ===================== 封面 =====================
\begin{titlepage}
\centering
\vspace*{2cm}
{\Huge\bfseries\color{cblue} 深度学习与大语言模型}\\[0.8cm]
{\Large 前沿技术全景报告}\\[1.5cm]
\rule{0.5\textwidth}{0.4pt}\\[1cm]
{\large 版本：2026 年 9 月}\\[0.3cm]
{\large 定位：系统性技术参考}\\[0.3cm]
{\large 覆盖：从基础数学到 2026 年前沿架构}\\[3cm]
{\small\itshape 本报告的章节结构：12 个主章节 + 术语附录，涵盖深度学习数学基础、Transformer 架构、LLM 训练流水线、MoE/GQA/MLA 前沿架构、推理优化、推理模型、多模态、安全对齐、可扩展性定律及 2025--2026 行业动态。}
\end{titlepage}
% ===================== 目录 =====================
\tableofcontents
\newpage
% ===================== 摘要与导读 =====================
\section*{摘要与导读}

本报告系统梳理了从深度学习数学基础到大语言模型（LLM）前沿技术的完整技术栈，回答三个核心问题：\textbf{大模型如何被构建、如何被训练、如何被对齐与治理}。全文按“基础原理 $\to$ 架构演进 $\to$ 训练系统 $\to$ 推理系统 $\to$ 能力扩展 $\to$ 安全治理 $\to$ 产业动态”的逻辑链条组织：第 1--2 章奠定数学基础并回顾架构演进史；第 3--4 章深入剖析 Transformer 机制与 LLM 训练流水线；第 5--6 章覆盖 MoE、GQA/MLA 等前沿架构与推理优化体系；第 7--8 章讨论推理模型（Reasoning Models）与多模态大模型；第 9--10 章面向安全对齐、可扩展性定律与长期趋势，其中 10.3 节引入\textbf{能力收敛假说}（CCH）——“表征收敛由规模免费赠予，能力收敛必须由访问结构购买”——作为贯通架构（第 5 章）、推理系统（第 6--7 章）与趋势判断（第 10、12 章）的理论主线；第 11--12 章汇总 2025--2026 年行业动态并给出总结与展望。

\textbf{阅读建议}（按读者画像）：
\begin{itemize}
  \item \textbf{初次接触者}：按顺序通读第 1$\to$3 章建立直觉，再按兴趣深入后续章节；
  \item \textbf{训练/推理工程师}：重点阅读第 4、6 章的分布式系统与推理优化细节；
  \item \textbf{算法研究者}：重点关注第 5、7、10 章的前沿架构、推理训练与 Scaling Law 分析；
  \item \textbf{技术与产品决策者}：重点阅读本摘要、第 10--11 章趋势判断与第 12 章展望。
\end{itemize}

\begin{infobox}{方法与数据来源说明}
本报告以公开论文、技术报告与官方文档为主要信息源。关键定量结论（如 DeepSeek-V3 的 671B 总参数 / 37B 激活参数 / 2.79M H800 GPU$\cdot$hours 训练成本、Qwen3-235B-A22B 的 128 专家 / 激活 8、LLaMA-2 70B 的 GQA 实验等）均引用自原始技术报告或论文。行业动态截至 2026 年 9 月：2025 年的主要事件（GPT-5、Gemini 3、Claude Opus 4.5、Qwen3、Kimi K2、Sora 2 等）与 NVIDIA 收购 Hugging Face（约 129 亿美元，2026 年 8 月官宣）均经多个公开报道交叉核实；2026 年上半年个别事件（如 Gemini 3.7 等）的量化细节以最终官方披露为准，读者引用时建议回溯原始来源核实。
\end{infobox}

\begin{figure}[htbp]
\centering
% fig01_macro_architecture.tex — 图1 报告宏观架构
% 《深度学习与大语言模型：前沿技术全景报告》图示库
% 由 dl_llm_report.tex 抽取，经 % fig01_macro_architecture.tex — 图1 报告宏观架构
% 《深度学习与大语言模型：前沿技术全景报告》图示库
% 由 dl_llm_report.tex 抽取，经 % fig01_macro_architecture.tex — 图1 报告宏观架构
% 《深度学习与大语言模型：前沿技术全景报告》图示库
% 由 dl_llm_report.tex 抽取，经 \input{figs/fig01_macro_architecture} 引用（caption/label 在主文件）
% 注意：本文件为片段，依赖主文件导言区的 tikz/pgfplots 宏包、颜色定义与 tikzset 样式，不可单独编译。

\begin{tikzpicture}[node distance=4mm,
  layerbox/.style={lay, text width=106mm, minimum height=10mm, align=center}]
\node[layerbox] (top) {\textbf{生态与展望}（第 11--12 章）\\ 行业动态 · 竞争格局 · 总结与趋势判断};
\node[layerbox, below=of top] (gov) {\textbf{规律与治理}（第 9--10 章）\\ 安全对齐与防御纵深 · Scaling Laws · 能力收敛假说（CCH）};
\node[layerbox, below=of gov] (cap) {\textbf{能力扩展}（第 7--8 章）\\ 推理模型（System 2） · 多模态大模型};
\node[layerbox, below=of cap] (sys) {\textbf{模型与系统}（第 3--6 章）\\ Transformer · 训练流水线 · MoE/GQA/MLA · 推理优化};
\node[layerbox, below=of sys] (base) {\textbf{数理与架构基础}（第 1--2 章）\\ 数学与优化 · CNN/RNN $\to$ Transformer};
\draw[brr] ($(base.west)+(-8mm,0)$) -- ($(top.west)+(-8mm,0)$);
\node[lab, rotate=90, anchor=south] at ($(base.west)!0.5!(top.west)+(-8mm,0)$) {自底向上 · 逐层依赖};
\end{tikzpicture}
 引用（caption/label 在主文件）
% 注意：本文件为片段，依赖主文件导言区的 tikz/pgfplots 宏包、颜色定义与 tikzset 样式，不可单独编译。

\begin{tikzpicture}[node distance=4mm,
  layerbox/.style={lay, text width=106mm, minimum height=10mm, align=center}]
\node[layerbox] (top) {\textbf{生态与展望}（第 11--12 章）\\ 行业动态 · 竞争格局 · 总结与趋势判断};
\node[layerbox, below=of top] (gov) {\textbf{规律与治理}（第 9--10 章）\\ 安全对齐与防御纵深 · Scaling Laws · 能力收敛假说（CCH）};
\node[layerbox, below=of gov] (cap) {\textbf{能力扩展}（第 7--8 章）\\ 推理模型（System 2） · 多模态大模型};
\node[layerbox, below=of cap] (sys) {\textbf{模型与系统}（第 3--6 章）\\ Transformer · 训练流水线 · MoE/GQA/MLA · 推理优化};
\node[layerbox, below=of sys] (base) {\textbf{数理与架构基础}（第 1--2 章）\\ 数学与优化 · CNN/RNN $\to$ Transformer};
\draw[brr] ($(base.west)+(-8mm,0)$) -- ($(top.west)+(-8mm,0)$);
\node[lab, rotate=90, anchor=south] at ($(base.west)!0.5!(top.west)+(-8mm,0)$) {自底向上 · 逐层依赖};
\end{tikzpicture}
 引用（caption/label 在主文件）
% 注意：本文件为片段，依赖主文件导言区的 tikz/pgfplots 宏包、颜色定义与 tikzset 样式，不可单独编译。

\begin{tikzpicture}[node distance=4mm,
  layerbox/.style={lay, text width=106mm, minimum height=10mm, align=center}]
\node[layerbox] (top) {\textbf{生态与展望}（第 11--12 章）\\ 行业动态 · 竞争格局 · 总结与趋势判断};
\node[layerbox, below=of top] (gov) {\textbf{规律与治理}（第 9--10 章）\\ 安全对齐与防御纵深 · Scaling Laws · 能力收敛假说（CCH）};
\node[layerbox, below=of gov] (cap) {\textbf{能力扩展}（第 7--8 章）\\ 推理模型（System 2） · 多模态大模型};
\node[layerbox, below=of cap] (sys) {\textbf{模型与系统}（第 3--6 章）\\ Transformer · 训练流水线 · MoE/GQA/MLA · 推理优化};
\node[layerbox, below=of sys] (base) {\textbf{数理与架构基础}（第 1--2 章）\\ 数学与优化 · CNN/RNN $\to$ Transformer};
\draw[brr] ($(base.west)+(-8mm,0)$) -- ($(top.west)+(-8mm,0)$);
\node[lab, rotate=90, anchor=south] at ($(base.west)!0.5!(top.west)+(-8mm,0)$) {自底向上 · 逐层依赖};
\end{tikzpicture}

\caption{报告宏观架构：五个层次自底向上堆叠，每层依赖下一层；括号内为对应章节。}
\label{fig:macro}
\end{figure}

图~\ref{fig:macro} 给出全文的宏观结构：\textbf{基础层}回答“深度学习为什么可行”，\textbf{模型与系统层}回答“大模型如何被构建、训练与高效部署”，\textbf{能力层}回答“如何让它更会推理、看见世界”，\textbf{规律与治理层}回答“边界在哪里、如何被约束”，\textbf{生态层}回答“谁在做什么、往哪里去”。各章末尾的小结框（infobox/keybox）均可独立阅读，作为该章的“电梯演讲”版本。

\newpage
% ===================== 第一章 =====================
\section{深度学习基础：从神经元到反向传播}
% ---------- 1.1 ----------
\subsection{神经元与计算图}
深度学习的原子单元是\textbf{人工神经元}。一个神经元接收输入向量 $\mathbf{x}\in\mathbb{R}^n$，计算加权和 $z=\mathbf{W}\cdot\mathbf{x}+b$，再通过\textbf{激活函数} $\sigma$ 得到输出 $a=\sigma(z)$。整个网络就是成千上万个这样的单元按特定拓扑连接而成的\textbf{计算图}（Computation Graph）。

现代 LLM 中，单个“神经元”的概念已弱化，取而代之的是矩阵运算和块状结构（block），但底层数学不变：
\[
\text{output} \;=\; \sigma\!\left(\sum_{i} W_i\, x_i + b\right)
\]
在 GPU 上，一次前向传播本质上是若干\textbf{大矩阵乘法}（GEMM）与逐元素激活的组合。这也是深度学习天然适合并行硬件的原因。

\begin{center}
\textbf{从神经元到矩阵：现代 LLM 的视角}\\[4pt]
$\text{单个神经元：} a = \sigma(\mathbf{w}^\top \mathbf{x} + b)$\\[2pt]
$\text{一层线性变换（n 个神经元）：} \mathbf{z} = \mathbf{W}\mathbf{x} + \mathbf{b}, \quad \mathbf{W}\in\mathbb{R}^{n\times d}$\\[2pt]
$\text{一个 LLM Block：} \mathbf{h} = \mathbf{x} + \text{Attn}(\text{RMSNorm}(\mathbf{x})) \;\longrightarrow\; \mathbf{h} + \text{FFN}(\text{RMSNorm}(\mathbf{h}))$
\end{center}

换言之，现代大语言模型中的“线性层”就是这样一个“神经元阵列”：以 LLaMA-2 7B 为例，每个注意力 QKV 投影矩阵约 $4096 \times 4096$，一个前馈 SwiGLU 层约有 3 个 $4096 \times 11008$ 的大矩阵——模型的“智能”正是分布在这数百个巨型矩阵的数值之中。

把视角再抬高一层的直观理解是：\textbf{学习就是在一个巨大的计算图上搜索“最优的旋钮组合”}。以 7B 模型为例，约 70 亿个参数就是 70 亿个旋钮；训练的过程，是让数万亿 token 依次流过计算图，用每次预测的误差（梯度）反向微调所有旋钮，直到整个网络对“下一个 token”的猜测足够准。这句话里藏着深度学习的三个支点：\emph{可微}（损失对参数处处可导，梯度存在）、\emph{可反向传播}（链式法则把误差按贡献归因到每个参数）、\emph{可并行}（大矩阵乘法主导，天然适配 GPU）。第 4--6 章的所有工程奇迹，本质上都是在放大这三个支点。图~\ref{fig:neuron} 给出了这一全景。

\begin{figure}[htbp]
\centering
% fig02_neuron_views.tex — 图2 深度学习的三个基本视图
% 《深度学习与大语言模型：前沿技术全景报告》图示库
% 由 dl_llm_report.tex 抽取，经 % fig02_neuron_views.tex — 图2 深度学习的三个基本视图
% 《深度学习与大语言模型：前沿技术全景报告》图示库
% 由 dl_llm_report.tex 抽取，经 % fig02_neuron_views.tex — 图2 深度学习的三个基本视图
% 《深度学习与大语言模型：前沿技术全景报告》图示库
% 由 dl_llm_report.tex 抽取，经 \input{figs/fig02_neuron_views} 引用（caption/label 在主文件）
% 注意：本文件为片段，依赖主文件导言区的 tikz/pgfplots 宏包、颜色定义与 tikzset 样式，不可单独编译。

\begin{tikzpicture}
% ---------- (a) 单神经元 ----------
\node[gry, minimum width=9mm] (x1) at (1.7, 0.85) {$x_1$};
\node[gry, minimum width=9mm] (x2) at (1.7, 0) {$x_2$};
\node[gry, minimum width=9mm] (x3) at (1.7, -0.85) {$x_3$};
\node[blk, align=center] (sum) at (3.9, 0) {加权和\\ $z=\mathbf{w}^{\top}\mathbf{x}+b$};
\node[grn, minimum width=9mm] (act) at (6.0, 0) {$\sigma$};
\node[gry, minimum width=9mm] (out) at (7.7, 0) {$a$};
\draw[arr] (x1.east) -- (sum.west);
\draw[arr] (x2.east) -- (sum.west);
\draw[arr] (x3.east) -- (sum.west);
\draw[arr] (sum) -- (act);
\draw[arr] (act) -- (out);
\node[lab] at (4.65, -1.55) {(a) 人工神经元：加权和 + 非线性};
% ---------- (b) 计算图与反向传播 ----------
\node[gry, minimum width=9mm] (bx) at (9.2, 0) {$\mathbf{x}$};
\node[blk, minimum width=13mm] (bl) at (10.8, 0) {线性层};
\node[grn, minimum width=11mm] (ba) at (12.5, 0) {激活};
\node[blk, minimum width=13mm] (bo) at (14.2, 0) {输出层};
\node[orn, minimum width=13mm] (bl2) at (15.9, 0) {损失 $L$};
\draw[arr] (bx)--(bl); \draw[arr] (bl)--(ba); \draw[arr] (ba)--(bo); \draw[arr] (bo)--(bl2);
\draw[rrr] (bl2.north) to[bend right=28] node[lab, above, pos=0.5] {反向：$\partial L/\partial\theta$ 沿链式法则逐层回传} (bl.north);
\node[lab] at (12.4, -1.55) {(b) 计算图：前向算值，反向归因};
% ---------- (c) 多层感知机（与上行整体居中） ----------
\node[lab] at (6.6, -2.7) {输入层};
\node[lab] at (8.9, -2.7) {隐藏层};
\node[lab] at (11.2, -2.7) {输出层};
\foreach \i in {1,2,3} \node[circle, draw=cgray, fill=cgray!8, minimum size=7mm] (mi\i) at (6.6, -3.35-0.95*\i) {};
\foreach \i in {1,2,3,4} \node[circle, draw=cblue, fill=cblue!8, minimum size=7mm] (mh\i) at (8.9, -2.9-0.95*\i) {};
\foreach \i in {1,2} \node[circle, draw=cgreen!70!black, fill=cgreen!8, minimum size=7mm] (mo\i) at (11.2, -3.35-1.4*\i) {};
\foreach \a in {1,2,3} \foreach \b in {1,2,3,4} \draw[cgray!50, thin] (mi\a) -- (mh\b);
\foreach \a in {1,2,3,4} \foreach \b in {1,2} \draw[cgray!50, thin] (mh\a) -- (mo\b);
\node[lab] at (8.9, -7.75) {(c) 多层感知机：层间全连接，“深度”来自层层复合};
\end{tikzpicture}
 引用（caption/label 在主文件）
% 注意：本文件为片段，依赖主文件导言区的 tikz/pgfplots 宏包、颜色定义与 tikzset 样式，不可单独编译。

\begin{tikzpicture}
% ---------- (a) 单神经元 ----------
\node[gry, minimum width=9mm] (x1) at (1.7, 0.85) {$x_1$};
\node[gry, minimum width=9mm] (x2) at (1.7, 0) {$x_2$};
\node[gry, minimum width=9mm] (x3) at (1.7, -0.85) {$x_3$};
\node[blk, align=center] (sum) at (3.9, 0) {加权和\\ $z=\mathbf{w}^{\top}\mathbf{x}+b$};
\node[grn, minimum width=9mm] (act) at (6.0, 0) {$\sigma$};
\node[gry, minimum width=9mm] (out) at (7.7, 0) {$a$};
\draw[arr] (x1.east) -- (sum.west);
\draw[arr] (x2.east) -- (sum.west);
\draw[arr] (x3.east) -- (sum.west);
\draw[arr] (sum) -- (act);
\draw[arr] (act) -- (out);
\node[lab] at (4.65, -1.55) {(a) 人工神经元：加权和 + 非线性};
% ---------- (b) 计算图与反向传播 ----------
\node[gry, minimum width=9mm] (bx) at (9.2, 0) {$\mathbf{x}$};
\node[blk, minimum width=13mm] (bl) at (10.8, 0) {线性层};
\node[grn, minimum width=11mm] (ba) at (12.5, 0) {激活};
\node[blk, minimum width=13mm] (bo) at (14.2, 0) {输出层};
\node[orn, minimum width=13mm] (bl2) at (15.9, 0) {损失 $L$};
\draw[arr] (bx)--(bl); \draw[arr] (bl)--(ba); \draw[arr] (ba)--(bo); \draw[arr] (bo)--(bl2);
\draw[rrr] (bl2.north) to[bend right=28] node[lab, above, pos=0.5] {反向：$\partial L/\partial\theta$ 沿链式法则逐层回传} (bl.north);
\node[lab] at (12.4, -1.55) {(b) 计算图：前向算值，反向归因};
% ---------- (c) 多层感知机（与上行整体居中） ----------
\node[lab] at (6.6, -2.7) {输入层};
\node[lab] at (8.9, -2.7) {隐藏层};
\node[lab] at (11.2, -2.7) {输出层};
\foreach \i in {1,2,3} \node[circle, draw=cgray, fill=cgray!8, minimum size=7mm] (mi\i) at (6.6, -3.35-0.95*\i) {};
\foreach \i in {1,2,3,4} \node[circle, draw=cblue, fill=cblue!8, minimum size=7mm] (mh\i) at (8.9, -2.9-0.95*\i) {};
\foreach \i in {1,2} \node[circle, draw=cgreen!70!black, fill=cgreen!8, minimum size=7mm] (mo\i) at (11.2, -3.35-1.4*\i) {};
\foreach \a in {1,2,3} \foreach \b in {1,2,3,4} \draw[cgray!50, thin] (mi\a) -- (mh\b);
\foreach \a in {1,2,3,4} \foreach \b in {1,2} \draw[cgray!50, thin] (mh\a) -- (mo\b);
\node[lab] at (8.9, -7.75) {(c) 多层感知机：层间全连接，“深度”来自层层复合};
\end{tikzpicture}
 引用（caption/label 在主文件）
% 注意：本文件为片段，依赖主文件导言区的 tikz/pgfplots 宏包、颜色定义与 tikzset 样式，不可单独编译。

\begin{tikzpicture}
% ---------- (a) 单神经元 ----------
\node[gry, minimum width=9mm] (x1) at (1.7, 0.85) {$x_1$};
\node[gry, minimum width=9mm] (x2) at (1.7, 0) {$x_2$};
\node[gry, minimum width=9mm] (x3) at (1.7, -0.85) {$x_3$};
\node[blk, align=center] (sum) at (3.9, 0) {加权和\\ $z=\mathbf{w}^{\top}\mathbf{x}+b$};
\node[grn, minimum width=9mm] (act) at (6.0, 0) {$\sigma$};
\node[gry, minimum width=9mm] (out) at (7.7, 0) {$a$};
\draw[arr] (x1.east) -- (sum.west);
\draw[arr] (x2.east) -- (sum.west);
\draw[arr] (x3.east) -- (sum.west);
\draw[arr] (sum) -- (act);
\draw[arr] (act) -- (out);
\node[lab] at (4.65, -1.55) {(a) 人工神经元：加权和 + 非线性};
% ---------- (b) 计算图与反向传播 ----------
\node[gry, minimum width=9mm] (bx) at (9.2, 0) {$\mathbf{x}$};
\node[blk, minimum width=13mm] (bl) at (10.8, 0) {线性层};
\node[grn, minimum width=11mm] (ba) at (12.5, 0) {激活};
\node[blk, minimum width=13mm] (bo) at (14.2, 0) {输出层};
\node[orn, minimum width=13mm] (bl2) at (15.9, 0) {损失 $L$};
\draw[arr] (bx)--(bl); \draw[arr] (bl)--(ba); \draw[arr] (ba)--(bo); \draw[arr] (bo)--(bl2);
\draw[rrr] (bl2.north) to[bend right=28] node[lab, above, pos=0.5] {反向：$\partial L/\partial\theta$ 沿链式法则逐层回传} (bl.north);
\node[lab] at (12.4, -1.55) {(b) 计算图：前向算值，反向归因};
% ---------- (c) 多层感知机（与上行整体居中） ----------
\node[lab] at (6.6, -2.7) {输入层};
\node[lab] at (8.9, -2.7) {隐藏层};
\node[lab] at (11.2, -2.7) {输出层};
\foreach \i in {1,2,3} \node[circle, draw=cgray, fill=cgray!8, minimum size=7mm] (mi\i) at (6.6, -3.35-0.95*\i) {};
\foreach \i in {1,2,3,4} \node[circle, draw=cblue, fill=cblue!8, minimum size=7mm] (mh\i) at (8.9, -2.9-0.95*\i) {};
\foreach \i in {1,2} \node[circle, draw=cgreen!70!black, fill=cgreen!8, minimum size=7mm] (mo\i) at (11.2, -3.35-1.4*\i) {};
\foreach \a in {1,2,3} \foreach \b in {1,2,3,4} \draw[cgray!50, thin] (mi\a) -- (mh\b);
\foreach \a in {1,2,3,4} \foreach \b in {1,2} \draw[cgray!50, thin] (mh\a) -- (mo\b);
\node[lab] at (8.9, -7.75) {(c) 多层感知机：层间全连接，“深度”来自层层复合};
\end{tikzpicture}

\caption{深度学习的三个基本视图：(a) 人工神经元；(b) 计算图上的前向传播与反向传播；(c) 多层感知机——现代 LLM 就是把 (c) 的宽度与深度放大七个数量级、再把隐藏层换成注意力与 SwiGLU 的产物。}
\label{fig:neuron}
\end{figure}

\begin{infobox}{核心洞察：深度学习的计算本质}
深度学习的所有操作——卷积、注意力、前馈——在底层都可以归结为\textbf{矩阵乘法}和\textbf{逐元素变换}。GPU 的 Tensor Core 专为大矩阵乘法优化，这使得深度学习训练/推理能在数百到数千张 GPU 上并行执行，吞吐量达数十 PFLOPS。
\end{infobox}
% ---------- 1.2 ----------
\subsection{激活函数}
激活函数引入\textbf{非线性}，使网络能拟合任意复杂函数。非线性之所以不可或缺，是因为\emph{线性变换的复合仍是线性变换}——若没有激活函数，再深的网络也只是单个线性映射 $y=(W_L\cdots W_1)x+\mathbf{b}$，表达能力与单层无异。因此激活函数是“深度”产生意义的根本前提。

工程上，激活函数的选择需权衡三点：\textbf{计算代价}（$\tanh$ 类需要指数运算，ReLU/GELU 近似更廉价）、\textbf{梯度特性}（负区间梯度是否消失）、\textbf{平滑性}（ReLU 在 0 点不可导，训练大模型时 GELU/SiLU 的平滑性更友好）。下表汇总了主流激活函数：

\begin{table}[htbp]
\centering
\small
\begin{tabular}{@{}llp{5.5cm}@{}}
\toprule
\textbf{函数} & \textbf{公式} & \textbf{特点}\\
\midrule
Sigmoid & $1/(1+e^{-x})$ & 输出 $(0,1)$，梯度消失严重\\
Tanh & $(e^x-e^{-x})/(e^x+e^{-x})$ & 输出 $(-1,1)$，零中心\\
ReLU & $\max(0,x)$ & 计算快，但存在“死亡神经元”\\
\textbf{GELU} & $x\cdot\Phi(x)$ & LLaMA、GPT 系列主流，平滑\\
\textbf{SwiGLU} & $\text{Swish}(xW_1)\odot(xW_2)$ & LLaMA/LLM 前馈层首选\\
SiLU / Swish & $x\cdot\sigma(x)$ & 门控激活\\
GEGLU & $\text{GeLU}(xW_1)\odot(xW_2)$ & GLM、Falcon 使用\\
\bottomrule
\end{tabular}
\caption{主流激活函数对比}
\end{table}

\begin{keybox}{趋势：门控激活成为 LLM 标配}
现代 LLM 普遍采用 \textbf{SwiGLU} 或 \textbf{GEGLU} 作为 FFN 激活。相比 ReLU 有 1--3\% 的困惑度（perplexity）优势，代价是参数量增加 50\%（FFN 中有两个门控矩阵 $W_1$ 和 $W_2$）。这一“用参数换精度”的权衡在大规模训练中非常划算。
\end{keybox}
% ---------- 1.3 ----------
\subsection{反向传播与优化器}
\textbf{反向传播}（Backpropagation）利用链式法则，从损失函数出发逐层计算梯度 $\partial L/\partial W$：
\[
L = \mathcal{L}(\hat{y},\, y), \qquad
\frac{\partial L}{\partial W_k} = \frac{\partial L}{\partial \mathbf{a}_{k+1}} \cdot \frac{\partial \mathbf{a}_{k+1}}{\partial \mathbf{z}_{k+1}} \cdot \frac{\partial \mathbf{z}_{k+1}}{\partial W_k}
\]
核心优化器对比：

\begin{table}[htbp]
\centering
\small
\begin{tabular}{@{}lll@{}}
\toprule
\textbf{优化器} & \textbf{核心思想} & \textbf{使用场景}\\
\midrule
SGD + Momentum & 经典梯度 + 惯性 & 计算机视觉\\
\textbf{Adam / AdamW} & 自适应学习率 + 权重解耦 & \textbf{LLM 主流}\\
Lion & 仅用符号更新，省显存 & 研究前沿\\
Sophia & 二阶近似，收敛更快 & 研究前沿\\
Muon & 2025 新兴，Transformer 表现优异 & 前沿探索\\
\bottomrule
\end{tabular}
\caption{主流优化器对比}
\end{table}

\textbf{反向传播的直觉}可以概括为“误差的归因过程”：前向传播把输入层层加工成预测；反向传播则追问“这次预测错了，每个参数该负多少责任”——沿计算图逆行，用链式法则把责任（梯度）按贡献比例分摊到每个参数。现代框架（PyTorch autograd）已把这一过程完全自动化，工程师只需定义前向计算；但理解它依然必要——第 4 章的激活检查点（以重算换显存）、1.4 节的梯度裁剪，以及“深层网络为何需要残差连接”（第 2 章），全部建立在这张计算图的拓扑之上。

LLM 预训练的典型 AdamW 配置：$\beta_1=0.9,\;\beta_2=0.95,\;\varepsilon=10^{-8}$，学习率余弦衰减，warmup 2000--10000 步。

\textbf{为什么是 AdamW 而不是 Adam}：Adam 将权重衰减正则化混入自适应学习率中，$\lambda$ 会被二阶矩估计 $v_t$ 除小，导致对大梯度参数正则化不足。AdamW（Loshchilov \& Hutter, 2019）将权重衰减\emph{解耦}为独立项 $w \leftarrow w - \eta(\text{Adam 更新} + \lambda w)$，恢复了 L2 正则的数学语义。这一细节对 Transformer 的泛化与训练稳定性有可测量的影响，因此成为 LLM 预训练的默认选择。

\textbf{学习率调度}同样关键：warmup 阶段让 Adam 的二阶矩估计先“热身”，避免初始步长过大导致 loss spike；其后采用余弦衰减至 $10\%$ 左右，在后期保持精细收敛。经验上，\emph{warmup 比例、峰值学习率与 batch size}三者强耦合——batch 每翻一倍，峰值学习率约可随之翻倍（线性缩放法则），这也是大规模训练调参的第一杠杆。
% ---------- 1.4 ----------
\subsection{正则化与训练稳定性}
深度学习训练是“高维非凸优化 + 有限精度 + 大规模并行”的叠加问题，稳定性工程往往与算法本身同等重要。以下逐项说明各技术在 LLM 训练中的角色：
\begin{itemize}
  \item \textbf{Dropout}：随机丢弃神经元，防过拟合（LLM 中通常 $\text{rate}=0.1$）
  \item \textbf{Weight Decay}：$L_2$ 正则化（AdamW 中解耦实现）
  \item \textbf{归一化}：
  \begin{itemize}
    \item \textbf{RMSNorm} 已成为 LLM 标配（比 LayerNorm 少一次均值计算，数值更稳定）
    \item 原始 Transformer 使用 Post-Norm（先子层后归一化），梯度需穿过归一化层才能回传，深层易不稳；现代 LLM 几乎全部改用 Pre-Norm（先归一化再进子层）
  \end{itemize}
  \item \textbf{Gradient Clipping}：裁剪梯度范数（$\|\mathbf{g}\| \leq 1.0$），防止梯度爆炸
  \item \textbf{Mixed Precision}：BF16 前向/反向 + FP32 主权重，2$\times$ 加速且显存减半；2024 年底 DeepSeek-V3 进一步把 FP8 混合精度训练推向 671B 规模（对嵌入、归一化等敏感模块保留高精度），标志低精度训练进入千亿级主流
  \item \textbf{Gradient Accumulation}：小 batch 累积梯度模拟大 batch
\end{itemize}

\begin{keybox}{工程要点：loss spike 是大规模训练的头号敌人}
数百张 GPU 上训练千亿参数模型时，偶发的\textbf{loss spike}（损失突然飙升）是最高频事故，常见诱因包括：学习率过高、数据中的异常批次、数值溢出（FP16 下更易发生，BF16 范围更大故更稳）、通信竞争条件。工业界的标准防御是：$\|g\|\le 1.0$ 梯度裁剪 + 损失异常检测（spike 检测即回滚 checkpoint）+ 保守 warmup。这解释了为什么 LLM 训练栈看起来“保守”——\emph{能稳定跑完}比单点性能更重要。
\end{keybox}
% ---------- 1.5 ----------
\subsection{损失函数}
LLM 预训练的核心损失是\textbf{交叉熵}（Cross-Entropy），等价于\textbf{负对数似然}（NLL）：
\[
\mathcal{L} \;=\; -\sum_{t=1}^{T} \log P\!\left(x_t \;\middle|\; x_{<t},\, \theta\right)
\]
\textbf{训练目标}：在给定前文条件下，\textbf{最大化下一个 token 的预测概率}。

整个预训练过程可以抽象为：
\begin{center}
\texttt{数据 $\to$ 前向传播 $\to$ 计算损失 $\to$ 反向传播 $\to$ 更新参数 $\to$ 循环}
\end{center}

与交叉熵直接关联的评估指标是\textbf{困惑度}（Perplexity, PPL）：
\[
\text{PPL} \;=\; \exp\!\left(-\frac{1}{T}\sum_{t=1}^{T}\log P(x_t \mid x_{<t})\right)
\]
直观含义是“模型在每个 token 上平均在多少个候选词之间犹豫”。PPL=100 意味着平均不确定性相当于在 100 个等概率选项中选一个。训练监控中，\emph{训练损失与验证损失同时下降}是健康的信号；验证损失先降后升则提示过拟合或数据配比失衡——这是判断“数据墙”逼近的最朴素手段。

最后补充一个容易被忽视的前置环节——\textbf{分词}（Tokenization）：现代 LLM 普遍使用 BPE（Byte Pair Encoding）类算法把词表压缩到 32K--150K 规模。词表大小直接决定嵌入层与输出投影的参数规模，也决定了中文等多字节语言每字占用的 token 数（token 越省，同样的上下文预算可容纳的中文内容越多）。这解释了为何各家模型在词表设计上差异显著：GPT-4 系列约 100K、LLaMA-2 约 32K、Qwen 系列约 150K。
% ===================== 第二章 =====================
\section{经典架构演进：CNN $\to$ RNN $\to$ Transformer}
\subsection{CNN：局部特征提取}
卷积神经网络（Convolutional Neural Network）的核心思想是\textbf{局部感受野 + 权值共享 + 平移不变性}。

\textbf{关键组件}：卷积层 $\to$ 归一化 $\to$ 激活 $\to$ 池化。卷积操作可以看作一组滑动窗口滤波器：
\[
\text{out}[\mathbf{i}] \;=\; \sum_{\mathbf{j}} \text{in}[\mathbf{i}+\mathbf{j}] \cdot \text{kernel}[\mathbf{j}] + b
\]

\textbf{里程碑}：LeNet (1998) $\to$ AlexNet (2012) $\to$ VGG (2014) $\to$ \textbf{ResNet} (2015) $\to$ EfficientNet (2019)。

其中 \textbf{ResNet 的残差连接} $y = F(\mathbf{x}) + \mathbf{x}$ 解决了深层网络退化问题，这一思想后来被 Transformer 直接继承（每个 block 都有 skip connection，这也是 LLM 能稳定堆叠 80--100 层的关键）。

\emph{需要澄清的是：CNN 并未退出舞台}，而是转换了角色——它仍是计算机视觉、语音识别、推荐系统的主力架构，并且作为视觉编码器（ViT/ConvNeXt 谱系）成为 2025--2026 年多模态大模型（LLaVA、Qwen-VL、InternVL）的第一级组件，详见第 8 章。在纯文本 LLM 领域，CNN 确实让位于注意力架构。

\textbf{CNN 留给 LLM 时代的两份遗产}值得一提。其一是\emph{归纳偏置}（inductive bias）的设计哲学：局部感受野、权值共享、平移不变性，本质上是在架构里“预先写入”关于数据的先验知识——数据越少，先验越值钱。ViT（2020）证明在超大规模数据下纯注意力可以抛弃这些先验，但在端侧小数据、实时性要求高的场景（手机拍照、语音前端），CNN 的先验依然划算。其二是\emph{层级化特征}的思想（LeNet $\to$ AlexNet $\to$ ResNet 一路验证的“边缘 $\to$ 纹理 $\to$ 部件 $\to$ 物体”逐级抽象），Swin Transformer 等层级化 ViT 正是把这一思想移植回注意力架构——多模态模型的视觉塔至今仍在“纯 ViT”与“层级化/卷积混合”两条路线间权衡。
\subsection{RNN / LSTM / GRU：序列建模}
\textbf{核心思想}：循环连接，隐状态 $h_t$ 携带历史信息：
\[
h_t = \tanh(W_{xh}\, x_t + W_{hh}\, h_{t-1} + b)
\]

\textbf{LSTM} 引入三个门控机制（遗忘门 $f_t$、输入门 $i_t$、输出门 $o_t$），用细胞状态 $C_t$ 实现“选择性记忆”：
\[
f_t = \sigma(W_f [h_{t-1}, x_t] + b_f), \quad
C_t = f_t \odot C_{t-1} + i_t \odot \tilde{C}_t
\]

\textbf{GRU} 简化 LSTM，用更新门 $z_t$ 和重置门 $r_t$ 替代：
\[
z_t = \sigma(W_z [h_{t-1}, x_t]), \quad
h_t = (1-z_t) \odot h_{t-1} + z_t \odot \tilde{h}_t
\]

\begin{warnbox}{RNN 的三大致命缺陷}
\begin{enumerate}
  \item \textbf{串行计算}：$h_t$ 依赖 $h_{t-1}$，无法并行，长序列极慢
  \item \textbf{梯度消失/爆炸}：BPTT 链式传递中梯度指数衰减或增长
  \item \textbf{长距离依赖衰减}：信息经过 100+ 步后几乎完全丢失
\end{enumerate}
这三个缺陷使得 RNN 无法支撑百亿参数级的大规模训练，最终被 Transformer 取代。
\end{warnbox}

\textbf{梯度消失的数学根源}：BPTT 中，$h_t$ 对 $k$ 步之前隐状态的梯度形如
\[
\frac{\partial h_t}{\partial h_{t-k}} = \prod_{i=0}^{k-1} \frac{\partial h_{t-i}}{\partial h_{t-i-1}}
\]
即 $k$ 个 Jacobian 矩阵的连乘。当矩阵谱半径 $<1$ 时该乘积按 $\rho^k$ 指数衰减，$>1$ 时指数爆炸。LSTM 的门控（尤其是遗忘门对细胞状态的\emph{乘法直通}）本质上是给这条梯度链装上“阀门”，缓解了但不根除该问题。Transformer 的注意力让信息以 $O(1)$ 跳数直接传递，从机制上绕开了这条指数衰减链——这是架构层面的根治。

\textbf{RNN 的遗产与回潮}同样值得关注。一方面，RNN 的“常数状态、线性成本”在流式与端侧场景仍有生命力：语音识别的前端、IoT 设备上的唤醒词模型至今大量使用 GRU/LSTM。另一方面，2023 年以来 RWKV、Mamba 等线性 RNN/SSM 以“训练可并行 + 推理 $O(1)$ 递推”的姿态卷土重来，证明循环思想并未过时——它们与注意力的融合（混合架构）是第 5 章的重要议题。可以说，RNN 与注意力之争不是“谁消灭谁”，而是“精确检索 vs 压缩记忆”这一对矛盾在不同成本约束下的再平衡。

\subsection{Transformer：范式革命}
注意力的思想并非一夜出现：2014 年 Bahdanau 等人在机器翻译中引入注意力以缓解 RNN 的信息瓶颈（Luong et al., 2015 简化为其点积形式），2017 年 Google 论文 \textbf{“Attention is All You Need”}（Vaswani et al.）将注意力推广为\emph{全部}序列交互机制，彻底改变了深度学习格局：

\begin{table}[htbp]
\centering
\small
\begin{tabular}{@{}lll@{}}
\toprule
\textbf{维度} & \textbf{RNN/LSTM} & \textbf{Transformer}\\
\midrule
并行度 & 串行 $O(T)$ & 全并行 $O(1)$\\
长距离依赖 & 指数衰减 & $O(1)$ 跳数\\
可扩展性 & 受限于梯度 & 自然扩展\\
推理效率 & 无状态 & 需 KV Cache\\
硬件利用 & GPU 利用率低 & 高（矩阵乘为主）\\
\bottomrule
\end{tabular}
\caption{RNN 与 Transformer 架构对比}
\end{table}

\begin{keybox}{关键洞察}
注意力机制让任意两个位置之间的信息传递只需 \textbf{1 步}，而 RNN 需要 $O(T)$ 步。这使得并行训练成为可能，是深度学习进入“规模时代”的根本原因。没有 Transformer 的并行性，就没有 GPT、LLaMA 等百亿/千亿参数模型。
\end{keybox}

\begin{figure}[htbp]
\centering
\resizebox{\textwidth}{!}{%
% fig03_architecture_timeline.tex — 图3 架构演进时间线
% 《深度学习与大语言模型：前沿技术全景报告》图示库
% 由 dl_llm_report.tex 抽取，经 % fig03_architecture_timeline.tex — 图3 架构演进时间线
% 《深度学习与大语言模型：前沿技术全景报告》图示库
% 由 dl_llm_report.tex 抽取，经 % fig03_architecture_timeline.tex — 图3 架构演进时间线
% 《深度学习与大语言模型：前沿技术全景报告》图示库
% 由 dl_llm_report.tex 抽取，经 \input{figs/fig03_architecture_timeline} 引用（caption/label 在主文件）
% 注意：本文件为片段，依赖主文件导言区的 tikz/pgfplots 宏包、颜色定义与 tikzset 样式，不可单独编译。

\begin{tikzpicture}[x=1.25cm]
\draw[arr] (-0.4,0) -- (14.5,0) node[lab, below left=0mm and -2mm] {年份};
\foreach \i/\yr/\side/\sty/\txt in {
  0/1998/1/gry/{LeNet\\CNN 时代开启},
  1/2012/-1/gry/{AlexNet\\深度学习爆发},
  2/2014/1/gry/{Seq2Seq+注意力\\RNN 时代雏形},
  3/2015/-1/gry/{ResNet\\残差连接},
  4/2017/1/blk/{\textbf{Transformer}\\范式革命},
  5/2018/-1/blk/{BERT / GPT-1\\预训练+微调},
  6/2020/1/blk/{GPT-3\\上下文学习涌现},
  7/2022/-1/orn/{ChatGPT\\对齐走向产品},
  8/2023/1/blk/{LLaMA / GPT-4\\开源与多模态},
  9/2024/-1/orn/{o1 / DeepSeek-V2\\推理模型元年},
  10/2025/1/orn/{R1 / Qwen3 / GPT-5\\开源推理普惠},
  11/2026/-1/orn/{Agentic·世界模型\\从对话到行动}} {
  \node[\sty, font=\scriptsize, align=center] at (\i*1.25, \side*1.5) {\txt};
  \draw[cgray, thin] (\i*1.25,0) -- (\i*1.25, \side*0.55);
  \node[lab, font=\scriptsize] at (\i*1.25, -\side*0.4) {\yr};
}
\end{tikzpicture}
 引用（caption/label 在主文件）
% 注意：本文件为片段，依赖主文件导言区的 tikz/pgfplots 宏包、颜色定义与 tikzset 样式，不可单独编译。

\begin{tikzpicture}[x=1.25cm]
\draw[arr] (-0.4,0) -- (14.5,0) node[lab, below left=0mm and -2mm] {年份};
\foreach \i/\yr/\side/\sty/\txt in {
  0/1998/1/gry/{LeNet\\CNN 时代开启},
  1/2012/-1/gry/{AlexNet\\深度学习爆发},
  2/2014/1/gry/{Seq2Seq+注意力\\RNN 时代雏形},
  3/2015/-1/gry/{ResNet\\残差连接},
  4/2017/1/blk/{\textbf{Transformer}\\范式革命},
  5/2018/-1/blk/{BERT / GPT-1\\预训练+微调},
  6/2020/1/blk/{GPT-3\\上下文学习涌现},
  7/2022/-1/orn/{ChatGPT\\对齐走向产品},
  8/2023/1/blk/{LLaMA / GPT-4\\开源与多模态},
  9/2024/-1/orn/{o1 / DeepSeek-V2\\推理模型元年},
  10/2025/1/orn/{R1 / Qwen3 / GPT-5\\开源推理普惠},
  11/2026/-1/orn/{Agentic·世界模型\\从对话到行动}} {
  \node[\sty, font=\scriptsize, align=center] at (\i*1.25, \side*1.5) {\txt};
  \draw[cgray, thin] (\i*1.25,0) -- (\i*1.25, \side*0.55);
  \node[lab, font=\scriptsize] at (\i*1.25, -\side*0.4) {\yr};
}
\end{tikzpicture}
 引用（caption/label 在主文件）
% 注意：本文件为片段，依赖主文件导言区的 tikz/pgfplots 宏包、颜色定义与 tikzset 样式，不可单独编译。

\begin{tikzpicture}[x=1.25cm]
\draw[arr] (-0.4,0) -- (14.5,0) node[lab, below left=0mm and -2mm] {年份};
\foreach \i/\yr/\side/\sty/\txt in {
  0/1998/1/gry/{LeNet\\CNN 时代开启},
  1/2012/-1/gry/{AlexNet\\深度学习爆发},
  2/2014/1/gry/{Seq2Seq+注意力\\RNN 时代雏形},
  3/2015/-1/gry/{ResNet\\残差连接},
  4/2017/1/blk/{\textbf{Transformer}\\范式革命},
  5/2018/-1/blk/{BERT / GPT-1\\预训练+微调},
  6/2020/1/blk/{GPT-3\\上下文学习涌现},
  7/2022/-1/orn/{ChatGPT\\对齐走向产品},
  8/2023/1/blk/{LLaMA / GPT-4\\开源与多模态},
  9/2024/-1/orn/{o1 / DeepSeek-V2\\推理模型元年},
  10/2025/1/orn/{R1 / Qwen3 / GPT-5\\开源推理普惠},
  11/2026/-1/orn/{Agentic·世界模型\\从对话到行动}} {
  \node[\sty, font=\scriptsize, align=center] at (\i*1.25, \side*1.5) {\txt};
  \draw[cgray, thin] (\i*1.25,0) -- (\i*1.25, \side*0.55);
  \node[lab, font=\scriptsize] at (\i*1.25, -\side*0.4) {\yr};
}
\end{tikzpicture}

}%
\caption{深度学习架构演进时间线（1998--2026）：2017 年 Transformer 之前是“为特定任务设计归纳偏置”的时代，之后进入“通用架构 + 规模扩展”的时代；2024 年起竞争焦点转向推理能力与效率。}
\label{fig:timeline}
\end{figure}

图~\ref{fig:timeline} 汇总了这一演进脉络。值得注意的还有\textbf{“架构冻结”现象}：2017 年至今，“注意力 + 残差 + 前馈”这一骨架几乎未变，创新集中在四个“插槽”上——\emph{位置编码}（绝对 $\to$ RoPE）、\emph{注意力变体}（MHA $\to$ GQA/MLA）、\emph{归一化}（Post-LN $\to$ Pre-RMSNorm）、\emph{激活函数}（ReLU $\to$ SwiGLU）。理解这四个插槽，就理解了 LLM 架构史的八成（详见第 3、5 章）；2024 年后的 MoE 与 SSM 混合架构，也只是在这副骨架上做“稀疏化”与“线性化”的深加工。
% ===================== 第三章 =====================
\section{Transformer 深度剖析：注意力机制的数学本质}
\subsection{标准 Transformer 架构}
一个完整的 Transformer Block 包含\textbf{两个子层}：多头自注意力 + 前馈网络，每个子层都有\textbf{残差连接}和\textbf{归一化}：

\begin{figure}[htbp]
\centering
% fig04_transformer_block.tex — 图4 现代 decoder-only Transformer Block
% 《深度学习与大语言模型：前沿技术全景报告》图示库
% 由 dl_llm_report.tex 抽取，经 % fig04_transformer_block.tex — 图4 现代 decoder-only Transformer Block
% 《深度学习与大语言模型：前沿技术全景报告》图示库
% 由 dl_llm_report.tex 抽取，经 % fig04_transformer_block.tex — 图4 现代 decoder-only Transformer Block
% 《深度学习与大语言模型：前沿技术全景报告》图示库
% 由 dl_llm_report.tex 抽取，经 \input{figs/fig04_transformer_block} 引用（caption/label 在主文件）
% 注意：本文件为片段，依赖主文件导言区的 tikz/pgfplots 宏包、颜色定义与 tikzset 样式，不可单独编译。

\begin{tikzpicture}[node distance=4.5mm]
\node[gry, minimum width=62mm] (inp) {输入：Token Embedding + 位置信息（RoPE）};
\node[blk, below=of inp, minimum width=48mm] (n1) {RMSNorm（Pre-Norm）};
\node[grn, below=of n1, minimum width=48mm] (attn) {多头自注意力（因果掩码）};
\node[gry, below=of attn, minimum width=48mm, minimum height=5.5mm, inner sep=2pt] (add1) {$\oplus$ 残差合并};
\node[blk, below=of add1, minimum width=48mm] (n2) {RMSNorm（Pre-Norm）};
\node[grn, below=of n2, minimum width=48mm] (ffn) {前馈网络 SwiGLU FFN};
\node[gry, below=of ffn, minimum width=48mm, minimum height=5.5mm, inner sep=2pt] (add2) {$\oplus$ 残差合并};
\node[orn, below=of add2, minimum width=48mm] (out) {LM Head + Softmax $\to$ 下一 token 分布};
\draw[arr] (inp)--(n1); \draw[arr] (n1)--(attn); \draw[arr] (attn)--(add1);
\draw[arr] (add1)--(n2); \draw[arr] (n2)--(ffn); \draw[arr] (ffn)--(add2); \draw[arr] (add2)--(out);
\draw[brr] (inp.east) -- ++(1.4cm,0) |- (add1.east);
\draw[brr] (add1.east) -- ++(0.7cm,0) |- (add2.east);
\draw[decorate, decoration={brace, amplitude=7pt, mirror}, thick, cgray] ($(n1.north west)+(-3mm,0)$) -- ($(add2.south west)+(-3mm,0)$);
\node[lab, left=4mm of add1, align=center] {$\times\,N$ 层堆叠\\（LLaMA-2 70B：$N{=}80$）};
\node[lab, right=4mm of attn, align=left] {KV Cache 产生地\\（第 6 章优化的对象）};
\node[lab, right=4mm of ffn, align=left] {参数占全模型约 $2/3$\\（知识存储主体，MoE 改造对象）};
\end{tikzpicture}
 引用（caption/label 在主文件）
% 注意：本文件为片段，依赖主文件导言区的 tikz/pgfplots 宏包、颜色定义与 tikzset 样式，不可单独编译。

\begin{tikzpicture}[node distance=4.5mm]
\node[gry, minimum width=62mm] (inp) {输入：Token Embedding + 位置信息（RoPE）};
\node[blk, below=of inp, minimum width=48mm] (n1) {RMSNorm（Pre-Norm）};
\node[grn, below=of n1, minimum width=48mm] (attn) {多头自注意力（因果掩码）};
\node[gry, below=of attn, minimum width=48mm, minimum height=5.5mm, inner sep=2pt] (add1) {$\oplus$ 残差合并};
\node[blk, below=of add1, minimum width=48mm] (n2) {RMSNorm（Pre-Norm）};
\node[grn, below=of n2, minimum width=48mm] (ffn) {前馈网络 SwiGLU FFN};
\node[gry, below=of ffn, minimum width=48mm, minimum height=5.5mm, inner sep=2pt] (add2) {$\oplus$ 残差合并};
\node[orn, below=of add2, minimum width=48mm] (out) {LM Head + Softmax $\to$ 下一 token 分布};
\draw[arr] (inp)--(n1); \draw[arr] (n1)--(attn); \draw[arr] (attn)--(add1);
\draw[arr] (add1)--(n2); \draw[arr] (n2)--(ffn); \draw[arr] (ffn)--(add2); \draw[arr] (add2)--(out);
\draw[brr] (inp.east) -- ++(1.4cm,0) |- (add1.east);
\draw[brr] (add1.east) -- ++(0.7cm,0) |- (add2.east);
\draw[decorate, decoration={brace, amplitude=7pt, mirror}, thick, cgray] ($(n1.north west)+(-3mm,0)$) -- ($(add2.south west)+(-3mm,0)$);
\node[lab, left=4mm of add1, align=center] {$\times\,N$ 层堆叠\\（LLaMA-2 70B：$N{=}80$）};
\node[lab, right=4mm of attn, align=left] {KV Cache 产生地\\（第 6 章优化的对象）};
\node[lab, right=4mm of ffn, align=left] {参数占全模型约 $2/3$\\（知识存储主体，MoE 改造对象）};
\end{tikzpicture}
 引用（caption/label 在主文件）
% 注意：本文件为片段，依赖主文件导言区的 tikz/pgfplots 宏包、颜色定义与 tikzset 样式，不可单独编译。

\begin{tikzpicture}[node distance=4.5mm]
\node[gry, minimum width=62mm] (inp) {输入：Token Embedding + 位置信息（RoPE）};
\node[blk, below=of inp, minimum width=48mm] (n1) {RMSNorm（Pre-Norm）};
\node[grn, below=of n1, minimum width=48mm] (attn) {多头自注意力（因果掩码）};
\node[gry, below=of attn, minimum width=48mm, minimum height=5.5mm, inner sep=2pt] (add1) {$\oplus$ 残差合并};
\node[blk, below=of add1, minimum width=48mm] (n2) {RMSNorm（Pre-Norm）};
\node[grn, below=of n2, minimum width=48mm] (ffn) {前馈网络 SwiGLU FFN};
\node[gry, below=of ffn, minimum width=48mm, minimum height=5.5mm, inner sep=2pt] (add2) {$\oplus$ 残差合并};
\node[orn, below=of add2, minimum width=48mm] (out) {LM Head + Softmax $\to$ 下一 token 分布};
\draw[arr] (inp)--(n1); \draw[arr] (n1)--(attn); \draw[arr] (attn)--(add1);
\draw[arr] (add1)--(n2); \draw[arr] (n2)--(ffn); \draw[arr] (ffn)--(add2); \draw[arr] (add2)--(out);
\draw[brr] (inp.east) -- ++(1.4cm,0) |- (add1.east);
\draw[brr] (add1.east) -- ++(0.7cm,0) |- (add2.east);
\draw[decorate, decoration={brace, amplitude=7pt, mirror}, thick, cgray] ($(n1.north west)+(-3mm,0)$) -- ($(add2.south west)+(-3mm,0)$);
\node[lab, left=4mm of add1, align=center] {$\times\,N$ 层堆叠\\（LLaMA-2 70B：$N{=}80$）};
\node[lab, right=4mm of attn, align=left] {KV Cache 产生地\\（第 6 章优化的对象）};
\node[lab, right=4mm of ffn, align=left] {参数占全模型约 $2/3$\\（知识存储主体，MoE 改造对象）};
\end{tikzpicture}

\caption{现代 decoder-only Transformer Block（Pre-Norm 结构）：每个子层先归一化再计算，残差连接提供无衰减的梯度“高速公路”；整个 Block 以 $\times N$ 堆叠，末端 LM Head 将隐状态映射回词表。}
\label{fig:block}
\end{figure}

图~\ref{fig:block} 展示了这一结构。这里补充一个关键视角——\textbf{FFN 是模型的“键值记忆体”}：研究（Geva et al., 2021）表明，FFN 的第一层相当于一组“模式检测键”（检测输入中的某种语义模式），第二层是与之配对的“知识值”（把该模式映射为下一个词的倾向）。模型的成块事实知识（人物、 API、历史事件）大量存储在这一层层的键值对里——这解释了三件事：为什么 FFN 占参数大头、为什么 MoE（第 5 章）改造的是 FFN 而非注意力、以及为什么“知识编辑”研究（如 ROME/MEMIT，定位并改写 FFN 中的特定键值对）瞄准的也是 FFN。相应地，\textbf{注意力层负责“沟通”}：在 token 之间搬运与组合信息（指代、对齐、检索），两个子层一个管“存储”一个管“路由”，缺一不可。

\textbf{三个现代变体值得注意}：(1)\emph{Pre-Norm}——先归一化再做子层运算，原始 Transformer 使用 Post-Norm，而现代 LLM 几乎全部采用 Pre-Norm，因为梯度可沿残差通路近乎无衰减地回传，训练显著更稳；(2)\emph{RMSNorm} 替代 LayerNorm——省去均值中心化计算，数值上更鲁棒；(3)\emph{FFN 膨胀比}——前馈层通常把维度扩张至 $2.7\text{--}4\times$ 隐藏维度（LLaMA 取 $\frac{8}{3}\pi d$ 以配合 SwiGLU 的三路矩阵），FFN 往往占全模型 $\frac{2}{3}$ 以上的参数，是“知识存储”的主要载体，也是 MoE（第 5 章）改造的对象。

\subsection{自注意力机制：数学细节}
\subsubsection{单头注意力}
\[
\text{Attention}(Q, K, V) \;=\; \text{softmax}\!\left(\frac{Q\, K^\top}{\sqrt{d_k}}\right) V
\]
\begin{itemize}
  \item $Q, K, V$ 形状：$(\text{batch},\; \text{seq\_len},\; d_k)$
  \item $QK^\top$ 形状：$(\text{batch},\; n,\; n)$ —— 这就是 $O(n^2)$ 复杂度的来源
  \item $\sqrt{d_k}$ 缩放：防止点积值过大导致 softmax 饱和（当 $d_k$ 大时，$q\cdot k$ 的方差为 $d_k$，除以 $\sqrt{d_k}$ 使方差归一化）
  \item 时间复杂度：$O(n^2 \cdot d)$
  \item 空间复杂度：$O(n^2)$（注意力矩阵）
\end{itemize}

\begin{figure}[htbp]
\centering
% fig05_attention_pipeline.tex — 图5 自注意力计算流水线
% 《深度学习与大语言模型：前沿技术全景报告》图示库
% 由 dl_llm_report.tex 抽取，经 % fig05_attention_pipeline.tex — 图5 自注意力计算流水线
% 《深度学习与大语言模型：前沿技术全景报告》图示库
% 由 dl_llm_report.tex 抽取，经 % fig05_attention_pipeline.tex — 图5 自注意力计算流水线
% 《深度学习与大语言模型：前沿技术全景报告》图示库
% 由 dl_llm_report.tex 抽取，经 \input{figs/fig05_attention_pipeline} 引用（caption/label 在主文件）
% 注意：本文件为片段，依赖主文件导言区的 tikz/pgfplots 宏包、颜色定义与 tikzset 样式，不可单独编译。

\begin{tikzpicture}[node distance=6mm]
\node[gry] (x) {输入 $X$};
\node[blk, right=6mm of x, minimum width=27mm] (p) {三路线性投影\\ 得到 $Q,\;K,\;V$};
\node[blk, right=6mm of p, minimum width=27mm] (mm) {相似度矩阵\\ $QK^\top/\sqrt{d_k}$};
\node[grn, right=6mm of mm, minimum width=22mm] (sm) {Softmax\\ 注意力权重};
\node[grn, below=10mm of sm, minimum width=23mm] (av) {加权求和\\ $\text{Attn}\cdot V$};
\node[blk, left=6mm of av, minimum width=21mm] (h1) {单头输出};
\node[blk, left=6mm of h1, minimum width=26mm] (cat) {Concat（$h$ 头并行）};
\node[orn, left=6mm of cat, minimum width=22mm] (o) {$W^O$ 输出};
\draw[arr] (x)--(p); \draw[arr] (p)--(mm); \draw[arr] (mm)--(sm);
\draw[arr] (sm.south) -- (av.north);
\draw[arr] (av)--(h1); \draw[arr] (h1)--(cat); \draw[arr] (cat)--(o);
\node[lab, above=0.5mm of p] {$h$ 个头各自独立执行这整条流水线（互不干扰）};
\node[lab, below=1mm of mm] {$n\times n$ 矩阵：$O(n^2)$ 复杂度的来源};
\node[lab, below=1mm of o] {输出返回残差流};
\end{tikzpicture}
 引用（caption/label 在主文件）
% 注意：本文件为片段，依赖主文件导言区的 tikz/pgfplots 宏包、颜色定义与 tikzset 样式，不可单独编译。

\begin{tikzpicture}[node distance=6mm]
\node[gry] (x) {输入 $X$};
\node[blk, right=6mm of x, minimum width=27mm] (p) {三路线性投影\\ 得到 $Q,\;K,\;V$};
\node[blk, right=6mm of p, minimum width=27mm] (mm) {相似度矩阵\\ $QK^\top/\sqrt{d_k}$};
\node[grn, right=6mm of mm, minimum width=22mm] (sm) {Softmax\\ 注意力权重};
\node[grn, below=10mm of sm, minimum width=23mm] (av) {加权求和\\ $\text{Attn}\cdot V$};
\node[blk, left=6mm of av, minimum width=21mm] (h1) {单头输出};
\node[blk, left=6mm of h1, minimum width=26mm] (cat) {Concat（$h$ 头并行）};
\node[orn, left=6mm of cat, minimum width=22mm] (o) {$W^O$ 输出};
\draw[arr] (x)--(p); \draw[arr] (p)--(mm); \draw[arr] (mm)--(sm);
\draw[arr] (sm.south) -- (av.north);
\draw[arr] (av)--(h1); \draw[arr] (h1)--(cat); \draw[arr] (cat)--(o);
\node[lab, above=0.5mm of p] {$h$ 个头各自独立执行这整条流水线（互不干扰）};
\node[lab, below=1mm of mm] {$n\times n$ 矩阵：$O(n^2)$ 复杂度的来源};
\node[lab, below=1mm of o] {输出返回残差流};
\end{tikzpicture}
 引用（caption/label 在主文件）
% 注意：本文件为片段，依赖主文件导言区的 tikz/pgfplots 宏包、颜色定义与 tikzset 样式，不可单独编译。

\begin{tikzpicture}[node distance=6mm]
\node[gry] (x) {输入 $X$};
\node[blk, right=6mm of x, minimum width=27mm] (p) {三路线性投影\\ 得到 $Q,\;K,\;V$};
\node[blk, right=6mm of p, minimum width=27mm] (mm) {相似度矩阵\\ $QK^\top/\sqrt{d_k}$};
\node[grn, right=6mm of mm, minimum width=22mm] (sm) {Softmax\\ 注意力权重};
\node[grn, below=10mm of sm, minimum width=23mm] (av) {加权求和\\ $\text{Attn}\cdot V$};
\node[blk, left=6mm of av, minimum width=21mm] (h1) {单头输出};
\node[blk, left=6mm of h1, minimum width=26mm] (cat) {Concat（$h$ 头并行）};
\node[orn, left=6mm of cat, minimum width=22mm] (o) {$W^O$ 输出};
\draw[arr] (x)--(p); \draw[arr] (p)--(mm); \draw[arr] (mm)--(sm);
\draw[arr] (sm.south) -- (av.north);
\draw[arr] (av)--(h1); \draw[arr] (h1)--(cat); \draw[arr] (cat)--(o);
\node[lab, above=0.5mm of p] {$h$ 个头各自独立执行这整条流水线（互不干扰）};
\node[lab, below=1mm of mm] {$n\times n$ 矩阵：$O(n^2)$ 复杂度的来源};
\node[lab, below=1mm of o] {输出返回残差流};
\end{tikzpicture}

\caption{自注意力计算流水线：一个输入经三路投影产生 $Q,K,V$；相似度经缩放与 Softmax 归一化后对 $V$ 加权求和得到单头输出；$h$ 个头并行执行后拼接、经 $W^O$ 融合回残差流。}
\label{fig:attn}
\end{figure}

\subsubsection{多头注意力}
\[
\text{MHA}(X) = \text{Concat}(\text{head}_1, \ldots, \text{head}_h)\; W^O
\]
\[
\text{head}_i = \text{Attention}(X W_i^Q,\; X W_i^K,\; X W_i^V)
\]
\textbf{直觉理解}：把注意力看作一次\emph{软检索}——每个 token 的 $q$ 向量是“我要找什么”，所有 token 的 $k$ 向量是“我能提供什么标签”，$v$ 向量是“我的实际内容”。$QK^\top$ 计算所有标签匹配度，softmax 将其归一化为注意力分布，最后对内容 $V$ 做加权平均。整个机制不含任何显式规则，全部结构关系由训练中的梯度自动塑造。

每个头在不同子空间中学习不同的注意力模式：
\begin{itemize}
  \item 某些头关注\textbf{语法结构}（如主谓宾关系）
  \item 某些头关注\textbf{语义关系}（如同义替换、上下位）
  \item 某些头关注\textbf{指代消解}（“它”指代谁）
  \item 某些头关注\textbf{位置模式}（如相邻 token）
\end{itemize}

\begin{infobox}{注意力可视化的发现}
研究（Vig, 2019; Wang et al., 2022）表明：GPT-2 中约 10\% 的注意力头可被归类为“归纳头”（induction heads），它们学会了“当看到与上文相同 token 出现时，预测上文的下一个 token”。这是模型习得长距离模式匹配的关键机制。
\end{infobox}

\subsection{位置编码}
Transformer 本身不包含位置信息，需要外部注入。下表汇总了主要方案：

\begin{table}[htbp]
\centering
\small
\begin{tabular}{@{}lll@{}}
\toprule
\textbf{方法} & \textbf{原理} & \textbf{使用模型}\\
\midrule
绝对位置编码 & 正弦/余弦函数 & 原始 Transformer\\
可学习位置编码 & 直接学习位置向量 & BERT、GPT-2\\
\textbf{RoPE} & 将位置编码为旋转矩阵 & \textbf{LLaMA、Qwen、Mistral}\\
ALiBi & 注意力分数加距离偏置 & MPT、BLOOM\\
NTK-aware scaling & 外推 RoPE 到长序列 & LLaMA-2-128K\\
YaRN / NTK-by-parts & 分段外推 & 2025 年主流长上下文方案\\
\bottomrule
\end{tabular}
\caption{位置编码方案对比}
\end{table}

\textbf{RoPE（旋转位置编码）原理}——这是 LLM 领域最重要的创新之一：

对 query 和 key 向量，按位置 $m$ 施加\textbf{旋转矩阵}：
\[
f(q, m) = R_m \cdot q, \qquad
R_m = \begin{pmatrix}
\cos m\theta_1 & -\sin m\theta_1 \\
\sin m\theta_1 & \cos m\theta_1
\end{pmatrix} \oplus \cdots
\]
其中 $\theta_i = 10000^{-2i/d}$ 是频率参数，$\oplus$ 表示块对角拼接。

这样注意力分数自然包含了\textbf{相对位置}信息：
\[
q_m^\top k_n = q^\top R_{n-m}^{\top} R_n \, k = q^\top R_{n-m} \, k
\]

\begin{keybox}{RoPE 的优势}
\begin{enumerate}
  \item \textbf{天然支持相对位置}：注意力分数只依赖 $n-m$
  \item \textbf{推理时 KV Cache 友好}：不需要修改已缓存的 K
  \item \textbf{可外推}：通过调整 $\theta_i$ 或插值即可扩展上下文长度
  \item \textbf{无额外参数}：纯数学变换，不增加模型大小
\end{enumerate}
\end{keybox}

\textbf{长上下文外推}是 RoPE 的工程延伸：模型在 4K 上下文上训练后要在 128K 上使用，直接外推会导致注意力分数在训练时从未见过的相位组合上“失准”。NTK-aware scaling 通过降低基频 $\theta_i$（等效拉伸频谱）让远距位置落在训练分布的邻域内；YaRN（NTK-by-parts）进一步对高/低频分段处理，兼顾局部精度与长程外推，已成为 2025 年 128K--1M 长上下文方案的常用组件——配合训练时的长文本课程（curriculum）使用，效果最佳。

\subsection{解码器 vs 编码器}
\begin{table}[htbp]
\centering
\small
\begin{tabular}{@{}llll@{}}
\toprule
\textbf{特性} & \textbf{编码器} & \textbf{解码器} & \textbf{编码-解码}\\
& (BERT) & (GPT/LLaMA) & (T5)\\
\midrule
注意力 & 双向 & 因果 (Causal) & 编码双向 + 解码因果\\
预训练目标 & MLM & Next Token & Span Corruption\\
适用场景 & 分类、检索 & 生成、对话 & 翻译、摘要\\
推理速度 & 一次前向 & 自回归逐 token & 两阶段\\
LLM 主流 & -- & \checkmark & --\\
\bottomrule
\end{tabular}
\caption{三种 Transformer 变体对比}
\end{table}

\textbf{因果掩码}（Causal Mask）确保 token $t$ 只能看到 $\leq t$ 的信息：\[
M_{ij} = \begin{cases} 0 & \text{if } i \leq j \\ -\infty & \text{if } i > j \end{cases}
\]
在 softmax 中，$e^{-\infty}=0$，实现了严格的单向信息流。

因果掩码看似限制了信息流，实则带来两个关键工程红利：其一，\emph{训练时}同一 prompt 的所有位置可以并行计算（每个位置只依赖前文，互不冲突），一次前向即产生 $T$ 个训练样本；其二，\emph{推理时}位置 $t$ 的 Key/Value 只依赖前文，因此可缓存复用——这正是 KV Cache 的合法性来源（详见第 6 章）。可以说，自回归 LLM 的“可并行训练 + 可增量推理”这一看似矛盾的组合，统一由因果掩码保证。

% ===================== 第四章 =====================
\section{大语言模型（LLM）训练流水线}

LLM 的训练通常分为\textbf{三个核心阶段}：预训练 $\to$ 监督微调 $\to$ 人类对齐。下面逐一展开。

\begin{figure}[htbp]
\centering
% fig06_training_pipeline.tex — 图6 LLM 训练流水线三阶段
% 《深度学习与大语言模型：前沿技术全景报告》图示库
% 由 dl_llm_report.tex 抽取，经 % fig06_training_pipeline.tex — 图6 LLM 训练流水线三阶段
% 《深度学习与大语言模型：前沿技术全景报告》图示库
% 由 dl_llm_report.tex 抽取，经 % fig06_training_pipeline.tex — 图6 LLM 训练流水线三阶段
% 《深度学习与大语言模型：前沿技术全景报告》图示库
% 由 dl_llm_report.tex 抽取，经 \input{figs/fig06_training_pipeline} 引用（caption/label 在主文件）
% 注意：本文件为片段，依赖主文件导言区的 tikz/pgfplots 宏包、颜色定义与 tikzset 样式，不可单独编译。

\begin{tikzpicture}[node distance=9mm]
\node[blk, minimum width=46mm, minimum height=23mm, align=center] (pre) {\textbf{阶段一：预训练}\\ 万亿级 token / 万卡集群\\ 目标：下一 token 预测\\ {\scriptsize 成本：百万 GPU$\cdot$小时级}};
\node[blk, right=10mm of pre, minimum width=42mm, minimum height=23mm, align=center] (sft) {\textbf{阶段二：监督微调}\\ $10^4$--$10^6$ 条指令-回答对\\ 目标：遵循指令\\ {\scriptsize 成本：数千 GPU$\cdot$小时级}};
\node[blk, right=10mm of sft, minimum width=42mm, minimum height=23mm, align=center] (rl) {\textbf{阶段三：对齐}\\ RLHF / DPO / RLVR\\ 目标：有用且无害\\ {\scriptsize 成本：数万 GPU$\cdot$小时级起}};
\node[gry, below=9mm of sft, minimum width=104mm, minimum height=8mm] (cap) {能力演进：语言与知识 $\to$ 指令遵循 $\to$ 偏好对齐与推理};
\draw[arr] (pre)--(sft); \draw[arr] (sft)--(rl);
\draw[arr] (rl.south) -- ++(0,-20mm) -| node[lab, pos=0.25, below] {数据飞轮：用户交互 $\to$ 高质量数据回流} (pre.south);
\end{tikzpicture}
 引用（caption/label 在主文件）
% 注意：本文件为片段，依赖主文件导言区的 tikz/pgfplots 宏包、颜色定义与 tikzset 样式，不可单独编译。

\begin{tikzpicture}[node distance=9mm]
\node[blk, minimum width=46mm, minimum height=23mm, align=center] (pre) {\textbf{阶段一：预训练}\\ 万亿级 token / 万卡集群\\ 目标：下一 token 预测\\ {\scriptsize 成本：百万 GPU$\cdot$小时级}};
\node[blk, right=10mm of pre, minimum width=42mm, minimum height=23mm, align=center] (sft) {\textbf{阶段二：监督微调}\\ $10^4$--$10^6$ 条指令-回答对\\ 目标：遵循指令\\ {\scriptsize 成本：数千 GPU$\cdot$小时级}};
\node[blk, right=10mm of sft, minimum width=42mm, minimum height=23mm, align=center] (rl) {\textbf{阶段三：对齐}\\ RLHF / DPO / RLVR\\ 目标：有用且无害\\ {\scriptsize 成本：数万 GPU$\cdot$小时级起}};
\node[gry, below=9mm of sft, minimum width=104mm, minimum height=8mm] (cap) {能力演进：语言与知识 $\to$ 指令遵循 $\to$ 偏好对齐与推理};
\draw[arr] (pre)--(sft); \draw[arr] (sft)--(rl);
\draw[arr] (rl.south) -- ++(0,-20mm) -| node[lab, pos=0.25, below] {数据飞轮：用户交互 $\to$ 高质量数据回流} (pre.south);
\end{tikzpicture}
 引用（caption/label 在主文件）
% 注意：本文件为片段，依赖主文件导言区的 tikz/pgfplots 宏包、颜色定义与 tikzset 样式，不可单独编译。

\begin{tikzpicture}[node distance=9mm]
\node[blk, minimum width=46mm, minimum height=23mm, align=center] (pre) {\textbf{阶段一：预训练}\\ 万亿级 token / 万卡集群\\ 目标：下一 token 预测\\ {\scriptsize 成本：百万 GPU$\cdot$小时级}};
\node[blk, right=10mm of pre, minimum width=42mm, minimum height=23mm, align=center] (sft) {\textbf{阶段二：监督微调}\\ $10^4$--$10^6$ 条指令-回答对\\ 目标：遵循指令\\ {\scriptsize 成本：数千 GPU$\cdot$小时级}};
\node[blk, right=10mm of sft, minimum width=42mm, minimum height=23mm, align=center] (rl) {\textbf{阶段三：对齐}\\ RLHF / DPO / RLVR\\ 目标：有用且无害\\ {\scriptsize 成本：数万 GPU$\cdot$小时级起}};
\node[gry, below=9mm of sft, minimum width=104mm, minimum height=8mm] (cap) {能力演进：语言与知识 $\to$ 指令遵循 $\to$ 偏好对齐与推理};
\draw[arr] (pre)--(sft); \draw[arr] (sft)--(rl);
\draw[arr] (rl.south) -- ++(0,-20mm) -| node[lab, pos=0.25, below] {数据飞轮：用户交互 $\to$ 高质量数据回流} (pre.south);
\end{tikzpicture}

\caption{LLM 训练流水线三阶段：成本与数据规模逐级递减，而“行为塑造”的杠杆逐级增大；虚线反馈构成持续迭代的数据飞轮。}
\label{fig:pipeline}
\end{figure}

图~\ref{fig:pipeline} 概括了本章结构。三个阶段的成本结构截然不同：预训练吃掉 95\% 以上的总算力，却只决定模型的“智力上限”；SFT 与对齐花费不到 5\% 的算力，却几乎完全决定模型的“可用性”——这种不对称性是理解大模型产业分工的钥匙（基座厂商与后训练/应用团队的价值链切分）。

\subsection{阶段一：预训练（Pre-training）}

\textbf{目标}：学习语言的统计规律和世界知识。

\textbf{数据构成}：
\begin{itemize}
  \item 互联网爬取文本（CommonCrawl 等，数万亿 token）
  \item 书籍、学术论文、代码（GitHub）
  \item 高质量语料筛选：去重（MinHash/SimHash）、质量过滤、毒性过滤
  \item \textbf{数据配比}至关重要：通用网页 50--70\%，代码 10--20\%，书籍/学术 10--20\%
\end{itemize}

\textbf{数据工程是预训练的第一竞争力}。一个典型的工业级数据管线包含五个环节：\emph{(1) 原始获取}——CommonCrawl 等来源按快照（月级）批量下载，单次快照可达数百 TB；\emph{(2) 解析清洗}——HTML 解析为纯文本，剔除导航栏、脚本、乱码；\emph{(3) 质量过滤}——用小分类器或启发式规则（长度、符号比、重复 n-gram 占比）打分，各家经验是\emph{删掉一半低质数据后效果更好}；\emph{(4) 去重}——文档级去重（精确哈希）+ 段落/句级近似去重（MinHash + LSH，典型保留 85--95\%），重复数据会诱导模型“背诵”而非“理解”；\emph{(5) 安全与合规}——毒性过滤、PII 脱敏、版权排除列表。DeepSeek-V3 报告约 14.8T token 训练语料，其中合成与代码占比显著提升——\emph{数据的边际价值在 2024 年后已超过参数规模}，这已成为行业共识。

\textbf{两个常被低估的细节}：其一，\emph{数据配比是最大的超参数}——通用网页、代码、书籍、学术、多语言的比例直接决定模型的能力画像（代码数据提升的不只是编程，还有逻辑推理；多语言数据的清理质量决定“人格”是否统一），头部团队的做法是在训练中周期性跑评测、动态调整配比（data curriculum）；其二，\emph{tokenizer 决定语言经济学}——以 Qwen 的 15 万词表为例，中文约 1.5--2 字/token，而英文约 0.25 词/token，同样的上下文窗口里中文可容纳的信息量因此翻倍级差异；中文语料还面临“清洗后存量明显小于英文”的事实，这推动国内团队更早转向合成数据与多语言混合训练。

\textbf{训练目标}：Next Token Prediction（因果语言建模），即
\[
\max_\theta \; \mathbb{E}_{(x_1,\ldots,x_T)} \left[ \sum_{t=1}^{T} \log P_\theta(x_t \mid x_{<t}) \right]
\]

\subsubsection{分布式训练策略}
大规模 LLM 训练需要数千张 GPU 协同工作，核心策略包括：

\begin{table}[htbp]
\centering
\small
\begin{tabular}{@{}lll@{}}
\toprule
\textbf{策略} & \textbf{原理} & \textbf{典型框架}\\
\midrule
\textbf{数据并行 (DP)} & 每张卡独立梯度，AllReduce 同步 & PyTorch DDP、FSDP\\
\textbf{张量并行 (TP)} & 单个矩阵拆到多卡 & Megatron-LM\\
\textbf{流水线并行 (PP)} & 不同层放在不同卡上 & GPipe / 1F1B 调度\\
\textbf{序列并行 (SP)} & 长序列拆分 & Ring Attention\\
\textbf{ZeRO 优化} & 分片优化器状态/梯度/参数 & DeepSpeed ZeRO-1/2/3\\
\textbf{3D/4D/5D 混合} & DP + TP + PP (+ SP + EP) & 生产级训练标配\\
\bottomrule
\end{tabular}
\caption{分布式训练策略}
\end{table}

\textbf{硬件与软件栈}：
\begin{itemize}
  \item \textbf{硬件}：NVIDIA H100/H200/B200、Google TPU v5p/v6、自研 ASIC
  \item \textbf{通信}：NCCL、NVLink、GPUDirect RDMA
  \item \textbf{软件}：PyTorch DDP / FSDP / DeepSpeed / Megatron-LM / JAX
\end{itemize}

\textbf{如何组合这些并行维度}是大规模训练的核心设计决策，判据是\emph{通信频率与拓扑匹配}：
\begin{itemize}
  \item \textbf{张量并行}通信最频繁（每层前向/反向各需一次 allreduce/allgather），必须放在 NVLink 域内（单机 8 卡）；
  \item \textbf{流水线并行}只在 micro-batch 边界通信，通信量小，适合跨节点（InfiniBand/RoCE）；气泡（bubble）用 1F1B + 足够 micro-batch 数压制；
  \item \textbf{数据并行}每步一次梯度同步，放最外层横向扩展吞吐；ZeRO/FSDP 在其内部进一步分片显存；
  \item \textbf{专家并行}（MoE 专用）按专家分片，配合 all-to-all 通信，是第 5 章 MoE 训练的必备维度。
\end{itemize}
生产级训练典型拓扑是：\emph{节点内 TP=8 $\times$ 节点间 PP=4--8 $\times$ 最外层 DP/ZeRO}，即 3D/4D/5D 混合并行。以 70B 模型为例，BF16 参数约 140GB，Adam 状态再乘 3 倍以上——单卡 80GB 远放不下，TP 切参数 + ZeRO 切状态是标准解法。

\textbf{2025 年的效率标杆}：DeepSeek-V3 在 2048 块 H800 上以 \emph{DualPipe} 双向流水线（前向/反向计算与 all-to-all 通信深度重叠，几乎消除流水线气泡）、FP8 混合精度与节点受限路由的组合，把 14.8T token 的训练压缩到约 2.79M GPU$\cdot$hours——印证了\emph{架构-系统协同设计}比单纯堆卡更能决定训练经济学。

\subsubsection{典型超参数}（GPT-4 级别参考）
\begin{itemize}
  \item 学习率：$\sim 3 \times 10^{-4}$（cosine decay to $1 \times 10^{-5}$）
  \item Batch size：$\sim 4\text{M}$ tokens（通过梯度累积实现）
  \item 训练步数：数十万步量级（如 Llama 2 70B 以约 4M token 的全局 batch 训练约 2T token，对应数十万步）
  \item 精度：BF16 前向/反向 + FP32 主权重
  \item Warmup：2000--10000 步
  \item Weight decay：$0.1$
  \item Gradient clip：$\|g\| \leq 1.0$
\end{itemize}

\subsubsection{容错与检查点}
训练 100B+ 参数模型通常需要 3--12 个月，\textbf{容错机制}至关重要：
\begin{itemize}
  \item \textbf{分片式 checkpoint}：每个 rank 只持久化本地分片（参数、优化器状态、数据加载器游标），恢复时 AllGather 重建全局状态，避免全量镜像的单点与带宽瓶颈
  \item 故障检测与自动恢复（heartbeat + 重试）
  \item 弹性训练：节点故障时自动重新分配
  \item 典型频率：每 100--500 步保存一次 checkpoint
\end{itemize}

\subsection{阶段二：监督微调（SFT）}

\textbf{目标}：从“补全文本”转向“遵循指令、回答问题”。

\textbf{数据构造}：
\begin{itemize}
  \item 高质量指令-回答对（数千到数十万条）
  \item 来源：人工标注、模型蒸馏、规则生成
  \item 覆盖：问答、代码、翻译、数学、写作、推理
\textbf{核心原则}：数据质量 $\gg$ 数量（10K 高质量 $>$ 1M 低质量）
\end{itemize}

\textbf{SFT 到底教了模型什么}是一个值得澄清的问题：主流研究（如 LIMA）支持“SFT 是对齐/格式化，而非灌输知识”——能力几乎全部来自预训练，SFT 教会的是\emph{调用能力的接口}（听懂指令、组织答案、切换角色）。因此 SFT 数据的关键不是数量而是\emph{多样性-质量-难度}的三角平衡：覆盖足够多的任务形态（防止“会 A 不会 B”）、每条都经人工或强模型把关（防止学坏习惯）、并保留足够比例的复杂推理样本（防止能力退化）。2024 年后，SFT 数据中合成数据的占比持续上升（模型蒸馏 + 人工审核的流水线），“用强模型造数据、以人工守边界”成为标准做法；但纯合成数据的“分布收窄”风险（模型输出越来越像训练它的模型）也促使各家保留真实用户query与专业作者样本。

\textbf{自动数据构造技术}：
\begin{itemize}
  \item \textbf{Self-Instruct}：让模型自己生成指令
  \item \textbf{Evol-Instruct}：迭代进化指令复杂度（简单 $\to$ 复杂）
  \item \textbf{拒绝采样}：生成 $k$ 个回答，选最好的
\end{itemize}

\textbf{训练配置}：
\begin{itemize}
  \item 学习率：$10^{-5} \sim 5 \times 10^{-5}$（比预训练低 10--100 倍）
  \item Epoch：1--3
  \item 损失只计算\textbf{回答部分}（prompt 部分 mask 掉）
  \item 可选：LoRA 低秩适配（$r=16\text{--}64$，仅训练 $\sim 1\%$ 参数）
\end{itemize}

\textbf{LoRA 的数学形式}：冻结原权重 $W\in\mathbb{R}^{d\times k}$，旁路一个低秩增量
\[
W' = W + \Delta W = W + B A, \qquad B\in\mathbb{R}^{d\times r},\; A\in\mathbb{R}^{r\times k},\; r \ll \min(d,k)
\]
可训练参数量从 $dk$ 降至 $r(d+k)$（$r=16$ 时约为原来的 $\frac{1}{20}$）。经验上 LoRA 应加在注意力的 $W^Q, W^V$（甚至 $W^K, W^O$）而非 FFN 上效果更稳。其工程价值在于：\emph{同一基座可维护多个低成本领域适配器}（医疗、法律、代码），且推理时可合并权重、零额外延迟。QLoRA 进一步把基座量化到 4-bit（NF4）+ 分页优化器，使 65B 模型可在单张 48GB 卡上微调——这是个人研究者参与大模型微调的起点。

\subsection{阶段三：人类对齐（Alignment）}

\subsubsection{RLHF：基于人类反馈的强化学习}

经典三组件架构：
\begin{center}
\texttt{SFT 模型 (Policy) $\longrightarrow$ Reward Model (人类偏好训练) $\longrightarrow$ PPO Agent (策略优化)}
\end{center}

\textbf{步骤 1：收集偏好数据}

给定同一个 prompt $x$，模型生成 $k$ 个回答 $\{y_1, \ldots, y_k\}$，人类排序选出最优 $y_w$ 和最差 $y_l$。

\textbf{步骤 2：训练奖励模型}

使用 Bradley-Terry 模型：
\[
P(y_w \succ y_l \mid x) = \sigma\!\left(r_\phi(x, y_w) - r_\phi(x, y_l)\right)
\]
损失函数：
\[
\mathcal{L}_{\text{RM}} = -\mathbb{E}_{(x, y_w, y_l) \sim \mathcal{D}} \left[ \log \sigma\!\left(r_\phi(x, y_w) - r_\phi(x, y_l)\right) \right]
\]

\textbf{步骤 3：PPO 优化策略}

\[
\max_\pi \; \mathbb{E}_{x \sim \mathcal{D},\; y \sim \pi_\theta(y|x)} \left[ r_\phi(x,y) - \beta \cdot \text{KL}\!\left(\pi_\theta \;\|\; \pi_{\text{ref}}\right) \right]
\]
\begin{itemize}
  \item $r_\phi$：奖励模型给出的分数
  \item $\beta$：KL 约束强度（典型值 $0.01\text{--}0.2$）
  \item $\pi_{\text{ref}}$：SFT 模型（防止偏离太远）
\end{itemize}

\begin{warnbox}{RLHF 的工程挑战}
\begin{enumerate}
  \item 需要\textbf{同时维护 4 个模型}：Policy、Reference、Reward、Value Critic
  \item 显存需求是 SFT 的 \textbf{4 倍}
  \item 训练不稳定，超参数（学习率、$\beta$、GAE $\lambda$）非常敏感
  \item 奖励模型容易被“hack”：模型生成极端文本骗取高分
  \item 离线偏好数据存在分布偏移
\end{enumerate}
\end{warnbox}

\subsubsection{DPO：直接偏好优化}

2023 年提出，\textbf{去掉了 Reward Model}，直接从偏好对优化策略：
\[
\mathcal{L}_{\text{DPO}} = -\mathbb{E}_{(x, y_w, y_l)} \left[ \log \sigma\!\left( \beta \log \frac{\pi_\theta(y_w|x)}{\pi_{\text{ref}}(y_w|x)} - \beta \log \frac{\pi_\theta(y_l|x)}{\pi_{\text{ref}}(y_l|x)} \right) \right]
\]

\textbf{DPO 的直觉}：注意 $\beta \log \frac{\pi_\theta(y|x)}{\pi_{\text{ref}}(y|x)}$ 恰好是 RLHF 目标下\emph{隐式奖励}的闭式解（RLHF 的最优策略满足 $\pi^*(y|x) \propto \pi_{\text{ref}}(y|x)\, e^{r(x,y)/\beta}$，反解即得 $r = \beta \log \frac{\pi^*}{\pi_{\text{ref}}} + \text{const}$）。因此 DPO 相当于把奖励模型“吸收”进策略网络：不再单独训练 RM，而是直接以“相对参考模型的隐式偏好差”为监督信号优化策略。代价是它只能利用\emph{离线偏好对}，无法在线采样——当分布偏移大时效果弱于在线 RL（如 GRPO），这也是 2025 年“在线 RL 回归”趋势的背景（见第 7 章 DeepSeek-R1）。

\begin{table}[htbp]
\centering
\small
\begin{tabular}{@{}llll@{}}
\toprule
\textbf{方法} & \textbf{核心思想} & \textbf{优势} & \textbf{劣势}\\
\midrule
PPO & 经典 RL + RM & 灵活、表达力强 & 不稳定、显存大\\
\textbf{DPO} & 无 RM，直接优化 & 简单、稳定、省显存 & 表达力稍弱\\
IPO & DPO 改进 & 避免过拟合 & —\\
KTO & 只需二元反馈 & 不需成对数据 & 信息量较少\\
ORPO & SFT + 偏好两合一 & 简化流程 & —\\
SimPO & 长度归一化 & 减少长度偏差 & —\\
RLOO & 无 Critic 的 RL & 方差更低 & —\\
GRPO & 组相对策略优化 & DeepSeek 2025 & —\\
\bottomrule
\end{tabular}
\caption{对齐方法对比}
\end{table}

\subsubsection{Constitutional AI（Anthropic）}

用“宪法”（一组原则列表）指导模型\textbf{自我批评}：
\begin{enumerate}
  \item 模型生成回答
  \item 模型根据宪法原则自我批评
  \item 模型根据批评修改回答
  \item 用修改后的数据做 SFT
\end{enumerate}
\textbf{RLAIF}（RL from AI Feedback）进一步用 AI 反馈替代人类反馈，大幅降低成本。

\textbf{对齐的三条经验规律}值得记录。第一，\emph{对齐税}（alignment tax）真实存在但可控：过度保守的 RM 会让模型变得过度谨慎、拒答泛滥，工程上用“有用性/无害性”双奖励加权来平衡。第二，\emph{Goodhart 定律是 RLHF 的宿命}：RM 只是人类偏好的代理指标，优化久了必然出现 reward hacking（谄媚、冗长、堆砌格式），因此 RM 需要持续用新鲜偏好数据再训。第三，\emph{可验证奖励（RLVR）正在接管推理域}：数学有标准答案、代码能跑测试，这类任务直接用规则判分（第 7 章 R1 的做法），绕开了 RM 的全部毛病——“有标准答案的交给规则，没有的才交给人类”。

\subsection{持续学习与更新}

\begin{itemize}
  \item \textbf{灾难性遗忘}：新知识覆盖旧知识
  \item \textbf{解决方案}：混合新旧数据回放、LoRA 增量适配、弹性权重固化（EWC）
  \item \textbf{持续预训练（CPT）}：在特定领域数据上继续预训练
  \item \textbf{RAG（检索增强生成）}：外挂知识库，不修改权重
\end{itemize}

\textbf{四种“更新知识”手段的选型判据}：
\begin{itemize}
  \item \textbf{RAG}：知识更新快（改库即可）、可溯源、适合私有/高频变化知识；缺点是受检索质量制约，跨文档推理弱；
  \item \textbf{LoRA/CPT}：把领域\emph{风格与隐性知识}内化进权重，推理时无需检索、延迟低；代价是每更新一轮需重新训练与评估；
  \item \textbf{全量 CPT}：领域能力上限最高（如医疗、法律、代码专精模型），但成本高、易伤通用能力，需混入通用数据“回防”。
\end{itemize}
生产系统的常见组合是 \emph{RAG（事实层）+ SFT/LoRA（行为层）}：让模型“知道怎么做”靠微调，“知道当前事实”靠检索——两者正交、可叠加。

% ===================== 第五章 =====================
\section{前沿架构创新}

\subsection{混合专家（Mixture of Experts, MoE）}

\textbf{核心思想}：FFN 层拆分为 $N$ 个“专家”，每个 token 只激活 Top-$K$ 个专家，实现\textbf{稀疏激活}。

\begin{figure}[htbp]
\centering
\resizebox{0.97\textwidth}{!}{%
% fig07_moe_layer.tex — 图7 MoE 层工作机制
% 《深度学习与大语言模型：前沿技术全景报告》图示库
% 由 dl_llm_report.tex 抽取，经 % fig07_moe_layer.tex — 图7 MoE 层工作机制
% 《深度学习与大语言模型：前沿技术全景报告》图示库
% 由 dl_llm_report.tex 抽取，经 % fig07_moe_layer.tex — 图7 MoE 层工作机制
% 《深度学习与大语言模型：前沿技术全景报告》图示库
% 由 dl_llm_report.tex 抽取，经 \input{figs/fig07_moe_layer} 引用（caption/label 在主文件）
% 注意：本文件为片段，依赖主文件导言区的 tikz/pgfplots 宏包、颜色定义与 tikzset 样式，不可单独编译。

\begin{tikzpicture}
\node[gry, minimum width=20mm] (x) {输入 token $\mathbf{x}$};
\node[orn, right=7mm of x, minimum width=22mm, align=center] (r) {Router\\ 打分 + Top-$K$ 选择};
\foreach \i/\xa/\ya/\st in {
  1/2.55/1.35/gry, 2/3.95/1.35/orn, 3/2.55/0.45/gry, 4/3.95/0.45/gry,
  5/2.55/-0.45/orn, 6/3.95/-0.45/gry, 7/2.55/-1.35/gry, 8/3.95/-1.35/gry} {
  \node[\st, minimum width=12mm, minimum height=7.5mm, font=\small] (e\i) at ($(r.east)+(\xa,\ya)$) {专家\i};
}
\node[grn, minimum width=24mm, font=\small] (sh) at ($(r.east)+(3.25,2.45)$) {共享专家（恒激活）};
\node[gry, minimum width=29mm, align=center] (agg) at ($(r.east)+(6.7,0)$) {加权合并\\ $\sum_{i\in\text{Top-}K}\alpha_i E_i(\mathbf{x})$};
\node[gry, minimum width=15mm, right=4mm of agg] (o) {输出};
\foreach \i in {1,3,4,6,7,8} \draw[cgray!55, thin, dotted] (r.east) -- (e\i.west);
\draw[rrr] (r.east) -- (e2.west);
\draw[rrr] (r.east) -- (e5.west);
\draw[brr] (r.east) -- ++(0.9cm,0) |- (sh.west);
\draw[rrr] (e2.east) -- (agg.west);
\draw[rrr] (e5.east) -- (agg.west);
\draw[brr] (sh.east) -- (agg.north west);
\draw[arr] (x)--(r); \draw[arr] (agg)--(o);
\node[lab, below=3.5mm of e7, xshift=5mm, align=center] {示意：8 选 2 + 1 共享；\\ DeepSeek-V3 实际为 256 选 8 + 1 共享\\ （Qwen3 为 128 选 8、无共享专家）};
\end{tikzpicture}
 引用（caption/label 在主文件）
% 注意：本文件为片段，依赖主文件导言区的 tikz/pgfplots 宏包、颜色定义与 tikzset 样式，不可单独编译。

\begin{tikzpicture}
\node[gry, minimum width=20mm] (x) {输入 token $\mathbf{x}$};
\node[orn, right=7mm of x, minimum width=22mm, align=center] (r) {Router\\ 打分 + Top-$K$ 选择};
\foreach \i/\xa/\ya/\st in {
  1/2.55/1.35/gry, 2/3.95/1.35/orn, 3/2.55/0.45/gry, 4/3.95/0.45/gry,
  5/2.55/-0.45/orn, 6/3.95/-0.45/gry, 7/2.55/-1.35/gry, 8/3.95/-1.35/gry} {
  \node[\st, minimum width=12mm, minimum height=7.5mm, font=\small] (e\i) at ($(r.east)+(\xa,\ya)$) {专家\i};
}
\node[grn, minimum width=24mm, font=\small] (sh) at ($(r.east)+(3.25,2.45)$) {共享专家（恒激活）};
\node[gry, minimum width=29mm, align=center] (agg) at ($(r.east)+(6.7,0)$) {加权合并\\ $\sum_{i\in\text{Top-}K}\alpha_i E_i(\mathbf{x})$};
\node[gry, minimum width=15mm, right=4mm of agg] (o) {输出};
\foreach \i in {1,3,4,6,7,8} \draw[cgray!55, thin, dotted] (r.east) -- (e\i.west);
\draw[rrr] (r.east) -- (e2.west);
\draw[rrr] (r.east) -- (e5.west);
\draw[brr] (r.east) -- ++(0.9cm,0) |- (sh.west);
\draw[rrr] (e2.east) -- (agg.west);
\draw[rrr] (e5.east) -- (agg.west);
\draw[brr] (sh.east) -- (agg.north west);
\draw[arr] (x)--(r); \draw[arr] (agg)--(o);
\node[lab, below=3.5mm of e7, xshift=5mm, align=center] {示意：8 选 2 + 1 共享；\\ DeepSeek-V3 实际为 256 选 8 + 1 共享\\ （Qwen3 为 128 选 8、无共享专家）};
\end{tikzpicture}
 引用（caption/label 在主文件）
% 注意：本文件为片段，依赖主文件导言区的 tikz/pgfplots 宏包、颜色定义与 tikzset 样式，不可单独编译。

\begin{tikzpicture}
\node[gry, minimum width=20mm] (x) {输入 token $\mathbf{x}$};
\node[orn, right=7mm of x, minimum width=22mm, align=center] (r) {Router\\ 打分 + Top-$K$ 选择};
\foreach \i/\xa/\ya/\st in {
  1/2.55/1.35/gry, 2/3.95/1.35/orn, 3/2.55/0.45/gry, 4/3.95/0.45/gry,
  5/2.55/-0.45/orn, 6/3.95/-0.45/gry, 7/2.55/-1.35/gry, 8/3.95/-1.35/gry} {
  \node[\st, minimum width=12mm, minimum height=7.5mm, font=\small] (e\i) at ($(r.east)+(\xa,\ya)$) {专家\i};
}
\node[grn, minimum width=24mm, font=\small] (sh) at ($(r.east)+(3.25,2.45)$) {共享专家（恒激活）};
\node[gry, minimum width=29mm, align=center] (agg) at ($(r.east)+(6.7,0)$) {加权合并\\ $\sum_{i\in\text{Top-}K}\alpha_i E_i(\mathbf{x})$};
\node[gry, minimum width=15mm, right=4mm of agg] (o) {输出};
\foreach \i in {1,3,4,6,7,8} \draw[cgray!55, thin, dotted] (r.east) -- (e\i.west);
\draw[rrr] (r.east) -- (e2.west);
\draw[rrr] (r.east) -- (e5.west);
\draw[brr] (r.east) -- ++(0.9cm,0) |- (sh.west);
\draw[rrr] (e2.east) -- (agg.west);
\draw[rrr] (e5.east) -- (agg.west);
\draw[brr] (sh.east) -- (agg.north west);
\draw[arr] (x)--(r); \draw[arr] (agg)--(o);
\node[lab, below=3.5mm of e7, xshift=5mm, align=center] {示意：8 选 2 + 1 共享；\\ DeepSeek-V3 实际为 256 选 8 + 1 共享\\ （Qwen3 为 128 选 8、无共享专家）};
\end{tikzpicture}

}%
\caption{MoE 层工作机制：Router 为每个 token 打分并选出 Top-$K$ 个专家，仅被选中的专家（橙色）参与计算，输出按门控权重 $\alpha_i$ 加权合并；共享专家每 token 恒激活，负责“通用知识”。}
\label{fig:moe}
\end{figure}

\textbf{细粒度专家 + 共享专家}的设计逻辑值得展开。传统 MoE（如 Mixtral）用 8 个大专家选 2 个，组合数有限（$C_8^2=28$）；DeepSeek-V3 把专家切细到 256 个、每个专家更小，选 8 个的组合数达 $10^{14}$ 量级——\emph{同样的计算预算下，知识可以按更细的粒度组织}（“数学-积分-符号处理”这样的专项专家成为可能）；而 1 个共享专家恒定激活，承担所有 token 都需要的通用计算，让路由专家可以更“偏科”。代价同样明显：\emph{全部专家必须常驻显存}（671B 参数即使只激活 37B，加载与分布式缓存的仍是全量），MoE 推理服务对显存容量（而非算力）更敏感，这是 MoE 部署经济学与稠密模型的本质差异；训练侧则引入\emph{专家并行}（EP）与 all-to-all 通信，是 4D/5D 并行中最难调的一维。

\[
\text{MoE}(\mathbf{x}) = \sum_{i \in \text{Top-}K} \alpha_i \cdot \text{Expert}_i(\mathbf{x}), \qquad
\alpha_i = \text{softmax}(\text{Router}(\mathbf{x}))_i
\]

\begin{table}[htbp]
\centering
\small
\begin{tabular}{@{}llll@{}}
\toprule
\textbf{模型} & \textbf{总参数} & \textbf{激活参数} & \textbf{专家结构}\\
\midrule
Mixtral 8$\times$7B & 46.7B & 12.9B & 8 专家 / 激活 2\\
Qwen3-235B-A22B & 235B & 22B & 128 专家 / 激活 8（无共享专家，KV 头 4）\\
DeepSeek-V3 & 671B & 37B & 256 路由专家 / 激活 8（+1 共享专家）\\
Llama 4 Scout & $\sim$109B & 17B & 16 专家 / 激活 1\\
Kimi K2 & $\sim$1T & 32B & 384 专家 / 激活 8（+1 共享专家）\\
GPT-4（估计） & 1T+ & 100--200B & 未公开\\
\bottomrule
\end{tabular}
\caption{MoE 模型参数规模（总参数决定“知识容量”，激活参数决定单 token 计算成本）}
\end{table}

\textbf{为什么 MoE 成为旗舰模型主流架构}：核心是\emph{解耦“知识容量”与“单 token 计算量”}。稠密模型中，参数规模直接绑定推理成本；MoE 让参数量（可存储的知识）继续扩张，而每个 token 只付费激活的 $K$ 个专家。DeepSeek-V3 把这一点推到极致：671B 总参数中仅 37B 参与前向，训练成本约 \textbf{2.79M H800 GPU$\cdot$hours}（折合约 \textbf{\$5.5M}）——以 2024 年底的价格训练出对标 GPT-4 级别的模型，直接改写了大模型的经济学。

\textbf{路由策略}：Top-$K$ 路由（$K=2$ 最常见）、Noisy Top-K（加噪声鼓励探索）、Expert Choice（专家选 token）。

\textbf{谱系}：MoE 并非新事物——2017 年 Shazeer 等的稀疏门控 LSTM 层即其雏形；GShard（2020）与 Switch Transformer（2021，1.6T 参数）首次将其带入 Transformer 大规模训练；Mixtral 8$\times$7B（2024）证明 MoE 可以兼顾开源与工程友好；DeepSeek-V3、Qwen3 与 Kimi K2（2025）则把“细粒度专家 + 共享专家 + 无辅助损失/精细负载均衡”打磨成旗舰标配。

\textbf{负载均衡}辅助损失（Switch Transformer 形式）：$\mathcal{L}_{\text{aux}} = \alpha \cdot N \cdot \sum_{j} f_j \cdot P_j$，其中 $f_j$ 是第 $j$ 个专家被选中的 token 频率、$P_j$ 是路由器分配给它的平均概率——最小化该损失会促使负载 $f_j \to 1/N$ 均匀。

\textbf{DeepSeek-V3 的 MoE 创新}（细粒度专家 + 共享专家 + 无辅助损失均衡）：
\begin{itemize}
  \item \textbf{无辅助损失负载均衡}：为每个专家引入偏置项 $b_i$，它只参与 Top-$K$ 路由的\emph{选择}、不进入门控加权值，并按负载动态增减——在不引入“为平衡而牺牲质量”的梯度干扰的前提下实现均衡；辅以序列级互补辅助损失与\textbf{节点受限路由}（每个 token 的专家限定在至多若干节点内，压低 all-to-all 通信量）；
  \item \textbf{Multi-Token Prediction (MTP)}：附加一个轻量 MTP 模块同时预测下下个 token，密集训练信号改善表征；推理时该模块可整体丢弃，或直接充当\emph{自草拟}的推测解码器（见第 6 章）；
  \item \textbf{训练成本}：约 2.79M H800 GPU$\cdot$hours（约 \$5.6M，2024 年 12 月技术报告口径），且 \textbf{FP8 混合精度}首次在 671B 规模全程使用。
\end{itemize}

\subsection{注意力机制变体}

\subsubsection{GQA（Grouped Query Attention）}

$Q$ 有 $h$ 个头，KV 有 $g$ 个头（$g < h$），多个 Q 头共享一个 KV 头。KV Cache 减少 $h/g$ 倍（典型 4--8$\times$），质量几乎不降。用于 LLaMA 2 70B、Mistral 7B、Gemma 2。

\textbf{算一笔显存账}：KV Cache 大小 $= 2(\text{K/V}) \times L \times n_{\text{kv}} \times d_h \times T \times \text{batch} \times \text{bytes}$。以 LLaMA-2 70B 为例（80 层、64 个 Q 头、每头 128 维、FP16），单条 4K 上下文序列的 KV Cache 为 $2 \times 80 \times 64 \times 128 \times 4096 \times 2\,\text{B} \approx 10.7\,\text{GB}$——batch 32 时约 343GB，远超模型权重本身（约 140GB）。GQA 将 KV 头从 64 降到 8，同条件下降到约 43GB（单序列约 1.3GB），恰好回到“可服务”量级。LLaMA-2 技术报告的结论是：70B 上 GQA-8 的建模质量与 MHA 基本持平，而解码吞吐向 MQA 看齐——这一“近乎免费”的压缩是 GQA 成为 2024 年后几乎所有新模型默认配置的原因。

\subsubsection{MLA（Multi-head Latent Attention）}

DeepSeek-V2/V3 创新，将 K、V 联合压缩到低维隐空间：
\[
\hat{K} = W_{KV}^D \cdot (W^C \cdot \mathbf{x}), \qquad
\hat{V} = W_{KV}^U \cdot \hat{K}
\]
KV Cache 压缩比比 GQA 高 \textbf{5--10 倍}。关键洞察：K 和 V 信息冗余度高，可共享压缩空间。

\textbf{直觉展开}：MLA 观察到，对推理成本起决定作用的不是“KV 有多少个头”，而是“每个 token 需要缓存多少维”。于是把 $(K, V)$ 联合压缩到一个低维隐向量 $c_{KV} \in \mathbb{R}^{d_c}$（DeepSeek-V2 中 $d_c=512$，而 LLaMA2-70B 每 token 的 KV 是 $64 \times 2 \times 128 = 16384$ 维），缓存 $c_{KV}$（512 维）加少量解耦 RoPE 键（64 维）替代缓存全部 K/V——DeepSeek-V2 官方报告生成阶段 KV Cache \textbf{降低 93.3\%}，推理成本同步大幅下降。位置信息则通过独立的“解耦 RoPE”注入，避免压缩破坏旋转位置结构。这一设计证明了\emph{架构级省显存}比\emph{系统级省显存}（量化、分页）更根本——它改变了“多大批量可服务”的经济学上限。

\subsubsection{滑动窗口注意力}

Mistral 7B：窗口大小 4096，每 4 层一次全局注意力。显存 $O(\text{window})$，长序列推理友好，但窗口外信息丢失。

\subsubsection{线性注意力与状态空间模型}

\textbf{Linear Transformer}：用核函数 $\phi$ 近似 softmax，复杂度 $O(n d^2)$。

\textbf{Mamba / Mamba-2}：选择性状态空间模型（SSM），$O(n)$ 复杂度：
\[
h_t = \bar{A}\, h_{t-1} + \bar{B}\, x_t, \qquad y_t = C\, h_t
\]
长序列推理极快，但表达能力可能不如标准注意力。2025 趋势：Mamba-Transformer 混合架构。

\textbf{Mamba 的关键创新是“选择性”（selectivity）}：早期线性 SSM（如 S4）的状态转移矩阵 $A, B, C$ 是与输入无关的常量，等价于一个固定滤波器，无法实现“看内容决定记什么”。Mamba 让 $A, B, C$ 成为当前输入 $x_t$ 的函数（选择性扫描，selective scan），并以硬件感知的并行 scan 算法实现——兼顾 RNN 的 $O(1)$ 递推推理与 GPU 上的并行训练。实际部署中的共识是\emph{混合架构}：用少量注意力层负责“精确检索”（注意力是精确查找表，SSM 是压缩记忆），如 Jamba、Zamba 以及 2025 年 NVIDIA Nemotron-H 等“注意力 + SSM 交替”设计，在长上下文成本与检索精度之间取得平衡。

回看本节，KV 侧的三代优化构成一条清晰的演进线：\emph{MQA}（全部 Q 头共享 1 个 KV 头，省到极致但质量受损）$\to$ \emph{GQA}（分组共享，质量与开销折中）$\to$ \emph{MLA}（把 KV 联合压缩进低维隐空间，本质是对 KV 做“表示学习”）。每一代都在回答同一个问题——\textbf{“KV 里到底有多少信息是真正必要的？”}——并给出更激进的答案；配合量化与分页（第 6 章），架构级与系统级的省显存手段相互正交、可以叠加。

\subsection{代表性架构创新}

\begin{table}[htbp]
\centering
\small
\begin{tabular}{@{}lp{5.6cm}p{4.6cm}@{}}
\toprule
\textbf{模型} & \textbf{核心创新} & \textbf{亮点}\\
\midrule
DeepSeek-V2 & MLA + 细粒度 MoE & KV Cache 压缩 93.3\%\\
DeepSeek-V3 & 无辅助损失 MoE + MTP + FP8 & 约 \$5.6M 训练 671B\\
DeepSeek-R1 & 纯 RL 推理 + GRPO & 无 CoT 标注涌现推理\\
Kimi K2 & 1T 级 MoE 全量开源 & 开源模型榜第一（LMArena）\\
Qwen3-235B-A22B & MoE + 混合思考/思考预算 & 128 专家/激活 8，Apache 2.0\\
Gemini 3 Pro & 原生多模态 + Deep Think & LMArena 首破 1500 Elo\\
Gemini 3.7（2026，据报道） & Agentic 多 Agent & 多 Agent 数学/工程\\
Llama 4 & 大规模 MoE + 长上下文 & 开源旗舰（评测有争议）\\
\bottomrule
\end{tabular}
\caption{2024--2026 代表性架构创新}
\end{table}

\textbf{架构创新的一般方法论}：先在小模型上验证（数天内完成对比实验）、再放大到旗舰规模——Chinchilla 与 DeepSeek 的所有架构决策都遵循这一“小模型先导”原则，这也是 $\mu$P 类超参迁移技术（第 10 章）的价值所在。评价一项架构创新，看三件事：\emph{帕累托改进}（同质量下成本是否下降，而非仅“分数是否更高”）、\emph{硬件亲和性}（是否匹配显存带宽/互联拓扑的约束）、\emph{工程可落地性}（推理引擎多快能支持）。2024--2026 年的赢家——MoE、GQA/MLA、MTP、FP8——全部是这三条的优等生；而一些“分数好看”但硬件不友好的设计（如某些双向注意力变体）则停留在论文阶段。

% ===================== 第六章 =====================
\section{推理优化}

\subsection{推理瓶颈分析}

LLM 推理分两个阶段，瓶颈完全不同：

\begin{table}[htbp]
\centering
\small
\begin{tabular}{@{}llll@{}}
\toprule
\textbf{阶段} & \textbf{特征} & \textbf{瓶颈类型} & \textbf{优化方向}\\
\midrule
Prefill & 处理整个 prompt & 计算密集 & FlashAttention、算子融合\\
Decode & 逐 token 生成 & 访存密集 & KV Cache 优化、推测解码\\
\bottomrule
\end{tabular}
\caption{推理两阶段瓶颈}
\end{table}

\textbf{分析工具：算术强度}（Arithmetic Intensity, FLOPs/Byte）。Prefill 阶段，一个长度为 $n$ 的 prompt 前向时，权重只需读取一次却参与 $n$ 次运算，算术强度 $\propto n$，轻松超过硬件的“计算/访存平衡点”（如 H100 BF16 约 $989\,\text{TFLOPS} / 3.35\,\text{TB/s} \approx 295$ FLOP/Byte），属于\emph{计算受限}；Decode 阶段每个 token 都要把全部权重重新读一遍，算术强度 $\approx 1$（权重为主时），远低于平衡点，属于\emph{访存受限}。这一判据直接推出两条优化主线：\emph{Prefill 拼算力}（FlashAttention、算子融合、FP8），\emph{Decode 拼带宽}（KV Cache 压缩、量化、推测解码“一次读权重出多个 token”）。

\begin{figure}[htbp]
\centering
% fig08_inference_system.tex — 图8 LLM 推理系统全景
% 《深度学习与大语言模型：前沿技术全景报告》图示库
% 由 dl_llm_report.tex 抽取，经 % fig08_inference_system.tex — 图8 LLM 推理系统全景
% 《深度学习与大语言模型：前沿技术全景报告》图示库
% 由 dl_llm_report.tex 抽取，经 % fig08_inference_system.tex — 图8 LLM 推理系统全景
% 《深度学习与大语言模型：前沿技术全景报告》图示库
% 由 dl_llm_report.tex 抽取，经 \input{figs/fig08_inference_system} 引用（caption/label 在主文件）
% 注意：本文件为片段，依赖主文件导言区的 tikz/pgfplots 宏包、颜色定义与 tikzset 样式，不可单独编译。

\begin{tikzpicture}[node distance=6mm]
\node[gry, minimum width=18mm] (req) {用户请求};
\node[blk, right=6mm of req, minimum width=27mm] (pre) {Prefill 预填充\\ 整段 prompt 并行};
\node[orn, right=6mm of pre, minimum width=20mm] (kv) {KV Cache};
\node[grn, right=6mm of kv, minimum width=27mm] (dec) {Decode 解码\\ 逐 token 自回归};
\node[gry, right=6mm of dec, minimum width=14mm] (res) {响应};
\draw[arr] (req)--(pre); \draw[arr] (pre)--(kv); \draw[arr] (kv)--(dec); \draw[arr] (dec)--(res);
\draw[rrr] (dec.north) to[bend right=22] node[lab, above, pos=0.5] {每步读写 KV（访存受限）} (kv.north);
\node[lab, below=3mm of pre] {计算受限 · 拼算力：FlashAttention、算子融合、FP8};
\node[lab, below=3mm of dec] {访存受限 · 拼带宽：量化、投机解码};
\node[gry, minimum width=126mm, minimum height=9mm, below=17mm of kv, font=\small] (svc) {服务层：Continuous Batching · PagedAttention · Prefix 复用 · 量化 · 投机解码};
\node[lab, below=1.5mm of svc] {所有请求共享同一套引擎与显存池，调度策略决定吞吐与延迟的平衡};
\end{tikzpicture}
 引用（caption/label 在主文件）
% 注意：本文件为片段，依赖主文件导言区的 tikz/pgfplots 宏包、颜色定义与 tikzset 样式，不可单独编译。

\begin{tikzpicture}[node distance=6mm]
\node[gry, minimum width=18mm] (req) {用户请求};
\node[blk, right=6mm of req, minimum width=27mm] (pre) {Prefill 预填充\\ 整段 prompt 并行};
\node[orn, right=6mm of pre, minimum width=20mm] (kv) {KV Cache};
\node[grn, right=6mm of kv, minimum width=27mm] (dec) {Decode 解码\\ 逐 token 自回归};
\node[gry, right=6mm of dec, minimum width=14mm] (res) {响应};
\draw[arr] (req)--(pre); \draw[arr] (pre)--(kv); \draw[arr] (kv)--(dec); \draw[arr] (dec)--(res);
\draw[rrr] (dec.north) to[bend right=22] node[lab, above, pos=0.5] {每步读写 KV（访存受限）} (kv.north);
\node[lab, below=3mm of pre] {计算受限 · 拼算力：FlashAttention、算子融合、FP8};
\node[lab, below=3mm of dec] {访存受限 · 拼带宽：量化、投机解码};
\node[gry, minimum width=126mm, minimum height=9mm, below=17mm of kv, font=\small] (svc) {服务层：Continuous Batching · PagedAttention · Prefix 复用 · 量化 · 投机解码};
\node[lab, below=1.5mm of svc] {所有请求共享同一套引擎与显存池，调度策略决定吞吐与延迟的平衡};
\end{tikzpicture}
 引用（caption/label 在主文件）
% 注意：本文件为片段，依赖主文件导言区的 tikz/pgfplots 宏包、颜色定义与 tikzset 样式，不可单独编译。

\begin{tikzpicture}[node distance=6mm]
\node[gry, minimum width=18mm] (req) {用户请求};
\node[blk, right=6mm of req, minimum width=27mm] (pre) {Prefill 预填充\\ 整段 prompt 并行};
\node[orn, right=6mm of pre, minimum width=20mm] (kv) {KV Cache};
\node[grn, right=6mm of kv, minimum width=27mm] (dec) {Decode 解码\\ 逐 token 自回归};
\node[gry, right=6mm of dec, minimum width=14mm] (res) {响应};
\draw[arr] (req)--(pre); \draw[arr] (pre)--(kv); \draw[arr] (kv)--(dec); \draw[arr] (dec)--(res);
\draw[rrr] (dec.north) to[bend right=22] node[lab, above, pos=0.5] {每步读写 KV（访存受限）} (kv.north);
\node[lab, below=3mm of pre] {计算受限 · 拼算力：FlashAttention、算子融合、FP8};
\node[lab, below=3mm of dec] {访存受限 · 拼带宽：量化、投机解码};
\node[gry, minimum width=126mm, minimum height=9mm, below=17mm of kv, font=\small] (svc) {服务层：Continuous Batching · PagedAttention · Prefix 复用 · 量化 · 投机解码};
\node[lab, below=1.5mm of svc] {所有请求共享同一套引擎与显存池，调度策略决定吞吐与延迟的平衡};
\end{tikzpicture}

\caption{LLM 推理系统全景：Prefill 与 Decode 瓶颈不同、优化路径不同，中间的 KV Cache 是两条路径的交汇点；底座是所有主流引擎共享的服务层技术栈（详见 6.2--6.6 节）。}
\label{fig:infer}
\end{figure}

\textbf{衡量推理服务的四个指标}是后续所有优化的靶子：\emph{TTFT}（Time To First Token，首 token 延迟，主要由 Prefill 决定，交互体验的第一印象）、\emph{TPOT}（Time Per Output Token，单 token 生成间隔，主要由 Decode 决定，决定“打字机”的流畅度）、\emph{吞吐量}（tokens/s，决定单位算力成本）与\emph{并发容量}（同时服务的请求数，由 KV Cache 显存余量决定）。四者相互牵制——批越大吞吐越高，但 TPOT 与 TTFT 随之变差；KV Cache 压得越省并发越高，但可能牺牲质量。推理系统的全部艺术，就是按照业务 SLO（如“TTFT $<$ 400ms、TPOT $<$ 30ms”）在这四个数字之间找最优工作点。

\subsection{KV Cache 优化}

显存占用：$\text{Size} = 2 \times L \times n_{\text{kv}} \times d_{\text{kv}} \times n_{\text{batch}} \times n_{\text{bytes}}$

\textbf{优化手段}：
\begin{itemize}
  \item \textbf{GQA / MLA}：减少 KV 头数（4--10$\times$ 压缩）
  \item \textbf{PagedAttention (vLLM)}：分页管理，减少碎片
  \item \textbf{KV Cache 量化}：FP16 $\to$ INT8 $\to$ INT4
  \item \textbf{Sliding Window}：只保留最近窗口
  \item \textbf{H$_2$O / StreamingLLM}：保留高注意力 token
  \item \textbf{Prefix Caching}：共享 system prompt 的 KV Cache
\end{itemize}

\subsection{量化技术}

\begin{table}[htbp]
\centering
\small
\begin{tabular}{@{}llll@{}}
\toprule
\textbf{精度} & \textbf{显存占比} & \textbf{质量损失} & \textbf{场景}\\
\midrule
FP16 / BF16 & 100\% & 无 & 服务器\\
FP8 & 50\% & $<$0.5\% & H100 原生支持\\
INT8 & 50\% & $<$1\% & 边缘部署\\
INT4 & 25\% & 1--3\% & 移动端\\
\bottomrule
\end{tabular}
\caption{量化精度对比}
\end{table}

\textbf{主流量化算法}及其思路：
\begin{itemize}
  \item \textbf{GPTQ}：基于 Hessian 二阶信息的\emph{离线逐层最优量化}——逐列量化后用已量化列修正剩余列，误差累积可控；
  \item \textbf{AWQ}：\emph{激活感知}——统计输入激活幅度，对“显著通道”施加保护性缩放再量化，避免异常值（outlier）主导误差；
  \item \textbf{SmoothQuant}：针对\emph{激活值}的量化难题，把激活的异常难度“迁移”一部分到权重上（逐通道等价变换），使 W8A8 同时可行；
  \item \textbf{GGUF / llama.cpp}：面向 CPU/混合推理的分块量化格式，支持 2--8 bit 混合精度与 Q4/Q5 档位，是端侧生态的事实标准。
\end{itemize}
选择经验：服务端大流量场景优先 FP8/W8A8（质量损失 $<$1\%，且 H100 原生支持、算力翻倍）；边缘/移动端用 INT4/AWQ；CPU 用 GGUF。

\subsection{推测解码（Speculative Decoding）}

小模型快速草拟 $k$ 个 token，大模型一次验证：
\[
\alpha = \min\!\left(1,\; \frac{P_{\text{target}}(t \mid \text{ctx})}{P_{\text{draft}}(t \mid \text{ctx})}\right)
\]
\textbf{为什么质量无损}：上面的 $\alpha$ 正是\emph{接受-拒绝采样}的接受概率——被拒绝的 token 按修正分布 $\propto \max(0, P_{\text{target}} - P_{\text{draft}})$ 重新采样，可以严格证明整个过程的输出分布与直接对目标模型采样\textbf{完全一致}。这就是推测解码“免费加速”的数学保证：加速来自“一次权重读取验证多个 token”，而非任何近似。

工程上期望加速比 2--3$\times$（取决于草拟模型与目标模型的分布接近程度）。草拟模型的选择有三条路线：\emph{小模型}（经典方案）、\emph{同模型多头}（Medusa 在主干上加解码头并行猜多 token、Self-Draft）、\emph{特征级自草拟}（EAGLE/EAGLE-2 利用倒数第二层隐状态直接生成候选，2025 年主流）。变体还有 Lookahead Decoding（ Jacobi 迭代，无需任何草拟模型）。

\subsection{推理引擎与系统优化}

\begin{table}[htbp]
\centering
\small
\begin{tabular}{@{}ll@{}}
\toprule
\textbf{引擎} & \textbf{核心特点}\\
\midrule
\textbf{vLLM} & PagedAttention + 连续批处理\\
TensorRT-LLM & NVIDIA 专用，极致优化\\
SGLang & 结构化输出加速、RadixAttention\\
llama.cpp & CPU/边缘推理，GGUF 格式\\
MLX & Apple Silicon 优化\\
\bottomrule
\end{tabular}
\caption{主流推理引擎}
\end{table}

\textbf{关键系统优化}：
\begin{itemize}
  \item \textbf{Continuous Batching}：请求完成即释放，新请求随时插入，吞吐提升 2--5$\times$
  \item \textbf{FlashAttention 3}：IO-aware 分块 + Hopper TMA/WGMMA，注意力 2--4$\times$ 加速
  \item \textbf{算子融合 / CudaGraph}：减少显存读写和 CPU-GPU 同步
  \item \textbf{结构化输出}：Guided Decoding / JSON Mode（Outlines、XGrammar）
\end{itemize}

\textbf{两个核心机制展开}：

\textbf{FlashAttention} 的关键洞察是注意力慢的根因不在计算而在\emph{显存读写}：朴素实现要把 $n \times n$ 的注意力矩阵写入 HBM 再读回，显存访问量 $O(n^2)$。FlashAttention 把 $Q, K, V$ 切成驻留 SRAM 的块（tile），在片上完成“块内 softmax + 累加”，并用 online-softmax 技术避免保存完整注意力矩阵——复杂度不变（仍是 $O(n^2)$ 计算），但显存降到 $O(n)$、访存降至理论下界，实测 2--4$\times$ 加速。FlashAttention-2 改进了块内并行划分，FlashAttention-3 利用 Hopper 的 TMA/WGMMA 异步流水再提速。

\textbf{Continuous Batching}（连续批处理 / in-flight batching）改变了推理服务的基本调度单位：从“静态 batch 全部完成才能出结果”变为“每个 decode step 动态重组 batch”——先完成的请求立即出队、新请求即刻插入。配合 KV Cache 分页（PagedAttention，类比操作系统虚拟内存：物理块不要求连续、逻辑视图连续），显存碎片从 60\%+ 降至 4\% 以下，系统吞吐提升 2--5$\times$ 以上。这两项技术是 vLLM 成为开源推理事实标准的核心，2024 年后已成为所有主流引擎的标配。

\textbf{引擎竞争格局（2026）}：vLLM 与 SGLang 双雄并立——前者生态最全、硬件覆盖最广，后者以 RadixAttention 前缀复用与结构化输出见长；TensorRT-LLM 是 NVIDIA 官方的极致性能路线（编译式优化，灵活性换速度）；llama.cpp 统治 CPU/端侧（GGUF 生态），MLX 服务 Apple Silicon。值得关注的是\textbf{模型厂商反向开源基础设施}的趋势：DeepSeek 2025 年初陆续开源 FlashMLA（MLA 专用注意力核）、DeepGEMM（FP8 GEMM 库）、3FS（训练数据存储）等组件，把“架构-系统协同设计”的成果回馈社区；国内则围绕华为 Ascend 等国产算力形成 vLLM-Ascend、MindIE 等适配层——推理引擎已经成为“模型-硬件”之间的必争中间层。

% ===================== 第七章 =====================
\section{推理模型（Reasoning Models）}

\subsection{从 CoT 到 System 2}

传统 LLM 是“System 1”——快速、直觉式输出。推理模型引入“System 2”——慢思考、逐步推理。

\textbf{思维链技术演进}：
\begin{itemize}
  \item \textbf{Zero-shot CoT}：“Let's think step by step”
  \item \textbf{Few-shot CoT}：提供推理示例
  \item \textbf{Tree of Thoughts (ToT)}：多路径探索 + 回溯
  \item \textbf{Self-Consistency}：采样多条 CoT，投票选答案
\end{itemize}

\textbf{范式转变：从“提示技巧”到“训练能力”}。2022 年的 CoT 提示属于 System 1 的“外骨骼”——能力上限仍受基座模型制约，且无法在简单任务上关闭。2024 年底 OpenAI o1 的出现标志着推理变成一种\emph{训练出来的能力}：模型在训练中习得“何时思考、思考多久、如何自我修正”，思考长度成为可配置的资源——这即 \textbf{Test-time Compute}（测试时计算）：把原本投入“训练”的算力部分转移到“推理时”，用更多 token 换更高准确率。OpenAI、DeepSeek 均报告了\emph{推理 token 数与任务难度/准确率的幂律关系}，即“扩展定律在时间维度的延伸”：
\[
\text{准确率} \;\propto\; f(\text{训练算力}) \cdot g(\text{推理算力}), \qquad g \text{ 同样近似幂律}
\]
这一“第三扩展轴”（参数 $N$、数据 $D$、推理算力 $R$）是 2025--2026 年大模型竞争的主战场（详见第 10 章）。

\begin{figure}[htbp]
\centering
% fig09_testtime_compute.tex — 图9 Test-time Compute 收益曲线
% 《深度学习与大语言模型：前沿技术全景报告》图示库
% 由 dl_llm_report.tex 抽取，经 % fig09_testtime_compute.tex — 图9 Test-time Compute 收益曲线
% 《深度学习与大语言模型：前沿技术全景报告》图示库
% 由 dl_llm_report.tex 抽取，经 % fig09_testtime_compute.tex — 图9 Test-time Compute 收益曲线
% 《深度学习与大语言模型：前沿技术全景报告》图示库
% 由 dl_llm_report.tex 抽取，经 \input{figs/fig09_testtime_compute} 引用（caption/label 在主文件）
% 注意：本文件为片段，依赖主文件导言区的 tikz/pgfplots 宏包、颜色定义与 tikzset 样式，不可单独编译。

\begin{tikzpicture}
\begin{semilogxaxis}[width=0.8\textwidth, height=7.2cm, clip=false,
  xlabel={思考 token 预算（对数刻度）}, ylabel={任务准确率（示意）},
  xtick=\empty, ytick=\empty, axis lines=left, xmin=1, xmax=100000, ymin=0, ymax=100]
\addplot[very thick, cblue, domain=1:100000, samples=200] {8 + 80*(1 - exp(-(ln(x))^1.5/28))};
\addplot[thick, cgray, dashed, domain=1:100000] {90};
\node[lab, anchor=south west, align=left] at (axis cs:2,92.8) {任务难度上界（可望不可即）};
\node[lab, anchor=south east, align=left] at (axis cs:95000,82) {饱和区：成本翻倍、收益骤减};
\node[lab, anchor=west, align=left] at (axis cs:60,10) {近似幂律区间：\\预算翻倍 $\to$ 准确率可观提升};
\end{semilogxaxis}
\end{tikzpicture}
 引用（caption/label 在主文件）
% 注意：本文件为片段，依赖主文件导言区的 tikz/pgfplots 宏包、颜色定义与 tikzset 样式，不可单独编译。

\begin{tikzpicture}
\begin{semilogxaxis}[width=0.8\textwidth, height=7.2cm, clip=false,
  xlabel={思考 token 预算（对数刻度）}, ylabel={任务准确率（示意）},
  xtick=\empty, ytick=\empty, axis lines=left, xmin=1, xmax=100000, ymin=0, ymax=100]
\addplot[very thick, cblue, domain=1:100000, samples=200] {8 + 80*(1 - exp(-(ln(x))^1.5/28))};
\addplot[thick, cgray, dashed, domain=1:100000] {90};
\node[lab, anchor=south west, align=left] at (axis cs:2,92.8) {任务难度上界（可望不可即）};
\node[lab, anchor=south east, align=left] at (axis cs:95000,82) {饱和区：成本翻倍、收益骤减};
\node[lab, anchor=west, align=left] at (axis cs:60,10) {近似幂律区间：\\预算翻倍 $\to$ 准确率可观提升};
\end{semilogxaxis}
\end{tikzpicture}
 引用（caption/label 在主文件）
% 注意：本文件为片段，依赖主文件导言区的 tikz/pgfplots 宏包、颜色定义与 tikzset 样式，不可单独编译。

\begin{tikzpicture}
\begin{semilogxaxis}[width=0.8\textwidth, height=7.2cm, clip=false,
  xlabel={思考 token 预算（对数刻度）}, ylabel={任务准确率（示意）},
  xtick=\empty, ytick=\empty, axis lines=left, xmin=1, xmax=100000, ymin=0, ymax=100]
\addplot[very thick, cblue, domain=1:100000, samples=200] {8 + 80*(1 - exp(-(ln(x))^1.5/28))};
\addplot[thick, cgray, dashed, domain=1:100000] {90};
\node[lab, anchor=south west, align=left] at (axis cs:2,92.8) {任务难度上界（可望不可即）};
\node[lab, anchor=south east, align=left] at (axis cs:95000,82) {饱和区：成本翻倍、收益骤减};
\node[lab, anchor=west, align=left] at (axis cs:60,10) {近似幂律区间：\\预算翻倍 $\to$ 准确率可观提升};
\end{semilogxaxis}
\end{tikzpicture}

\caption{Test-time Compute 的典型收益曲线（示意）：在对数思考预算下，准确率先呈近似幂律爬升，随后进入饱和区——这决定了“思考预算”产品必须按任务难度动态定价，而非无脑拉满。}
\label{fig:ttc}
\end{figure}

\textbf{2025 年的实证注脚}：o3 在发布时公布的高算力模式取得 \textbf{ARC-AGI 87.5\%}（此前两年所有模型长期在 30\% 以下徘徊），但单任务成本与时延大幅上升；o4-mini 在 AIME 2025 达到 \textbf{99.5\%} pass@1，并首次实现推理模型的\emph{全工具链调用}（搜索、Python、图像生成随思考过程穿插使用）。“Test-time compute 换准确率”由此从论文结论变成可售卖的产品开关（low/medium/high 思考档位），也把“每 token 成本”的定价模型改写为“每任务成本”。

\subsection{代表性推理模型}

\begin{table}[htbp]
\centering
\small
\begin{tabular}{@{}llp{5cm}@{}}
\toprule
\textbf{模型} & \textbf{核心方法} & \textbf{亮点}\\
\midrule
OpenAI o1 & 长 CoT + RL & 用户不可见推理过程\\
OpenAI o3 / o4-mini & 可配置思考预算 + 全工具链 & ARC-AGI 87.5\%；AIME 2025 99.5\%\\
DeepSeek-R1 & 纯 RL + GRPO + 规则奖励 & RL 成本 29.4 万美元，开源\\
Qwen3 & 混合思考 + 思考预算 & 单模型双模：/think 与 /no\_think\\
Gemini Thinking & 类 o1 推理 & 原生多模态推理\\
\bottomrule
\end{tabular}
\caption{代表性推理模型}
\end{table}

\subsubsection{DeepSeek-R1 的核心创新}

\begin{figure}[htbp]
\centering
% fig10_r1_pipeline.tex — 图10 DeepSeek-R1 四阶段训练流水线
% 《深度学习与大语言模型：前沿技术全景报告》图示库
% 由 dl_llm_report.tex 抽取，经 % fig10_r1_pipeline.tex — 图10 DeepSeek-R1 四阶段训练流水线
% 《深度学习与大语言模型：前沿技术全景报告》图示库
% 由 dl_llm_report.tex 抽取，经 % fig10_r1_pipeline.tex — 图10 DeepSeek-R1 四阶段训练流水线
% 《深度学习与大语言模型：前沿技术全景报告》图示库
% 由 dl_llm_report.tex 抽取，经 \input{figs/fig10_r1_pipeline} 引用（caption/label 在主文件）
% 注意：本文件为片段，依赖主文件导言区的 tikz/pgfplots 宏包、颜色定义与 tikzset 样式，不可单独编译。

\begin{tikzpicture}[node distance=4mm,
  rbox/.style={text width=56mm, align=center}]
\node[gry, rbox] (base) {DeepSeek-V3-Base 基座};
\node[blk, rbox, below=of base] (cs) {阶段一：冷启动 SFT（数千条长 CoT）};
\node[grn, rbox, below=of cs] (rl1) {阶段二：推理 RL（GRPO + 规则奖励）};
\node[blk, rbox, below=of rl1] (rs) {阶段三：拒绝采样 + SFT（80 万条样本）};
\node[grn, rbox, below=of rs] (rl2) {阶段四：全场景 RL（有用性 + 无害性）};
\node[orn, rbox, below=of rl2] (r1) {DeepSeek-R1};
\draw[arr] (base)--(cs); \draw[arr] (cs)--(rl1); \draw[arr] (rl1)--(rs); \draw[arr] (rs)--(rl2); \draw[arr] (rl2)--(r1);
\draw[rrr] (base.west) -- ++(-0.8cm,0) |- node[lab, pos=0.2, left, align=center] {R1-Zero：跳过 SFT\\ 直接 RL（涌现推理\\ 但语言混杂）} (rl1.west);
\node[gry, rbox, below=4mm of r1, font=\small] (sm) {Qwen / Llama 1.5B--70B 共 6 个蒸馏模型（全部开源）};
\draw[brr] (r1.east) -- ++(0.9cm,0) |- node[lab, pos=0.35, right, align=left] {蒸馏：60 万推理样本\\ + 20 万非推理样本\\ SFT 至小模型} (sm.east);
\end{tikzpicture}
 引用（caption/label 在主文件）
% 注意：本文件为片段，依赖主文件导言区的 tikz/pgfplots 宏包、颜色定义与 tikzset 样式，不可单独编译。

\begin{tikzpicture}[node distance=4mm,
  rbox/.style={text width=56mm, align=center}]
\node[gry, rbox] (base) {DeepSeek-V3-Base 基座};
\node[blk, rbox, below=of base] (cs) {阶段一：冷启动 SFT（数千条长 CoT）};
\node[grn, rbox, below=of cs] (rl1) {阶段二：推理 RL（GRPO + 规则奖励）};
\node[blk, rbox, below=of rl1] (rs) {阶段三：拒绝采样 + SFT（80 万条样本）};
\node[grn, rbox, below=of rs] (rl2) {阶段四：全场景 RL（有用性 + 无害性）};
\node[orn, rbox, below=of rl2] (r1) {DeepSeek-R1};
\draw[arr] (base)--(cs); \draw[arr] (cs)--(rl1); \draw[arr] (rl1)--(rs); \draw[arr] (rs)--(rl2); \draw[arr] (rl2)--(r1);
\draw[rrr] (base.west) -- ++(-0.8cm,0) |- node[lab, pos=0.2, left, align=center] {R1-Zero：跳过 SFT\\ 直接 RL（涌现推理\\ 但语言混杂）} (rl1.west);
\node[gry, rbox, below=4mm of r1, font=\small] (sm) {Qwen / Llama 1.5B--70B 共 6 个蒸馏模型（全部开源）};
\draw[brr] (r1.east) -- ++(0.9cm,0) |- node[lab, pos=0.35, right, align=left] {蒸馏：60 万推理样本\\ + 20 万非推理样本\\ SFT 至小模型} (sm.east);
\end{tikzpicture}
 引用（caption/label 在主文件）
% 注意：本文件为片段，依赖主文件导言区的 tikz/pgfplots 宏包、颜色定义与 tikzset 样式，不可单独编译。

\begin{tikzpicture}[node distance=4mm,
  rbox/.style={text width=56mm, align=center}]
\node[gry, rbox] (base) {DeepSeek-V3-Base 基座};
\node[blk, rbox, below=of base] (cs) {阶段一：冷启动 SFT（数千条长 CoT）};
\node[grn, rbox, below=of cs] (rl1) {阶段二：推理 RL（GRPO + 规则奖励）};
\node[blk, rbox, below=of rl1] (rs) {阶段三：拒绝采样 + SFT（80 万条样本）};
\node[grn, rbox, below=of rs] (rl2) {阶段四：全场景 RL（有用性 + 无害性）};
\node[orn, rbox, below=of rl2] (r1) {DeepSeek-R1};
\draw[arr] (base)--(cs); \draw[arr] (cs)--(rl1); \draw[arr] (rl1)--(rs); \draw[arr] (rs)--(rl2); \draw[arr] (rl2)--(r1);
\draw[rrr] (base.west) -- ++(-0.8cm,0) |- node[lab, pos=0.2, left, align=center] {R1-Zero：跳过 SFT\\ 直接 RL（涌现推理\\ 但语言混杂）} (rl1.west);
\node[gry, rbox, below=4mm of r1, font=\small] (sm) {Qwen / Llama 1.5B--70B 共 6 个蒸馏模型（全部开源）};
\draw[brr] (r1.east) -- ++(0.9cm,0) |- node[lab, pos=0.35, right, align=left] {蒸馏：60 万推理样本\\ + 20 万非推理样本\\ SFT 至小模型} (sm.east);
\end{tikzpicture}

\caption{DeepSeek-R1 四阶段训练流水线：橙色虚线是 R1-Zero 的对照路径（验证“纯 RL 涌现推理”）；蓝色实线展示蒸馏路线——小模型无需自己跑 RL，直接继承 R1 的推理模式。}
\label{fig:r1}
\end{figure}

\textbf{GRPO（Group Relative Policy Optimization）}是 R1 的关键算法：对同一 prompt 采样一组 $G$ 个回答 $\{y_1, \ldots, y_G\}$，用组内统计量直接计算相对优势
\[
\hat{A}_i = \frac{r_i - \operatorname{mean}(\{r_j\}_{j=1}^G)}{\operatorname{std}(\{r_j\}_{j=1}^G)}
\]
再套用 PPO 的 clip 目标 + KL 约束优化策略。与 PPO 的本质区别是\emph{不需要 Critic 网络}——优势估计由组内对比免费给出，显存占用与工程复杂度大幅下降（PPO 需要同时维护 Policy / Reference / Reward / Critic 四个模型）。在数学、代码等可自动判分任务上，规则奖励恰好提供了 $r_i$，使整套流水线摆脱了对人类标注的依赖。

\textbf{四阶段训练流水线}（R1 技术报告）：
\begin{enumerate}
  \item \textbf{冷启动 SFT}：用数千条高质量长 CoT 样本教会模型“长思考”的基本行为模式，抑制 R1-Zero 的语言混杂与可读性问题；
  \item \textbf{大规模 RL（GRPO）}：用规则奖励在数学/代码/STEM 上强化；
  \item \textbf{拒绝采样 + SFT}：用 R1 生成并经规则验证的约 80 万条样本（60 万推理 + 20 万非推理）反哺基座重新 SFT——R1 主模型即此产物；
  \item \textbf{全场景 RL}：在推理数据之外加入有用性/无害性奖励，覆盖通用对话场景。
\end{enumerate}
R1-Zero（跳过冷启动、直接 RL）实验证明推理能力可从基座“涌现”，但出现语言混杂与可读性退化——说明少量冷启动数据仍是塑造“思考风格”的有效手段。

\textbf{这套流水线的三个要点}：其一，\textbf{规则奖励替代神经奖励模型}——数学用答案校验、代码用测试用例，从源头规避 reward hacking；其二，训练过程出现教科书级的\textbf{涌现行为}——模型自发学会“等等，让我重新检查”“我之前的方法可能有误”式的自我验证、回溯与替代方案探索（R1-Zero 训练曲线上出现标志性的 \emph{“aha moment”}：反思类 token 随训练显著增多，基准准确率同步跳变）；其三，\textbf{能力可蒸馏}——把 R1 生成的推理样本直接 SFT 到 Qwen/Llama（1.5B--70B 共 6 个小模型），小模型在数学基准上甚至优于同规模直接 RL，说明推理模式可以“教会”而不必“重新发现”。

\textbf{成本}：官方 2025 年 9 月披露，R1 的 RL 增量训练成本仅约 \textbf{29.4 万美元}（V3 基座预训练约 \$5.6M）；作为对照，o1 的训练成本从未公开，外界普遍估计在数千万美元量级。R1 以开源权重 + 极低训练成本证明：\emph{推理能力不再是大算力公司的专属品}——这直接触发了 2025 年初的“DeepSeek 冲击波”（见第 11 章）。

\textbf{蒸馏与 RL 的边界}是 R1 论文留下的重要工程结论：把 R1 的推理样本蒸馏给 1.5B--7B 小模型，效果\emph{优于}小模型自己跑 RL——说明推理“模式”可以低成本迁移；但小模型蒸馏版的天花板被老师锁死，要想“超越老师”，仍然需要更大的基座 + 完整的 RL 过程。这为业界画了一张实用路线图：\emph{应用型小模型走蒸馏（便宜、快），前沿能力探索走 RL（昂贵、但能产生新知识）}。此外，思维链的\emph{可见性}也成了产品决策：DeepSeek 公开完整思考链（利于研究、审计与信任），OpenAI 只展示摘要（防蒸馏、降低安全暴露面），Qwen3 干脆把“是否思考”做成用户可切的开关——同一技术，三种产品哲学。

\subsection{过程奖励模型（PRM）}

\begin{table}[htbp]
\centering
\small
\begin{tabular}{@{}ll@{}}
\toprule
\textbf{类型} & \textbf{特点}\\
\midrule
ORM（结果奖励） & 只看最终答案对不对\\
\textbf{PRM（过程奖励）} & 对每一步推理打分\\
\bottomrule
\end{tabular}
\caption{奖励模型类型}
\end{table}

\textbf{ORM vs PRM 的本质差异}：ORM 只对“最终答案”打分，模型可以“蒙对”（答案正确但推理错误），奖励信号是稀疏的；PRM 对每一步中间状态打分，奖励稠密、能定位错误发生的位置，代价是需要大量\emph{过程标注}（如 PRM800K 数据集约 80 万条逐步标注）。推理时 PRM 通常与搜索算法联用：Beam Search / MCTS 展开多条推理路径，用 PRM 分数剪枝——这本质上把“单次生成”变成了“带评估的搜索”，是 Test-time Compute 的另一实现形态。2025 年后，随着 LLM-as-a-Judge 质量提升，\emph{自动化过程评估}（让强模型为弱模型的推理逐步打分）成为降低标注成本的主流手段。

\textbf{来自 R1 的反面证据}：DeepSeek-R1 论文把 PRM 与 MCTS 均列为\emph{失败的尝试}——显式 PRM 难以可靠判定中间步骤的正确性、天然易被 reward hacking，且每轮搜索迭代都需重训奖励模型；MCTS 则受制于指数级搜索空间与“如何把生成切成可搜索的步”这一未解难题。这解释了 2025 年主流路线为何收敛于“纯 RL + 规则奖励 + 自由长思考”：与其外挂搜索，不如让策略网络把搜索行为内化进参数。

\textbf{2025 趋势}：
\begin{itemize}
  \item \textbf{Automated PRM}：用 LLM 自动生成分步评估
  \item \textbf{MCTS + PRM}：AlphaGo 思路用于数学推理
  \item \textbf{Test-time Compute}：增加推理时间换准确率
\end{itemize}

\subsection{推理训练技术栈}

\[
\begin{aligned}
&\text{数据构造（CoT·自我博弈·拒绝采样）} \to \text{SFT} \to \text{RL（PPO/GRPO）}\\
&\to \text{测试时策略（ToT/MCTS）}
\end{aligned}
\]

\subsection{推理的边界与局限}

\begin{itemize}
  \item \textbf{幻觉}：知识性 vs 指令性；缓解：RAG、Self-Consistency
  \item \textbf{长 CoT 陷阱}：过长的思考链引入噪声；“思维链不忠实”现象
  \item \textbf{组合推理极限}：多步推理错误率累积
  \item \textbf{2025 前沿}：推理+工具（搜索/计算器/代码执行）、Agentic 推理、世界模型
\end{itemize}

\begin{warnbox}{必须清醒认识的三点局限}
\begin{enumerate}
  \item \textbf{CoT 忠实性问题}：Turpin et al. (2023) 等研究证实，模型输出的“思考过程”不一定反映其真实的决策依据——模型可能“先知道答案，再编造理由”。因此\emph{不能把 CoT 当作可审计的解释}，高可靠场景需要外部验证器（执行代码、对照知识库）。
  \item \textbf{奖励 hacking 的推理版}：当奖励来自 LLM 评审或宽松判分时，模型会学会“表演推理”（冗长但空洞的思考链）而非真正推理——奖励信号的可验证性（代码单测、数学答案比对）是推理 RL 质量的上限。
  \item \textbf{推理的边际收益递减}：Test-time Compute 的幂律在特定难度后会饱和，且成本线性增长；工程上需要“思考预算”机制（o3 的可配置预算、Qwen3 的 /think 与 /no\_think 双模式）在成本与准确率之间做显式权衡。
\end{enumerate}
\end{warnbox}

% ===================== 第八章 =====================
\section{多模态大模型}

\subsection{架构范式}

\begin{figure}[htbp]
\centering
\resizebox{\textwidth}{!}{%
% fig11_multimodal_architecture.tex — 图11 多模态大模型标准架构
% 《深度学习与大语言模型：前沿技术全景报告》图示库
% 由 dl_llm_report.tex 抽取，经 % fig11_multimodal_architecture.tex — 图11 多模态大模型标准架构
% 《深度学习与大语言模型：前沿技术全景报告》图示库
% 由 dl_llm_report.tex 抽取，经 % fig11_multimodal_architecture.tex — 图11 多模态大模型标准架构
% 《深度学习与大语言模型：前沿技术全景报告》图示库
% 由 dl_llm_report.tex 抽取，经 \input{figs/fig11_multimodal_architecture} 引用（caption/label 在主文件）
% 注意：本文件为片段，依赖主文件导言区的 tikz/pgfplots 宏包、颜色定义与 tikzset 样式，不可单独编译。

\begin{tikzpicture}
  % ---- 输入模态（左） ----
  \node[gry, minimum width=22mm] (img) at (0,1.7) {图像 / 文档};
  \node[gry, minimum width=22mm] (vid) at (0,0) {视频（帧序列）};
  \node[gry, minimum width=22mm] (aud) at (0,-1.7) {音频};
  % ---- 专用编码器 ----
  \node[blk, minimum width=26mm] (venc) at (3.5,1.1) {视觉编码器\\ ViT / ConvNeXt};
  \node[blk, minimum width=26mm] (aenc) at (3.5,-1.1) {音频编码器};
  % ---- 投影层 / LLM / 输出：阶段间隙 ≥0.7cm，避免贴挤 ----
  \node[orn, minimum width=26mm] (proj) at (6.9,0) {投影层\\ 对齐到词表空间};
  \node[grn, minimum width=28mm, minimum height=13mm] (llm) at (10.3,0) {LLM Backbone\\（文本 token 直连）};
  \node[orn, minimum width=28mm] (out) at (13.9,0) {输出\\ 文本 · 动作 · 工具调用};
  \node[gry, minimum width=22mm] (txt) at (6.9,2.9) {文本 token};
  % ---- 连线：编码器入口/投影层入口均分散锚点，箭头不再汇聚打结 ----
  \draw[arr] (img.east) -- ([yshift=2mm]venc.west);
  \draw[arr] (vid.east) -- ([yshift=-2mm]venc.west);
  \draw[arr] (aud.east) -- (aenc.west);
  \draw[arr] (venc.east) -- ([yshift=2mm]proj.west);
  \draw[arr] (aenc.east) -- ([yshift=-2mm]proj.west);
  \draw[arr] (txt) -- (proj);
  \draw[arr] (txt.east) to[bend left=22] (llm.70);
  \draw[arr] (proj)--(llm); \draw[arr] (llm)--(out);
  \node[lab, anchor=east] at (5.0,2.9) {拼接式（LLaVA 系）};
  \node[lab, anchor=west] at (10.3,2.1) {原生式（Gemini 系）};
  \node[lab] at (7.1,-3.1) {视觉 token 成本：一张 1024$^2$ 图 $\approx$ 数百至上千 token，是多模态请求成本的大头};
\end{tikzpicture}
 引用（caption/label 在主文件）
% 注意：本文件为片段，依赖主文件导言区的 tikz/pgfplots 宏包、颜色定义与 tikzset 样式，不可单独编译。

\begin{tikzpicture}
  % ---- 输入模态（左） ----
  \node[gry, minimum width=22mm] (img) at (0,1.7) {图像 / 文档};
  \node[gry, minimum width=22mm] (vid) at (0,0) {视频（帧序列）};
  \node[gry, minimum width=22mm] (aud) at (0,-1.7) {音频};
  % ---- 专用编码器 ----
  \node[blk, minimum width=26mm] (venc) at (3.5,1.1) {视觉编码器\\ ViT / ConvNeXt};
  \node[blk, minimum width=26mm] (aenc) at (3.5,-1.1) {音频编码器};
  % ---- 投影层 / LLM / 输出：阶段间隙 ≥0.7cm，避免贴挤 ----
  \node[orn, minimum width=26mm] (proj) at (6.9,0) {投影层\\ 对齐到词表空间};
  \node[grn, minimum width=28mm, minimum height=13mm] (llm) at (10.3,0) {LLM Backbone\\（文本 token 直连）};
  \node[orn, minimum width=28mm] (out) at (13.9,0) {输出\\ 文本 · 动作 · 工具调用};
  \node[gry, minimum width=22mm] (txt) at (6.9,2.9) {文本 token};
  % ---- 连线：编码器入口/投影层入口均分散锚点，箭头不再汇聚打结 ----
  \draw[arr] (img.east) -- ([yshift=2mm]venc.west);
  \draw[arr] (vid.east) -- ([yshift=-2mm]venc.west);
  \draw[arr] (aud.east) -- (aenc.west);
  \draw[arr] (venc.east) -- ([yshift=2mm]proj.west);
  \draw[arr] (aenc.east) -- ([yshift=-2mm]proj.west);
  \draw[arr] (txt) -- (proj);
  \draw[arr] (txt.east) to[bend left=22] (llm.70);
  \draw[arr] (proj)--(llm); \draw[arr] (llm)--(out);
  \node[lab, anchor=east] at (5.0,2.9) {拼接式（LLaVA 系）};
  \node[lab, anchor=west] at (10.3,2.1) {原生式（Gemini 系）};
  \node[lab] at (7.1,-3.1) {视觉 token 成本：一张 1024$^2$ 图 $\approx$ 数百至上千 token，是多模态请求成本的大头};
\end{tikzpicture}
 引用（caption/label 在主文件）
% 注意：本文件为片段，依赖主文件导言区的 tikz/pgfplots 宏包、颜色定义与 tikzset 样式，不可单独编译。

\begin{tikzpicture}
  % ---- 输入模态（左） ----
  \node[gry, minimum width=22mm] (img) at (0,1.7) {图像 / 文档};
  \node[gry, minimum width=22mm] (vid) at (0,0) {视频（帧序列）};
  \node[gry, minimum width=22mm] (aud) at (0,-1.7) {音频};
  % ---- 专用编码器 ----
  \node[blk, minimum width=26mm] (venc) at (3.5,1.1) {视觉编码器\\ ViT / ConvNeXt};
  \node[blk, minimum width=26mm] (aenc) at (3.5,-1.1) {音频编码器};
  % ---- 投影层 / LLM / 输出：阶段间隙 ≥0.7cm，避免贴挤 ----
  \node[orn, minimum width=26mm] (proj) at (6.9,0) {投影层\\ 对齐到词表空间};
  \node[grn, minimum width=28mm, minimum height=13mm] (llm) at (10.3,0) {LLM Backbone\\（文本 token 直连）};
  \node[orn, minimum width=28mm] (out) at (13.9,0) {输出\\ 文本 · 动作 · 工具调用};
  \node[gry, minimum width=22mm] (txt) at (6.9,2.9) {文本 token};
  % ---- 连线：编码器入口/投影层入口均分散锚点，箭头不再汇聚打结 ----
  \draw[arr] (img.east) -- ([yshift=2mm]venc.west);
  \draw[arr] (vid.east) -- ([yshift=-2mm]venc.west);
  \draw[arr] (aud.east) -- (aenc.west);
  \draw[arr] (venc.east) -- ([yshift=2mm]proj.west);
  \draw[arr] (aenc.east) -- ([yshift=-2mm]proj.west);
  \draw[arr] (txt) -- (proj);
  \draw[arr] (txt.east) to[bend left=22] (llm.70);
  \draw[arr] (proj)--(llm); \draw[arr] (llm)--(out);
  \node[lab, anchor=east] at (5.0,2.9) {拼接式（LLaVA 系）};
  \node[lab, anchor=west] at (10.3,2.1) {原生式（Gemini 系）};
  \node[lab] at (7.1,-3.1) {视觉 token 成本：一张 1024$^2$ 图 $\approx$ 数百至上千 token，是多模态请求成本的大头};
\end{tikzpicture}

}%
\caption{多模态大模型的标准架构：各模态先经专用编码器压缩为 token 序列，由投影层对齐进 LLM 词表空间后与文本统一处理——融合发生的位置（投影层 vs 第一层）区分了拼接式与原生式两大路线。}
\label{fig:mm}
\end{figure}

\textbf{视觉 token 的成本账}值得单独一算：一张 $1024\times1024$ 图片经 patch 切分与动态分辨率适配后，通常产出数百至上千个视觉 token，相当于数千字文本——多模态请求的实际推理成本常是纯文本的 5--10 倍，且大头在预填充。因此第 6 章的推理优化在多模态场景从“可选项”变成“必需品”：视觉 token 压缩（平均池化/查询压缩）、图像前缀缓存、按需分辨率调度，本质都是在“看得清”与“算得起”之间找平衡。

\textbf{三种融合范式}值得区分：
\begin{enumerate}
  \item \textbf{拼接式（Late Fusion + Adapter）}：模态编码器独立预训练，经投影层映射到 LLM 词表空间后拼接（LLaVA、Qwen-VL、InternVL 均为这一路线）——工程上最成熟，视觉编码器可冻结，训练成本低；
  \item \textbf{早期融合（Early Fusion）}：不同模态从第一层起就在同一 Transformer 内交互（Gemini 原生多模态路线）——跨模态推理上限更高，但训练成本高、需要海量配对数据；
  \item \textbf{统一自回归 / 扩散混合}：生成侧用扩散模型（图像/视频）、理解侧用 LLM（Sora、Veo、可灵的“理解 + 生成”双管线）。
\end{enumerate}
对齐（Projection）是多模态模型的关键瓶颈：视觉 token 与文本 token 的\emph{语义粒度天然不对齐}（一个 14$\times$14 patch 对应几十个语义维度），投影层设计（线性 vs MLP vs Q-Former/Perceiver Resampler）直接决定细粒度能力（OCR、图表、空间关系）的上限。

\subsection{视觉语言模型（VLM）}

\begin{table}[htbp]
\centering
\small
\begin{tabular}{@{}llp{4.5cm}@{}}
\toprule
\textbf{模型} & \textbf{架构} & \textbf{特点}\\
\midrule
LLaVA 系列 & CLIP ViT + 线性投影 + LLaMA & 开源标杆\\
InternVL 系列 & ViT + MLP 投影 + LLM & 开源多模态第一梯队\\
Qwen2.5-VL & 原生动态分辨率 + 视频 & OCR/文档理解标杆\\
GPT-4o & 原生多模态（非拼接） & 实时语音交互\\
Gemini 2.5/3 & 原生多模态 + Deep Think & 文本+图+音+视频统一\\
\bottomrule
\end{tabular}
\caption{代表性 VLM}
\end{table}

\textbf{VLM 的标准训练配方}（以 LLaVA 系、Qwen-VL 系为代表）分两阶段：
\begin{enumerate}
  \item \textbf{对齐阶段（Alignment）}：冻结视觉编码器与 LLM，只训练投影层；数据用“图文对 + 图像描述”，目标是把视觉 token 拉进语言空间——这一阶段决定跨模态“对得准不准”；
  \item \textbf{指令微调阶段（SFT）}：解冻 LLM（或 LoRA），用多任务数据（VQA、OCR、图表理解、文档问答、视频时序）联合训练——决定“用得对不对”。
\end{enumerate}
2024--2025 年的两个关键改进：(1)\textbf{动态分辨率}——放弃固定 $336\times336$ 裁剪，按图像长宽比自适应切 patch（Qwen2-VL 的 Naive Dynamic Resolution），直接带来 OCR 与文档理解能力的代际提升；(2)\textbf{视觉 CoT}——在推理模型时代（第 7 章），VLM 同样引入长思考（先描述再推理、先规划搜索区域再作答），Gemini 与 Qwen 的 thinking 视觉模型即为代表。

\subsection{音频/语音}

\begin{itemize}
  \item \textbf{ASR}：Whisper（约 100 种语言、68 万小时训练数据）、SenseVoice
  \item \textbf{TTS}：VALL-E、CosyVoice、F5-TTS、NaturalSpeech 3
  \item \textbf{实时语音对话}：GPT-4o Live、Gemini Live（2025 主流）
  \item \textbf{2025 趋势}：全双工语音、情感/语调可控 TTS、低延迟流式推理
\end{itemize}

\textbf{技术要点}：现代 TTS 的主流路线是\emph{声码 token 化}——用 EnCodec 类神经编解码器把语音压成离散 token 序列，TTS 就退化为“文本 token $\to$ 语音 token”的自回归/流匹配生成问题，从而复用 LLM 工具链。Whisper 则证明 68 万小时多语言数据 + 强数据质量足以训出通用 ASR。实时语音对话（GPT-4o Live、Gemini Live、国内豆包/通义电话级产品）的端到端延迟预算通常要求 $<$500ms，工程上依赖：流式 VAD（端点检测）+ 流式 ASR + 小模型快答 + 流式 TTS 首包优化；更激进的方向是“语音-文本同一模型”（end-to-end speech LLM），省掉 ASR/TTS 级联的信息损失，但成本与可控性仍是瓶颈。

\subsection{视频理解与生成}

\begin{itemize}
  \item \textbf{理解}：Qwen2-VL（视频）、Gemini（原生视频）
  \item \textbf{生成}（2025 竞争格局）：
  \begin{itemize}
    \item Sora 2（OpenAI，2025-09）：物理一致性大幅提升，同步推出独立 App 走短视频社区路线
    \item Veo 3 / 3.1（Google，2025-05/10）：首个\emph{音画同步}（原生生成对白与音效）的视频模型
    \item Kling/可灵 2.x（快手）：国产代表，首帧/多参考图控制成熟
    \item Seedance 1.0（字节，2025-06）：极致推理性价比；Runway Gen-4：商业工作流
    \item 开源阵营：通义万相 Wan2.2、HunyuanVideo、LTX-Video
  \end{itemize}
  \item \textbf{技术路线}：Diffusion Transformer (DiT)、3D 卷积 + 时间注意力、潜空间视频扩散
\end{itemize}

\textbf{视频生成的架构共识}（Sora、Veo、可灵、Runway 殊途同归）：(1)\textbf{3D 时空 VAE}——把视频压进潜空间（时间维压缩 4$\times$ 左右、空间维压缩 8--16$\times$），扩散/流匹配在潜空间进行，把计算量降低一到两个数量级；(2)\textbf{DiT}——用 Transformer 替代 U-Net 作为生成骨干，天然支持任意分辨率与时长的 scaling；(3)\textbf{条件注入}——文本经交叉注意力或联合序列注入。2025--2026 的竞争焦点已从“能不能生成”转向\emph{物理一致性}（遮挡、流体、碰撞）、\emph{长时一致性}（角色/场景不漂移）与\emph{音画同步}（Veo 3、可灵 2.x 的原生音频），以及成本——单条 5 秒视频 720p 的推理成本已降至数美分级别。此外，Google Genie 3（2025 年 8 月）代表另一条路线——\emph{可交互世界模型}：不生成固定视频，而是根据文本实时生成可自由探索的 3D 一致世界；它与 LeCun 的世界模型主张共同指向“视频生成之后的下一个范式”（见第 10、12 章）。

\subsection{多模态训练与推理}

\textbf{训练}：视觉-语言对齐与指令微调的两阶段配方（见上节）解决“看得懂”，而 2025 年起多模态后训练的重心转向“想得对、做得准”——(1)\textbf{多模态 RLHF/RL}：以人类偏好对齐图像质量与安全性，并用可验证奖励（定位框 IoU、OCR 精度）做 GRPO 类强化，Qwen-VL 系率先报告视觉定位能力的显著提升；(2)\textbf{视觉 CoT}：先描述图像、先框选感兴趣区域再推理——推理模型范式（第 7 章）自然延伸到视觉域，OpenAI o3 的“以图思考”（thinking with images）是代表，图像本身成为推理过程中可操作的中间变量；(3)\textbf{Agent 化}：屏幕截图理解 $\to$ 规划操作 $\to$ 执行反馈，GUI Agent 从 demo 走向产品——OpenAI Operator（2025 年 1 月，底层 CUA 模型）、Claude Computer Use、字节开源的 UI-TARS 与智谱 AutoGLM，以及 2025 年 3 月出圈的通用 Agent Manus。

\textbf{推理}：多模态请求（图文交错输入、视频帧采样、长上下文视觉 token）使 KV Cache 与预填充成本远高于纯文本，工程上依赖\textbf{视觉 token 压缩}（平均池化/查询压缩降低每图 token 数）、\textbf{前缀缓存}（共享图像的 KV 复用）与\textbf{动态分辨率}调度；vLLM/SGLang 均已原生支持图文交错输入与视觉前缀缓存，是 2025 年多模态服务化的标配。

% ===================== 第九章 =====================
\section{安全、对齐与治理}

\subsection{对齐技术栈}

对齐（Alignment）的目标是：让模型的\emph{输出分布}与人类意图、价值观保持一致，而非仅仅“预测下一个 token”。工业实践中没有银弹，而是\textbf{多层防御纵深}（defense in depth）——每一层都假设其他层可能失效：

\begin{enumerate}
  \item \textbf{Constitutional AI}（Anthropic）：用“宪法”原则指导自我批评 + RLAIF
  \item \textbf{System Prompt 工程}：安全指令 + 行为约束 + 角色设定
  \item \textbf{输出过滤}：分类器检测有害内容、关键词过滤、置信度阈值
  \item \textbf{红队测试（Red Teaming）}：自动化攻击（越狱、提示注入）+ 人工红队
  \item \textbf{评估与审计}：安全基准（SafetyBench、HarmBench、AdvBench）、透明度报告、第三方审计
\end{enumerate}

其中\emph{训练时对齐}（1）是根基，\emph{部署时防御}（2--3）是护栏，\emph{运营时对抗}（4--5）是持续更新机制。值得强调的是：任何单点防御都可被绕过（2024--2025 年的多项研究已系统性攻破单一 jailbreak 过滤器），只有纵深 + 持续红队才能维持安全水位——这与网络安全领域“没有绝对安全”的共识一致。

\textbf{2025--2026 年的四个新焦点}：(1)\textbf{Agent 与提示注入}——模型能“执行动作”之后，间接提示注入从“内容风险”升级为“执行风险”（被投毒的网页/文档/工具描述可诱导 Agent 泄露数据或执行恶意操作），OWASP LLM Top 10 与多国安全机构均将其列为一号威胁，防御转向权限最小化、人机确认（human-in-the-loop）、沙箱执行与动作审计；(2)\textbf{CoT 忠实度}——推理模型的可见思考链未必反映真实决策依据，把“监控思考链”本身作为安全机制存在系统性风险（Anthropic 2025 年委托的独立研究对此发出明确警告）；(3)\textbf{模型混淆与来源}——AI 生成内容的水印/显隐式标识（中国《标识办法》、C2PA 内容凭证、Google SynthID 类水印）从技术选项变成合规义务；(4)\textbf{攻防自动化}——LLM 驱动的自动化越狱与自动化红队同时成熟，安全水位从“发布时达标”转向“持续运营指标”。

\begin{figure}[htbp]
\centering
\resizebox{\textwidth}{!}{%
% fig12_defense_in_depth.tex — 图12 多层防御纵深
% 《深度学习与大语言模型：前沿技术全景报告》图示库
% 由 dl_llm_report.tex 抽取，经 % fig12_defense_in_depth.tex — 图12 多层防御纵深
% 《深度学习与大语言模型：前沿技术全景报告》图示库
% 由 dl_llm_report.tex 抽取，经 % fig12_defense_in_depth.tex — 图12 多层防御纵深
% 《深度学习与大语言模型：前沿技术全景报告》图示库
% 由 dl_llm_report.tex 抽取，经 \input{figs/fig12_defense_in_depth} 引用（caption/label 在主文件）
% 注意：本文件为片段，依赖主文件导言区的 tikz/pgfplots 宏包、颜色定义与 tikzset 样式，不可单独编译。

\begin{tikzpicture}[node distance=5.5mm]
\node[gry, minimum width=23mm] (u) {用户 / Agent};
\node[blk, right=5mm of u, minimum width=32mm, align=center] (in) {输入侧防御\\ 提示注入检测 · 权限最小化};
\node[grn, right=5mm of in, minimum width=36mm, align=center] (m) {模型\\（对齐已内化：SFT·RLHF·宪法）};
\node[blk, right=5mm of m, minimum width=32mm, align=center] (out2) {输出侧防御\\ 分类器 · 标识 · 置信度};
\node[gry, right=5mm of out2, minimum width=17mm] (r) {响应 / 执行};
\draw[arr] (u)--(in); \draw[arr] (in)--(m); \draw[arr] (m)--(out2); \draw[arr] (out2)--(r);
\node[orn, below=11mm of m, minimum width=96mm, align=center] (rt) {持续红队：自动化攻击库（越狱 / 注入模板）+ 人工专家对抗};
\draw[rrr] (rt.north) -- (in.south east);
\draw[rrr] (rt.north) -- (m.south);
\draw[rrr] (rt.north) -- (out2.south west);
\node[lab, below=1.5mm of rt] {发现的漏洞回流为训练数据与过滤规则——防御是循环，不是关卡};
\end{tikzpicture}
 引用（caption/label 在主文件）
% 注意：本文件为片段，依赖主文件导言区的 tikz/pgfplots 宏包、颜色定义与 tikzset 样式，不可单独编译。

\begin{tikzpicture}[node distance=5.5mm]
\node[gry, minimum width=23mm] (u) {用户 / Agent};
\node[blk, right=5mm of u, minimum width=32mm, align=center] (in) {输入侧防御\\ 提示注入检测 · 权限最小化};
\node[grn, right=5mm of in, minimum width=36mm, align=center] (m) {模型\\（对齐已内化：SFT·RLHF·宪法）};
\node[blk, right=5mm of m, minimum width=32mm, align=center] (out2) {输出侧防御\\ 分类器 · 标识 · 置信度};
\node[gry, right=5mm of out2, minimum width=17mm] (r) {响应 / 执行};
\draw[arr] (u)--(in); \draw[arr] (in)--(m); \draw[arr] (m)--(out2); \draw[arr] (out2)--(r);
\node[orn, below=11mm of m, minimum width=96mm, align=center] (rt) {持续红队：自动化攻击库（越狱 / 注入模板）+ 人工专家对抗};
\draw[rrr] (rt.north) -- (in.south east);
\draw[rrr] (rt.north) -- (m.south);
\draw[rrr] (rt.north) -- (out2.south west);
\node[lab, below=1.5mm of rt] {发现的漏洞回流为训练数据与过滤规则——防御是循环，不是关卡};
\end{tikzpicture}
 引用（caption/label 在主文件）
% 注意：本文件为片段，依赖主文件导言区的 tikz/pgfplots 宏包、颜色定义与 tikzset 样式，不可单独编译。

\begin{tikzpicture}[node distance=5.5mm]
\node[gry, minimum width=23mm] (u) {用户 / Agent};
\node[blk, right=5mm of u, minimum width=32mm, align=center] (in) {输入侧防御\\ 提示注入检测 · 权限最小化};
\node[grn, right=5mm of in, minimum width=36mm, align=center] (m) {模型\\（对齐已内化：SFT·RLHF·宪法）};
\node[blk, right=5mm of m, minimum width=32mm, align=center] (out2) {输出侧防御\\ 分类器 · 标识 · 置信度};
\node[gry, right=5mm of out2, minimum width=17mm] (r) {响应 / 执行};
\draw[arr] (u)--(in); \draw[arr] (in)--(m); \draw[arr] (m)--(out2); \draw[arr] (out2)--(r);
\node[orn, below=11mm of m, minimum width=96mm, align=center] (rt) {持续红队：自动化攻击库（越狱 / 注入模板）+ 人工专家对抗};
\draw[rrr] (rt.north) -- (in.south east);
\draw[rrr] (rt.north) -- (m.south);
\draw[rrr] (rt.north) -- (out2.south west);
\node[lab, below=1.5mm of rt] {发现的漏洞回流为训练数据与过滤规则——防御是循环，不是关卡};
\end{tikzpicture}

}%
\caption{多层防御纵深（defense in depth）：训练时对齐是根基，输入/输出防御是护栏，红队是持续校准机制——任何单层都可被绕过，安全水位取决于整条链与回流的闭环。}
\label{fig:defense}
\end{figure}

\subsection{关键安全挑战}

\begin{table}[htbp]
\centering
\small
\begin{tabular}{@{}ll@{}}
\toprule
\textbf{威胁} & \textbf{防御手段}\\
\midrule
越狱 (Jailbreak) & 输入过滤、多轮上下文监控、对抗训练\\
幻觉与虚假信息 & RAG、引用来源、置信度表达\\
偏见与公平性 & 去偏数据、公平性约束、多视角评估\\
隐私泄露 & 差分隐私、数据脱敏、本地化部署\\
提示注入 & 输入隔离、权限最小化、沙箱执行\\
AI 武器化 & 访问控制、速率限制、内容审计\\
\bottomrule
\end{tabular}
\caption{安全威胁与防御}
\end{table}

\textbf{2025--2026 新挑战}：
\begin{itemize}
  \item \textbf{AI 生成内容识别}：深度伪造与信任链危机——显式标识、隐式水印（SynthID 类）与内容凭证（C2PA）三条路线并行，中国《标识办法》已将显式 + 隐式标识定为强制义务；
  \item \textbf{模型窃取}：API 逆向蒸馏与模型提取攻击持续演化；蒸馏检测（输出分布指纹）技术尚不成熟，成为开源--闭源之争中的焦点议题；
  \item \textbf{供应链攻击}：恶意插件、MCP 工具描述投毒、模型仓库污染——权重与工具链本身成为新的攻击面；
  \item \textbf{行业自律升级}：2025 年 OpenAI、Anthropic、Google 等联合发布 AI 网络安全防御公开信（100+ 公司签署）；前沿实验室普遍建立“准备度框架”（Preparedness Framework），按能力风险等级设定部署门槛与红队强度。
\end{itemize}

\subsection{治理与监管}

AI 治理正在从“原则宣言”走向“可执行法规”。主要司法辖区的框架对比：

\begin{itemize}
  \item \textbf{EU AI Act}（2024 年 8 月生效，2025--2027 分阶段实施）：\emph{风险分级}监管——不可接受风险（如社会评分）禁止、高风险（招聘、医疗）严格合规义务（数据治理、日志、人工监督）、有限风险（透明度义务）、最小风险自由。全球首个综合性 AI 法，具有“布鲁塞尔效应”——2025 年 8 月起 GPAI（通用模型）义务生效，2026 年 8 月全面适用。
  \item \textbf{美国}：联邦层面以 NIST AI 风险管理框架（AI RMF 1.0/2.0）+ 行政令为主，2025 年政策转向“竞争力优先”，监管权下放至州（如加州 SB 1047 的立法博弈）；各机构（FDA、SEC、FTC）在各自领域并行监管。
  \item \textbf{中国}：《生成式人工智能服务管理暂行办法》（2023 年 8 月生效）+ 大模型备案制度（网信办“生成式 AI 服务”备案与深度合成算法备案双轨）+《人工智能安全治理框架》（1.0/2.0 版）；特色是\emph{发展与安全并重}的路径——备案制降低合规不确定性，同时推动行业安全基线（内容安全、训练数据合法性、个人信息保护）。2025 年起，面向公众的生成式服务需通过备案并履行安全评估义务；《人工智能法》已列入立法规划、处于草案研究阶段。
  \item \textbf{国际}：Bletchley AI 安全峰会（2023）$\to$ 首尔峰会（2024）$\to$ 2025 年 AI 安全治理进程延续；联合国 AI 治理对话、ISO/IEC 42001（AI 管理体系）等标准加速落地；OECD AI 原则成为多数国家立法的参照系；2025 年 8 月联合国大会决议设立\emph{独立国际科学小组}与\emph{治理全球对话}双机制，全球治理框架初具雏形。
\end{itemize}

\begin{keybox}{给从业者的合规要点（中国语境）}
面向公众提供生成式 AI 服务需完成：\textbf{算法备案}（深度合成/生成式）+ \textbf{安全评估} + \textbf{内容标识}（《人工智能生成合成内容标识办法》2025 年 9 月 1 日起施行，要求显式+隐式标识）；训练数据需合法来源、不得侵犯知识产权与个人信息；跨境数据传输遵守《数据安全法》《个人信息保护法》。这三条是 2025--2026 年国内 AI 产品上线的硬性门槛。
\end{keybox}

% ===================== 第十章 =====================
\section{可扩展性定律与未来方向}

\subsection{Scaling Laws 演进}

\subsubsection{Kaplan Scaling Laws (2020)}
\[
L \propto N^{-0.076} \cdot D^{-0.095} \cdot B^{-0.050}
\]
$N$=参数量，$D$=数据量，$B$=计算量。结论：等比扩展。

\textbf{幂律的含义}：$L \propto N^{-0.076}$ 意味着参数量每翻一倍，损失仅降低约 $1-2^{-0.076} \approx 5\%$——扩展是“有效但缓慢”的。幂律的工程价值在于\emph{可外推性}：用 1B--7B 小规模模型测得的曲线，可以相当可靠地预测 100B 级别的行为，这使“花小钱验证大钱决策”成为可能，也是后来所有 Scaling Law 研究的起点。Kaplan 报告的指数 $D$（0.095）$> N$（0.076），暗示“数据比参数更值钱”——Chinchilla 用严谨实验证实并放大了这一结论。

\subsubsection{Chinchilla Scaling Laws (2022)}
DeepMind 用更严格的控制变量（固定算力预算，扫描不同 $N/D$ 配比）修正了 Kaplan 的结论：\textbf{数据量应与参数量等比扩展，最优配比约为每 1B 参数训练 20B token}（即 $D \approx 20N$）。

例如：70B 模型 $\to$ $\sim$1.4T token。Chinchilla 70B 用 1.4T token 训练，效果优于 GPT-3 175B（300B token）——\emph{用 $\frac{1}{2.5}$ 的参数量、$\frac{1}{5}$ 的算力达到更好的效果}，核心贡献是指出了 GPT-3 被严重“欠训练于数据”。

但 Chinchilla 最优是针对\emph{训练成本最小化}的最优。若把\emph{推理成本}计入生命周期总账（模型要服务数年），结论会反转：训练更多 token、服务更小模型反而更划算。这正是 Llama 3 的选择：8B 模型训练 15T token（远超 $20N=160$B 的 Chinchilla 最优），405B 旗舰同样在约 15.6T token 上训练——业界称之为\emph{inference-aware scaling}。2025--2026 年的共识是：\textbf{Chinchilla 定律给出的是“训练侧”参考线，而产品经济学把最优推向“更少参数 + 更多数据 + 更多推理算力”的三角}。

\subsubsection{超越 Chinchilla}

\begin{itemize}
  \item \textbf{Llama 3}：8B/70B/405B 全系 15T+ token——为多年推理服务而“过量训练”；
  \item \textbf{Qwen3}：36T token、119 种语言（2025 年开源阵营最大规模预训练），显著调高多语言与 Agent 数据权重；
  \item \textbf{DeepSeek-V3}：671B MoE 仅激活 37B、14.8T token——MoE 让“每 FLOP 换多少有效知识”成为新的比较基准；
  \item \textbf{数据质量路线}：Phi 系列证明“教科书级合成数据 + 严格过滤”可用 14B 逼近更大稠密模型的推理能力；合成数据闭环（模型生成 $\to$ 验证器过滤 $\to$ 回炉再训练）成为绕开数据墙的主力工程；
  \item \textbf{超参迁移}（$\mu$P）：小模型调出的最优超参可按解析规则映射到任意规模，降低“烧大钱试错”的风险。
\end{itemize}
\textbf{综合判断}：Scaling Law 的“第四轴”——\emph{推理时算力}（第 7 章）已与参数、数据并列；2025--2026 年的共识公式是“\textbf{更少激活参数 + 更多高质量（含合成）数据 + 更长思考时间}”。

\begin{figure}[htbp]
\centering
\input{figs/fig13_scaling_laws}
\caption{Scaling Law 的示意曲线（对数-对数坐标下近似直线）：在同等算力预算下，“数据等比”路线（Chinchilla）优于“参数优先”路线；数据质量与合成数据可将整条曲线进一步下压。}
\label{fig:scaling}
\end{figure}

\textbf{两个方法论警告}：其一，“涌现能力”有相当部分是\emph{度量伪影}——不连续指标（答对/答错）在连续的能力增长上自然产生跳变（Schaeffer et al., 2023），换用连续指标后“涌现”趋于平滑；其二，\emph{基准污染}（benchmark contamination）——测试题经互联网泄漏进训练语料使分数虚高，2025 年后头部团队转向“私有题库 + 动态评测 + 人工复核”三件套。读榜单的正确姿势：同源对比、关注置信区间、警惕考古式刷榜；跨代比较时优先看 Agentic 与长周期任务（更难污染、更贴近真实价值）。

\subsection{数据墙与计算墙}

\begin{warnbox}{数据墙（Data Wall）}
高质量公开人类文本可能于 2026--2032 年间耗尽（Epoch AI 估计）。解决方向：合成数据、多模态数据、交互数据、数据飞轮（模型生成 $\to$ 人类筛选 $\to$ 再训练）。
\end{warnbox}

\begin{warnbox}{计算墙（Compute Wall）}
单次训练 100B+ 模型成本 \$30M--\$100M+，能源消耗巨大。效率方向：MoE 稀疏激活、FP8/INT8 训练、MLA/线性注意力架构效率、高效优化器。
\end{warnbox}

\textbf{算力的地缘维度}（中国语境）：2022 年起的高端 GPU 出口管制使 H100/B200 对华禁售、H800/H20 供应持续受限，中国训练栈事实上走出三条路线——其一，\emph{国产算力替代}（华为 Ascend 910B/C 系列等，CANN 软件生态与 CUDA 的差距逐年缩小，MindIE/vLLM-Ascend 等推理适配层快速补位）；其二，\emph{以效率换算力}（MoE、MLA、FP8、MTP 的激进组合，DeepSeek-V3 用 2048 块 H800 训出 671B 模型即是范本）；其三，\emph{算力调度优化}（东数西算、算力租赁市场）。理解这个背景才能理解中国阵营的技术偏好：\emph{大量“效率创新”是被约束逼出来的，而约束反而造就了可输出的方法论}——R1 的低成本 RL、Qwen3 的高性价比 MoE 都是例证。

\subsection{能力收敛假说（CCH）：能力由访问结构定价}

Scaling Laws 刻画的是\emph{训练侧}的能力增长；而“规模能否自动换来智能”这一问题，在 2024--2026 年出现了理论转折。柏拉图表征假说（Platonic Representation Hypothesis, PRH；Huh et al., 2024）主张：随规模与数据多样性增长，异构网络的内部表征收敛到共享的“现实统计模型”——在这一图景中，表征收敛是规模免费赠予的。然而 2026 年的实证工作揭示了 PRH 的“裂缝”：不同模型可以在表征空间中高度对齐，却在推理与决策行为上显著分歧（Usama 与 Chang, 2026）。

Chen et al.（2026，arXiv:2607.14144）据此提出\textbf{能力收敛假说}（Capability Convergence Hypothesis, CCH），作为 PRH 的“续篇与边界”：\textbf{在固定 per-token 推理预算下，表征收敛不蕴含能力收敛}。能力并不收敛到一个点，而是收敛到一个\emph{类}——访问完备混合体（access-complete hybrid）：凡同时持有两类信息通道的架构都落入该类。论文的核心论断可以压缩成一句话：

\begin{keybox}{CCH 核心论断}
\textbf{表征收敛由规模免费赠予；能力收敛必须通过访问结构购买。}智能（能力）不是规模的副产品，而是在固定每 token 预算下、由架构的\textbf{访问结构}（access structure，即模型能以何种通道访问历史信息）所约束的产物。
\end{keybox}

\subsubsection{能力的操作化定义：历史-查询通道}

先厘清两个常被混用的词。\textbf{表征}（representation）回答“看到什么”——模型对数据的内部地图；\textbf{能力}（capability）回答“能做到什么”——在有限的每 token 预算内检索具体事实、执行复合逻辑。CCH 把后者形式化为\textbf{历史-查询通道}（history-to-query channel）：模型读入一段可以无限长的文本流，真正能用于回答查询的只有\textbf{可访问状态} $\Sigma$——一个容量为 $B$ 比特的瓶颈。无论模型总参数多大、表征多“柏拉图”，超出 $\Sigma$ 容量的历史信息都对查询不可见。换言之，\emph{能力不是模型“知道什么”，而是它在预算内“能访问什么、能对访问到的东西计算什么”}——这可以看作对“智能”的一种操作化定义：智能被归结为预算约束下的可回答查询类，而非参数规模或表征质量本身。

\subsubsection{访问结构：双通道完备类}

CCH 断言，固定预算下的能力收敛到“访问完备混合体”这一类：任何同时持有以下两种通道的架构都是完备的；缺任何一条通道，都有某类查询被证明无法完成——

\begin{itemize}
  \item \textbf{压缩态通道}（compressive $O(1)$-state channel）：把整段历史压缩成常数大小的运行状态（SSM/RNN 式递归态、有限维汇总向量），天然擅长状态跟踪与聚合类查询，但保留不下任意久远事件的逐字细节；
  \item \textbf{逐字索引通道}（scalable verbatim-index channel）：可随历史增长的逐字索引（注意力的 KV Cache、外置检索库），擅长精确回查原文，但单独持有它并不解决串行复合与状态聚合。
\end{itemize}

\subsubsection{见证任务与三堵墙}

CCH 把上述论断锚定在一个\textbf{见证任务}上：\textbf{无限流中的牛顿苹果问题}——在无限延长的 token 流中，某个时刻一个“苹果落地”式的具体事件被写入历史；模型既要在任意久之后仍能\textbf{逐字回查}该事件的内容与位置，又要在整条流上维持\textbf{运行状态}（计数、聚合、状态跟踪）。“召回 + 状态跟踪”这一复合任务原语的识别来自 Rawat et al.（2026）的实证。在固定预算且无链式思考（no-CoT）的条件下，任何单一通道架构都会撞上至少一堵\textbf{资源之墙}：

\begin{table}[htbp]
\centering\small
\begin{tabular}{@{}>{\raggedright\arraybackslash}p{1.8cm}>{\raggedright\arraybackslash}p{4.8cm}>{\raggedright\arraybackslash}p{7.2cm}@{}}
\toprule
\textbf{墙} & \textbf{挡住的架构} & \textbf{理论依据}\\
\midrule
香农墙 & 一切 $o(N_b)$ 状态架构：逐字保留长历史不可能 & 数据处理不等式 + Fano（无条件）\\
视界墙 & 一切固定窗口架构：窗外历史不可达 & 窗口外的信息物理上无法访问\\
电路墙 & 固定深度纯注意力栈：串行复合受限 & log-精度 Transformer $\subseteq \mathrm{TC}^0$（条件于 $\mathrm{TC}^0 \neq \mathrm{NC}^1$）\\
\bottomrule
\end{tabular}
\caption{CCH 的三堵资源之墙（Chen et al., 2026）：香农墙来自信息论、无条件成立；电路墙是复杂性条件结论。}
\label{tab:cchwalls}
\end{table}

在显式的\textbf{可分离性假设}下，双通道混合体“按每堵墙的价格付费”（分别支付状态容量、窗口覆盖、组合深度的工程代价）即可逐一穿越三墙——因此\textbf{能力在复合下严格超可加}（strictly super-additive）：两条各有短板的通道组合后，能完成任何单通道都无法完成的查询。这是“$1+1>2$”的架构定理版本，也为第 4 章“RAG 事实层 + SFT 行为层正交可叠加”的生产经验提供了理论注脚。

\textbf{预注册实验}：这篇论文最值得业界学习的是方法论纪律——\textbf{在数据冻结前冻结判据}。三项预测如期命中：\textbf{剪刀差}被精确测得（64 个标量的压缩态在精确回查任务上错误率 0.994，仅增加一层全局注意力后骤降至 0.000）；状态跟踪能力的分岔正好落在预注册边界上；合取见证任务显示出不可约的双通道解。另有一项预测失败且方向反转，作者如实报告。论文同样明确区分\emph{已证明的}与\emph{仍是猜想的}：访问完备性原理基于信息论下界与预注册实验，而“领域级能力收敛趋势”只是经济学动机的猜想——这种自我设限的写法，与 10.1 节的评测方法论（度量伪影、基准污染）互为呼应。

\begin{figure}[htbp]
\centering
\resizebox{0.97\textwidth}{!}{%
% fig14_cch_access_structure.tex — 图14 CCH 双通道访问结构与三堵墙
% 《深度学习与大语言模型：前沿技术全景报告》图示库
% 由 dl_llm_report.tex 抽取，经 % fig14_cch_access_structure.tex — 图14 CCH 双通道访问结构与三堵墙
% 《深度学习与大语言模型：前沿技术全景报告》图示库
% 由 dl_llm_report.tex 抽取，经 % fig14_cch_access_structure.tex — 图14 CCH 双通道访问结构与三堵墙
% 《深度学习与大语言模型：前沿技术全景报告》图示库
% 由 dl_llm_report.tex 抽取，经 \input{figs/fig14_cch_access_structure} 引用（caption/label 在主文件）
% 注意：本文件为片段，依赖主文件导言区的 tikz/pgfplots 宏包、颜色定义与 tikzset 样式，不可单独编译。

\begin{tikzpicture}[x=1cm, y=1cm]
  % ---- 左段：信息源与瓶颈 ----
  \node[blk, align=center] (stream) at (0,0)
    {\textbf{无限 token 流}\\ $\cdots\bullet$ 苹果落地 $\bullet\cdots$\\ {\scriptsize 查询：跟踪 $+$ 回查}};
  \node[orn, align=center] (sigma) at (3.6,0)
    {\textbf{可访问状态} $\Sigma$\\ 瓶颈容量 $B$ 比特\\ {\scriptsize 固定每 token 预算}};
  % ---- 中段：双通道（上下分层） ----
  \node[grn, align=center] (cha) at (7.7,1.3)
    {\textbf{压缩态通道}\\ $O(1)$：SSM/RNN 递归态\\ {\scriptsize 擅长跟踪，弱逐字回查}};
  \node[dgt, align=center] (chb) at (7.7,-1.3)
    {\textbf{逐字索引通道}\\ KV Cache · 外置索引\\ {\scriptsize 擅长回查，弱窗外复合}};
  % ---- 三堵墙：间距 1.15cm，名称水平标于墙顶（虚线不穿任何文字） ----
  \foreach \x/\nm in {10.35/香农墙, 11.5/视界墙, 12.65/电路墙}{
    \draw[dashed, very thick, red!65!black] (\x,-2.35) -- (\x,2.35);
    \node[font=\scriptsize\bfseries, text=red!65!black] at (\x,2.62) {\nm};
  }
  % ---- 右段：完备混合体 ----
  \node[grn, align=center] (hybrid) at (14.9,0)
    {\textbf{访问完备混合体}\\ 双通道，按墙付费\\ $\Rightarrow$ 能力超可加};
  % ---- 连线：双通道斜穿三墙（对应“按墙付费、逐一穿越”） ----
  \draw[brr] (stream) -- (sigma);
  \draw[brr] (sigma.east) -- (cha.west);
  \draw[brr] (sigma.east) -- (chb.west);
  \draw[brr] (cha.east) -- ([yshift=2.5mm]hybrid.west);
  \draw[brr] (chb.east) -- ([yshift=-2.5mm]hybrid.west);
\end{tikzpicture}
 引用（caption/label 在主文件）
% 注意：本文件为片段，依赖主文件导言区的 tikz/pgfplots 宏包、颜色定义与 tikzset 样式，不可单独编译。

\begin{tikzpicture}[x=1cm, y=1cm]
  % ---- 左段：信息源与瓶颈 ----
  \node[blk, align=center] (stream) at (0,0)
    {\textbf{无限 token 流}\\ $\cdots\bullet$ 苹果落地 $\bullet\cdots$\\ {\scriptsize 查询：跟踪 $+$ 回查}};
  \node[orn, align=center] (sigma) at (3.6,0)
    {\textbf{可访问状态} $\Sigma$\\ 瓶颈容量 $B$ 比特\\ {\scriptsize 固定每 token 预算}};
  % ---- 中段：双通道（上下分层） ----
  \node[grn, align=center] (cha) at (7.7,1.3)
    {\textbf{压缩态通道}\\ $O(1)$：SSM/RNN 递归态\\ {\scriptsize 擅长跟踪，弱逐字回查}};
  \node[dgt, align=center] (chb) at (7.7,-1.3)
    {\textbf{逐字索引通道}\\ KV Cache · 外置索引\\ {\scriptsize 擅长回查，弱窗外复合}};
  % ---- 三堵墙：间距 1.15cm，名称水平标于墙顶（虚线不穿任何文字） ----
  \foreach \x/\nm in {10.35/香农墙, 11.5/视界墙, 12.65/电路墙}{
    \draw[dashed, very thick, red!65!black] (\x,-2.35) -- (\x,2.35);
    \node[font=\scriptsize\bfseries, text=red!65!black] at (\x,2.62) {\nm};
  }
  % ---- 右段：完备混合体 ----
  \node[grn, align=center] (hybrid) at (14.9,0)
    {\textbf{访问完备混合体}\\ 双通道，按墙付费\\ $\Rightarrow$ 能力超可加};
  % ---- 连线：双通道斜穿三墙（对应“按墙付费、逐一穿越”） ----
  \draw[brr] (stream) -- (sigma);
  \draw[brr] (sigma.east) -- (cha.west);
  \draw[brr] (sigma.east) -- (chb.west);
  \draw[brr] (cha.east) -- ([yshift=2.5mm]hybrid.west);
  \draw[brr] (chb.east) -- ([yshift=-2.5mm]hybrid.west);
\end{tikzpicture}
 引用（caption/label 在主文件）
% 注意：本文件为片段，依赖主文件导言区的 tikz/pgfplots 宏包、颜色定义与 tikzset 样式，不可单独编译。

\begin{tikzpicture}[x=1cm, y=1cm]
  % ---- 左段：信息源与瓶颈 ----
  \node[blk, align=center] (stream) at (0,0)
    {\textbf{无限 token 流}\\ $\cdots\bullet$ 苹果落地 $\bullet\cdots$\\ {\scriptsize 查询：跟踪 $+$ 回查}};
  \node[orn, align=center] (sigma) at (3.6,0)
    {\textbf{可访问状态} $\Sigma$\\ 瓶颈容量 $B$ 比特\\ {\scriptsize 固定每 token 预算}};
  % ---- 中段：双通道（上下分层） ----
  \node[grn, align=center] (cha) at (7.7,1.3)
    {\textbf{压缩态通道}\\ $O(1)$：SSM/RNN 递归态\\ {\scriptsize 擅长跟踪，弱逐字回查}};
  \node[dgt, align=center] (chb) at (7.7,-1.3)
    {\textbf{逐字索引通道}\\ KV Cache · 外置索引\\ {\scriptsize 擅长回查，弱窗外复合}};
  % ---- 三堵墙：间距 1.15cm，名称水平标于墙顶（虚线不穿任何文字） ----
  \foreach \x/\nm in {10.35/香农墙, 11.5/视界墙, 12.65/电路墙}{
    \draw[dashed, very thick, red!65!black] (\x,-2.35) -- (\x,2.35);
    \node[font=\scriptsize\bfseries, text=red!65!black] at (\x,2.62) {\nm};
  }
  % ---- 右段：完备混合体 ----
  \node[grn, align=center] (hybrid) at (14.9,0)
    {\textbf{访问完备混合体}\\ 双通道，按墙付费\\ $\Rightarrow$ 能力超可加};
  % ---- 连线：双通道斜穿三墙（对应“按墙付费、逐一穿越”） ----
  \draw[brr] (stream) -- (sigma);
  \draw[brr] (sigma.east) -- (cha.west);
  \draw[brr] (sigma.east) -- (chb.west);
  \draw[brr] (cha.east) -- ([yshift=2.5mm]hybrid.west);
  \draw[brr] (chb.east) -- ([yshift=-2.5mm]hybrid.west);
\end{tikzpicture}

}%
\caption{CCH 的双通道访问结构与三堵墙（示意）：单一通道必被至少一堵墙挡住；同时持有压缩态与逐字索引双通道的混合体按每堵墙的价格（状态容量、窗口覆盖、组合深度）逐一穿越，能力在复合下严格超可加（Chen et al., 2026, arXiv:2607.14144）。}
\label{fig:cch}
\end{figure}

\subsubsection{用 CCH 回照本报告：访问结构的工程版图}

CCH 提供了一个把本报告多个章节串起来的统一视角——\emph{过去十年的许多“架构创新”，本质上都是在为这两条通道支付账单或扩展账本}：

\begin{table}[htbp]
\centering\small
\begin{tabular}{@{}lll@{}}
\toprule
\textbf{CCH 概念} & \textbf{工程对应} & \textbf{本报告章节}\\
\midrule
$O(1)$ 压缩态通道 & 线性注意力 / SSM（Mamba）；RNN 遗产 & 第 2、5 章\\
逐字索引通道 & KV Cache 与长上下文管理 & 第 6 章\\
外置逐字索引 & RAG（“事实层”检索增强） & 第 4 章\\
视界墙的工程价格 & 滑动窗口、YaRN 位置外推、长文本课程 & 第 5、6 章\\
电路墙的工程价格 & 链式思考：以测试时计算换组合深度 & 第 7 章\\
访问完备混合体 & SSM $\times$ 注意力混合架构；“RAG $+$ SFT”组合 & 第 4、5 章\\
能力的超可加性 & 事实靠检索、行为靠微调，两者正交可叠加 & 第 4 章\\
\bottomrule
\end{tabular}
\caption{用 CCH 回照本报告：访问结构概念的工程对应}
\label{tab:cchmap}
\end{table}

这张表也解释了 2025--2026 年架构竞争的真实格局：\textbf{纯 SSM、纯注意力、纯检索都不是终点，“混合 + 外置”才是通往访问完备的工程路线}——Jamba 类混合架构、长上下文与 RAG 并用的生产组合、把推理深度转移到测试时的推理模型（第 7 章），都是在同一本“访问结构账”上的不同记账方式。最后需要强调边界：CCH 的实验仍是小规模预注册测试，“领域级收敛”是猜想而非定理，且其结论限于无 CoT 的固定预算情形——但“能力由访问结构定价”这一视角，对理解第 11 章的算力--资本--生态格局与第 12 章的趋势判断，已经足够有力。

\subsection{未来方向（2025--2030）}

\begin{table}[htbp]
\centering
\small
\begin{tabular}{@{}ll@{}}
\toprule
\textbf{方向} & \textbf{描述}\\
\midrule
推理为王 & 从“更大模型”转向“更聪明推理”（Test-time Compute）\\
Agent 化 & LLM 从“回答者”变为“行动者”（工具调用、多步规划）\\
多模态统一 & 文本+图像+音频+视频+3D 统一在一个模型\\
效率革命 & MoE + 量化 + 推测解码，100B+ 可在消费级硬件运行\\
访问结构工程 & 混合架构与外置记忆：按 CCH 补齐能力通道（10.3 节）\\
AI for Science & 蛋白质设计、材料发现、数学定理证明\\
端侧 AI & 手机/PC 上的 1B--10B 模型成为标配\\
Agent 协议（MCP） & Agent 连接工具/数据的事实标准（2026）\\
具身智能 & LLM + 机器人，操作物理世界\\
世界模型 & AI 内部模拟物理/社会环境\\
\bottomrule
\end{tabular}
\caption{2025--2030 关键趋势}
\end{table}

其中\textbf{世界模型}与\textbf{具身智能}值得单独展开。世界模型路线（LeCun 的 JEPA 主张、Genie 3 的可交互世界生成、视频生成模型中被广泛讨论的“隐式物理理解”）的核心命题是：\emph{语言是对世界的有损压缩，先建立物理直觉再谈智能}——这与“纯语言训练触顶”的判断互为印证。具身智能则把 LLM 当“大脑”（任务规划与语义理解），把 VLA（视觉-语言-动作）小模型当“小脑”（低层控制），2025 年中国具身智能赛道 73 笔、257 亿元的融资热度与 2026 年“世界模型元年”的说法互为表里——两条线最终会在“能动手的智能体”上汇合。

% ===================== 第十一章 =====================
\section{2025--2026 行业前沿动态}

本章汇总 2025 年至 2026 年上半年的行业关键进展。\textbf{阅读提示}：本章内容主要来自厂商公告与行业公开报道，个别 2026 年事件（如 NVIDIA 收购 Hugging Face、Gemini 3.7、Google I/O 2026 主题等）的量化细节（金额、具体型号参数）以最终官方披露为准；引用这些数据用于正式文档前，建议回溯一手来源核实。

\textbf{本章三条主线}：(1)\textbf{推理与 Agent 成为产品主形态}——o1$\to$o3、R1、Gemini thinking 系列把“思考预算”做成产品开关，Operator/Computer Use 把“执行动作”做成标准能力；(2)\textbf{开源与闭源差距收敛}——Qwen3、Llama 4、DeepSeek 系列在核心基准上逼近同期闭源旗舰，开源阵营首次具备“生产级推理模型”选项；(3)\textbf{算力与生态整合加速}——NVIDIA 的纵向整合（芯片$\to$框架$\to$生态）、云厂商的自研芯片（TPU v6、Trainium 2/3）与数据中心扩张构成基础设施主战场。

\begin{table}[htbp]
\centering
\small
\begin{tabular}{@{}lp{6.2cm}p{4.4cm}@{}}
\toprule
\textbf{公司/项目} & \textbf{最新进展} & \textbf{亮点}\\
\midrule
OpenAI & o3/o4-mini（2025-04）、GPT-5（08）、GPT-5.2（12）、Sora 2、Operator & 推理 + Agent + 多模态统一\\
Google & Gemini 2.5 $\to$ Gemini 3 Pro（2025-11，LMArena 1501）、Veo 3/3.1、Genie 3 & 原生多模态 + 世界模型\\
Anthropic & Claude 4（05）$\to$ Opus 4.5（11，SWE-bench 80.9\%）、Computer Use & 编程与 Agent 执行标杆\\
Meta & Llama 4（2025-04，Scout/Maverick） & 大规模 MoE 开源（评测有争议）\\
xAI & Grok 3（02）$\to$ Grok 4（07，多智能体） & 10 万卡 Colossus 算力\\
DeepSeek & R1（01）、V3.1（08） & MLA、纯 RL 推理、RL 成本 29.4 万美元\\
阿里 & Qwen3（04）、Qwen3-Coder（07） & 36T token 预训练、开源矩阵最全\\
月之暗面 & Kimi K2（07） & 1T 级 MoE 全量开源，开源榜第一\\
智谱 & GLM-4.5（07）、2026-01 港股上市 & 大模型第一股\\
NVIDIA & B200/GB200、约 \$12.9B 收购 Hugging Face（2026-08） & 算力 + 生态纵向整合\\
\bottomrule
\end{tabular}
\caption{2025--2026 主要玩家动态}
\end{table}

\textbf{行业关键事件}：
\begin{itemize}
  \item NVIDIA 以近 \$130 亿收购 Hugging Face，算力+开源生态深度绑定
  \item OpenAI、Anthropic、Google 联合发布 AI 网络安全防御公开信（100+ 公司签署）
  \item Suno（AI 音乐生成）面临多位音乐人的风格模仿诉讼
  \item Adobe 工具深度集成 Slack，AI 辅助设计工作流成熟
  \item 大型数据中心环保争议（天然气发电、联邦法律违规指控）
  \item Google I/O 2026 宣布“Agentic Gemini 时代”
  \item ChatGPT 推出临时聊天、个性化、插件等持续迭代
\end{itemize}

\subsection{2025 年时间线：从 DeepSeek 冲击波到多极格局}

\begin{itemize}
  \item \textbf{1 月}：DeepSeek-R1 发布（对标 o1），数日内登顶中美应用商店，引发 NVIDIA 单日市值蒸发约 5890 亿美元（跌幅约 17\%）——DeepSeek 冲击波成为全年行业叙事的分水岭；OpenAI 发布 Operator（CUA 模型），打响 Agent 元年第一枪；
  \item \textbf{2 月}：xAI Grok 3 发布（10 万张 H100 的 Colossus 集群训练）；EU AI Act 禁止性条款开始适用；
  \item \textbf{3 月}：通用 Agent Manus 出圈（GAIA 基准超越 OpenAI DeepResearch）；GPT-4o 原生图像生成上线，一周内 1.3 亿用户生成超 7 亿张图；
  \item \textbf{4 月}：o3 / o4-mini 发布（推理模型首次全工具链调用）；Llama 4（Scout/Maverick，MoE）发布但评测口径陷入争议；Qwen3 发布（36T token、119 种语言、混合思考）；
  \item \textbf{5 月}：Claude 4 系列（Opus 4 / Sonnet 4）确立编程标杆；Google I/O 发布 Veo 3（首个音画同步视频生成）；
  \item \textbf{6--7 月}：o3-pro；\textbf{Kimi K2} 万亿参数全量开源（LMArena 开源第一）；Qwen3-Coder（Agentic Coding 开源 SOTA）；GLM-4.5（355B MoE）开源；Grok 4（多智能体 Heavy）；
  \item \textbf{8 月}：\textbf{GPT-5} 发布（快慢思考统一路由，幻觉率较 o3 大幅下降）；OpenAI 时隔多年重回开源（GPT-oss-120B/20B）；DeepSeek-V3.1（混合思考 + 强代码）；Google Genie 3 世界模型；
  \item \textbf{9 月}：Sora 2 + 独立 App（被称作视频界的 GPT-3.5 时刻）；中国《人工智能生成合成内容标识办法》施行；DeepSeek 经 Nature 封面论文披露 R1 训练成本仅 29.4 万美元；
  \item \textbf{10--12 月}：Claude Opus 4.5（SWE-bench 80.9\%，首个破 80）；\textbf{Gemini 3 Pro}（LMArena 首破 1500 Elo，GPQA 91.9\%）；Veo 3.1；GPT-5.1 / 5.2；据报道 Manus 被 Meta 收购；
  \item \textbf{2026 年 1 月}：智谱 AI 港股上市（全球大模型第一股）、MiniMax 上市首日涨 109\%；LeCun 离开 Meta 创立 AMI Labs 押注世界模型；三大实验室同步进军医疗 AI。
\end{itemize}

\subsection{中国力量：从追赶到并跑}

\textbf{开源阵营的主导权已明显东移}：2025 年全球开源模型的下载量与衍生生态中，中国模型占据多数份额；在生产级推理与 Agent 能力上，开源首次与闭源并肩。代表阵营：
\begin{itemize}
  \item \textbf{DeepSeek}：R1（RL 成本 29.4 万美元）$\to$ V3.1（混合思考、强代码）；以 MLA / MTP / FP8 / 无辅助损失 MoE 的架构-系统协同路线持续压低训练与推理成本，是全年行业成本曲线的锚点；
  \item \textbf{阿里 Qwen}：Qwen3 全系开源（0.6B--235B，Apache 2.0），36T token 预训练、119 种语言，衍生模型生态为全球最大；Qwen3-Coder 补齐 Agent 编程短板；
  \item \textbf{月之暗面 Kimi}：K2 以 1T 级 MoE 全量开源并登顶开源榜，验证 Agentic 数据工程 + 大 MoE 路线；
  \item \textbf{智谱 GLM}：GLM-4.5（355B MoE）开源，2026 年 1 月率先登陆港交所成为全球大模型第一股；
  \item \textbf{字节跳动}：豆包系产品用户量居国内首位，Seedream 4.0（图像生成/编辑一体化）与 Seedance 1.0（视频生成）在多模态侧形成竞争力；
  \item \textbf{MiniMax}：M1（长上下文混合架构）与海螺视频，2026 年 1 月上市首日市值破千亿港元。
\end{itemize}
产业与资本面：2025 年 1--10 月中国 AI 领域超亿元融资 139 笔、总额超 600 亿元人民币，其中具身智能 73 笔、257 亿元；监管侧备案制 + 内容标识双轨落地（见第 9 章），发展与规范并重。

\subsection{算力、资本与开源生态}

\textbf{算力}：NVIDIA Blackwell（B200/GB200）进入规模部署，2026 年 8 月官宣以约 \textbf{129 亿美元}收购 Hugging Face（含约 119 亿美元现金，进入 HSR 反垄断审查）——芯片 + 框架 + 模型生态的纵向整合引发开源社区对中立性的担忧；云厂商自研芯片（TPU、Trainium、MTIA）继续分流；吉瓦级数据中心的能源与环评（燃气发电、核电 PPA）成为新的瓶颈议题。\textbf{资本}：OpenAI 完成 400 亿美元融资（估值 3000 亿美元）、2026 年营收目标 300 亿美元；Anthropic F 轮 130 亿美元（估值 1830 亿美元）；行业在 AI 泡沫争议中进入以真实 ROI 论成败的价值证明之年。\textbf{开源生态}：MCP 协议捐赠开源基金会、成为 Agent 连接工具的事实标准；OpenAI 重回开源（GPT-oss）、Llama 4 争议动摇 Meta 的开源信誉、Qwen/Kimi/GLM 补位——开源重心从单点旗舰转向\emph{全栈生态}（权重 + 推理引擎 + Agent 框架 + 协议）。


% ===================== 第十二章 =====================
\section{总结与展望}

\subsection{技术栈全景}

\[
\boxed{\parbox{0.85\textwidth}{\centering
\textbf{应用层}：对话 · 代码 · Agent · 多模态 · 科学发现\\[2pt]
\textbf{协议层}：MCP 工具协议 · Function Calling · Agent 框架\\[2pt]
\textbf{模型层}：LLM · 推理模型 · 多模态 · MoE · 世界模型\\[2pt]
\textbf{训练层}：预训练 $\to$ SFT $\to$ RLHF/DPO/RL $\to$ 持续学习\\[2pt]
\textbf{基础设施}：GPU/TPU · 分布式训练 · 推理引擎 · 量化 · 边缘\\[2pt]
\textbf{安全治理}：对齐 · 红队 · 评估 · 内容标识 · 法规
}}
\]

图~\ref{fig:flywheel} 从系统视角把六层串成一个闭环：数据经训练变成模型，模型经部署触达应用，应用产生交互数据与安全信号回流——\emph{大模型竞争的本质是飞轮转速之争}。

\begin{figure}[htbp]
\centering
\input{figs/fig15_data_flywheel}
\caption{大模型技术栈的运行闭环：每一轮迭代都在为下一轮积累数据与评估资产。飞轮各环节分别对应第 4 章（训练）、第 5--6 章（模型与部署）、第 8 章（应用形态）与第 9 章（评测与安全）。}
\label{fig:flywheel}
\end{figure}

最后给一个\textbf{三遍读法}的建议：第一遍只读每章的 keybox/infobox 与全部图示（约 45 分钟），建立全景地图；第二遍带着手头的具体问题精读对应章节（本报告刻意让每章可独立成篇）；第三遍顺着“关键论文与延伸阅读”读原始文献——技术报告的价值在于把问题“问到点子上”，而答案永远在原始来源里持续更新。

\subsection{核心趋势}

\textbf{趋势判断的方法论}：以上趋势并非并列的“新方向”，而是沿\emph{算力分配}这条主线展开——训练算力（参数/数据扩展）边际收益递减后，算力向推理侧（test-time compute）、向架构效率（MoE/MLA/SSM）、向端侧（量化/小模型）重新配置；“安全内生”与“Agent 化”则是这一重新配置在治理与应用层的投影。理解这一主线，就能从趋势列表中推演出尚未出现的产品形态。

\begin{enumerate}
  \item \textbf{推理为王}：Test-time Compute 成为新扩展维度（思考预算已成产品开关）
  \item \textbf{Agent 化}：LLM 从“回答者”变为“行动者”，MCP 等协议成为互操作底座
  \item \textbf{多模态统一}：单一模型处理文本、图像、音频与视频（理解 + 生成）
  \item \textbf{效率革命}：MoE + 量化 + 推测解码，100B+ 模型可在消费级硬件运行
  \item \textbf{访问结构经济学}：能力由“压缩态 + 逐字索引”双通道的完备性定价，混合架构与外置记忆成为能力竞争的新前线（见 10.3 节）
  \item \textbf{AI for Science}：从工具到发现者（蛋白质、材料、数学）
  \item \textbf{安全内生}：从“事后过滤”到“训练时对齐 + 运营时审计”
  \item \textbf{开源东移}：Qwen、DeepSeek、Kimi、GLM 与 Llama 并立，开源首次覆盖生产级推理与 Agent 能力
  \item \textbf{端侧 AI}：隐私优先的本地推理（1B--10B 模型进手机/PC）
  \item \textbf{世界模型}：从语言走向物理（Genie 3、AMI Labs），为具身智能铺路
\end{enumerate}

2026 年值得盯住的三个变量：一是\emph{世界模型}能否从演示走向可用（Genie 3 类可交互世界、AMI Labs 路线）；二是\emph{上市潮}能否经受盈利检验（头部的 IPO 预期与港交所已上市公司的商业化数据）；三是\emph{监管全面生效}（EU AI Act 2026 年 8 月全面适用、中国标识义务常态化）对产品形态的实际约束。

\subsection{给学习者的建议}

\begin{itemize}
  \item \textbf{基础}：线性代数、概率论、优化理论、Python
  \item \textbf{框架}：PyTorch 精通 $\to$ JAX（研究前沿）
  \item \textbf{实践}（由浅入深）：
  \begin{itemize}
    \item 从 0 复现 Transformer（含 KV Cache 与 FlashAttention 思想）
    \item 训练 1B 小模型（预训练 $\to$ SFT），走完一条完整数据管线
    \item 部署 vLLM/SGLang 推理服务，动手实践量化与推测解码
    \item 复现 DPO 与 R1 风格 GRPO（规则奖励 + 可验证任务）
    \item 用 LoRA/QLoRA 做领域适配，并构建一个 MCP 工具调用 Agent
  \end{itemize}
  \item \textbf{前沿追踪}：arXiv (cs.LG, cs.CL)、HuggingFace、顶会 (NeurIPS/ICML/ICLR/ACL)；中文社区：机器之心、量子位、智源社区、ModelScope 与 GitHub 中文区
  \item \textbf{工具链}：Transformers + \allowbreak{}DeepSpeed/Megatron-LM + \allowbreak{}TRL/LLaMA-Factory + \allowbreak{}vLLM/SGLang + \allowbreak{}W\&B
\end{itemize}

\begin{infobox}{最后的话}
大模型技术栈的特点是\textbf{半衰期短}：本报告中的具体模型与数字（参数、成本、基准）会以季度为单位被刷新，但\emph{架构原理}（注意力、稀疏化、对齐）、\emph{系统原则}（并行切分、显存账本、算术强度）与\emph{方法论}（Scaling Law、控制变量、防御纵深）会以年为单位稳定。把时间投资在后两者上，是应对技术快速迭代的最优策略。
\end{infobox}

% ===================== 附录 =====================
\newpage
\appendix
\section{关键术语速查表}

\begin{longtable}{@{}>{\raggedright\arraybackslash}p{3.6cm}>{\raggedright\arraybackslash}p{4.6cm}p{6.8cm}@{}}
\toprule
\textbf{术语} & \textbf{英文} & \textbf{一句话解释}\\
\midrule
\endfirsthead
\multicolumn{3}{@{}l}{\small\itshape （续上页）}\\
\toprule
\textbf{术语} & \textbf{英文} & \textbf{一句话解释}\\
\midrule
\endhead
\midrule
\multicolumn{3}{r@{}}{\small\itshape （接下页）}\\
\endfoot
\bottomrule
\endlastfoot
注意力 & Attention & 让模型动态决定“看哪里”的机制\\
多头注意力 & Multi-Head Attention & 多组并行注意力捕捉不同关系\\
因果掩码 & Causal Mask & 确保只能看到前文，实现自回归\\
位置编码 & Positional Encoding & 注入 token 位置信息\\
RoPE & Rotary Position Embedding & 旋转矩阵编码位置，支持外推\\
MoE & Mixture of Experts & 稀疏激活专家网络\\
GQA & Grouped Query Attention & 分组 KV 头，减少 Cache\\
MLA & Multi-head Latent Attention & 联合压缩 KV，极致省显存\\
KV Cache & Key-Value Cache & 缓存已算 K/V 避免重算\\
FlashAttention & FlashAttention & IO 感知的分块注意力算法\\
推测解码 & Speculative Decoding & 小模型草拟 + 大模型验证\\
RLHF & RL from Human Feedback & 人类反馈强化学习对齐\\
DPO & Direct Preference Optimization & 无 RM 的偏好优化\\
CoT & Chain-of-Thought & 思维链推理\\
PRM & Process Reward Model & 过程奖励模型\\
RAG & Retrieval Augmented Generation & 检索增强生成\\
SFT & Supervised Fine-Tuning & 监督微调\\
LoRA & Low-Rank Adaptation & 低秩适配微调\\
量化 & Quantization & 降低数值精度省资源\\
连续批处理 & Continuous Batching & 动态批处理提升吞吐\\
数据墙 & Data Wall & 高质量数据枯竭\\
能力收敛假说 & Capability Convergence Hypothesis & 能力由访问结构定价，而非规模\\
访问结构 & Access Structure & 压缩态 + 逐字索引双通道的完备类\\
AGI & Artificial General Intelligence & 通用人工智能\\
具身智能 & Embodied AI & AI 与物理世界交互\\
世界模型 & World Model & AI 内部模拟环境\\
对齐 & Alignment & 让模型行为符合人类意图\\
红队 & Red Teaming & 对抗性安全测试\\
GRPO & Group Relative Policy Optimization & 组相对策略优化（免 Critic）\\
SSM & State Space Model & 状态空间模型（Mamba）\\
张量并行 & Tensor Parallelism & 单矩阵切多卡，NVLink 域内\\
流水线并行 & Pipeline Parallelism & 不同层放不同卡，跨节点\\
ZeRO & Zero Redundancy Optimizer & 分片优化器状态省显存\\
算术强度 & Arithmetic Intensity & FLOPs/Byte，判计算/访存受限\\
PagedAttention & Paged KV Cache & KV 分页管理，消碎片\\
MCP & Model Context Protocol & Agent 连接工具/数据的标准协议（2026 事实标准）\\
推测解码变体 & EAGLE / Medusa & 特征级 / 多头自草拟\\
MTP & Multi-Token Prediction & 多 token 预测，兼作草拟器\\
YaRN & Yet another RoPE extensioN & RoPE 长上下文分段外推\\
QLoRA & Quantized LoRA & 4-bit 基座 + LoRA 微调\\
NF4 & NormalFloat 4-bit & 正态感知 4-bit 量化格式\\
GPTQ / AWQ & — & 离线最优 / 激活感知 INT4 量化\\
SmoothQuant & — & 激活-权重量化难度迁移（W8A8）\\
RLAIF & RL from AI Feedback & AI 反馈替代人类反馈\\
越狱 & Jailbreak & 诱导模型违反安全约束的提示攻击\\
提示注入 & Prompt Injection & 经外部内容注入恶意指令\\
Test-time Compute & 测试时计算 & 用推理时长换准确率\\
MCTS & Monte Carlo Tree Search & 蒙特卡洛树搜索（推理路径）\\
CoT 忠实度 & CoT Faithfulness & 思考链是否反映真实决策依据\\
VAD & Voice Activity Detection & 语音端点检测\\
EnCodec & — & 神经声码编解码器（语音 token 化）\\
DiT & Diffusion Transformer & 扩散 Transformer 生成骨干\\
3D VAE & — & 视频时空潜空间压缩\\
GUI Agent & — & 屏幕截图理解 + 操作规划执行\\
数据配比 & Data Mixing & 各来源数据比例设计\\
灾难性遗忘 & Catastrophic Forgetting & 新训练覆盖旧知识\\
困惑度 & Perplexity & 平均“候选词犹豫数”，$\exp$ 损失\\
MMLU / C-Eval & — & 知识类通用基准（英/中）\\
GSM8K/MATH & — & 数学推理基准\\
HumanEval/MBPP & — & 代码能力基准\\
\end{longtable}

\newpage
\section{关键论文与延伸阅读}

\sloppy  % 参考文献条目含长 arXiv 编号，允许更宽松的断行以避免溢出版心

以下按主题列出本报告涉及的核心文献（均可经 arXiv 或官方博客获取），供延伸阅读与数据溯源：

\textbf{架构基础}
\begin{itemize}
  \item Vaswani et al., 2017. \emph{Attention Is All You Need}（arXiv:1706.03762）——Transformer 原始论文；
  \item He et al., 2015. \emph{Deep Residual Learning}（arXiv:1512.03385）——残差连接，被 Transformer 直接继承；
  \item Su et al., 2021. \emph{RoFormer: Rotary Position Embedding}（arXiv:2104.09864）——RoPE；
  \item Peng et al., 2023. \emph{YaRN}（arXiv:2309.00071）——RoPE 长上下文扩展；
  \item Gu and Dao, 2023. \emph{Mamba}（arXiv:2312.00752）——选择性状态空间模型。
\end{itemize}

\textbf{Scaling Law}
\begin{itemize}
  \item Kaplan et al., 2020. \emph{Scaling Laws for Neural Language Models}（arXiv:2001.08361）；
  \item Hoffmann et al., 2022. \emph{Training Compute-Optimal LLMs}（Chinchilla，arXiv:2203.15556）；
  \item Snell et al., 2024. \emph{Scaling LLM Test-Time Compute Optimally}（arXiv:2408.03314）——推理时扩展的开山之作；
  \item Chen et al., 2026. \emph{Capability from Access Structure, Not Scale: Lower Bounds and Pre-Registered Tests for Hybrid Sequence Models}（CCH，arXiv:2607.14144）——10.3 节的理论来源；
  \item Huh et al., 2024. \emph{The Platonic Representation Hypothesis}（arXiv:2405.07987）——CCH 所回应的表征收敛假说；
  \item Merrill and Sabharwal, 2023. \emph{The Parallelism Tradeoff: Limitations of Log-Precision Transformers}（arXiv:2207.00729）——“电路墙”的理论基础。
\end{itemize}

\textbf{MoE 与注意力变体}
\begin{itemize}
  \item Fedus et al., 2021. \emph{Switch Transformers}（arXiv:2101.03961）；
  \item Ainslie et al., 2023. \emph{GQA}（arXiv:2305.13245）；
  \item Jiang et al., 2024. \emph{Mixtral of Experts}（arXiv:2401.04088）；
  \item DeepSeek-AI, 2024. \emph{DeepSeek-V2}（arXiv:2405.04434，MLA）与 \emph{DeepSeek-V3}（arXiv:2412.19437，无辅助损失 MoE / MTP / FP8）。
\end{itemize}

\textbf{训练与对齐}
\begin{itemize}
  \item Loshchilov and Hutter, 2019. \emph{Decoupled Weight Decay Regularization}（AdamW，arXiv:1711.05101）；
  \item Hu et al., 2021. \emph{LoRA}（arXiv:2106.09685）；Dettmers et al., 2023. \emph{QLoRA}（arXiv:2305.14314）；
  \item Ouyang et al., 2022. \emph{InstructGPT}（arXiv:2203.02155）——RLHF 三部曲；
  \item Bai et al., 2022. \emph{Constitutional AI}（arXiv:2212.08073）；
  \item Rafailov et al., 2023. \emph{DPO}（arXiv:2305.18290）；
  \item Shao et al., 2024. \emph{DeepSeekMath}（GRPO，arXiv:2402.03300）。
\end{itemize}

\textbf{推理系统与推理模型}
\begin{itemize}
  \item Dao et al., 2022. \emph{FlashAttention}（arXiv:2205.14135）；
  \item Kwon et al., 2023. \emph{vLLM: PagedAttention}（arXiv:2309.06180）；
  \item Leviathan et al., 2023. \emph{Speculative Decoding}（arXiv:2211.17192）；Li et al., 2024. \emph{EAGLE}（arXiv:2401.15077）；
  \item Wei et al., 2022. \emph{Chain-of-Thought}（arXiv:2201.11903）；Wang et al., 2022. \emph{Self-Consistency}（arXiv:2203.11171）；
  \item Lightman et al., 2023. \emph{Let's Verify Step by Step}（PRM，arXiv:2305.20050）；
  \item DeepSeek-AI, 2025. \emph{DeepSeek-R1}（arXiv:2501.12948，Nature 封面论文）；
  \item Qwen Team, 2025. \emph{Qwen3 Technical Report}（arXiv:2505.09388）。
\end{itemize}

\textbf{模型体系与多模态}
\begin{itemize}
  \item Brown et al., 2020. \emph{GPT-3}（arXiv:2005.14165）；OpenAI, 2023. \emph{GPT-4 Technical Report}（arXiv:2303.08774）；
  \item Touvron et al., 2023. \emph{LLaMA 2}（arXiv:2307.09288）；Grattafiori et al., 2024. \emph{The Llama 3 Herd of Models}（arXiv:2407.21783）；
  \item Radford et al., 2021. \emph{CLIP}（arXiv:2103.00020）；Radford et al., 2022. \emph{Whisper}（arXiv:2212.04356）；
  \item Liu et al., 2023. \emph{LLaVA}（arXiv:2304.08485）；
  \item Peebles and Xie, 2023. \emph{Scalable Diffusion Models with Transformers (DiT)}（arXiv:2212.09748）。
\end{itemize}

\textbf{行业动态溯源}：各模型官方技术报告/系统卡（OpenAI、Anthropic、Google、Meta、DeepSeek、Qwen、Moonshot、智谱）；中文行业报道推荐机器之心、量子位与 36 氪 AI 频道；监管文本以网信办、欧盟委员会官网为准。

\vfill
\begin{center}
\large\textbf{--- 报告完 ---}\\[0.5cm]
{\small\itshape 本报告基于公开技术报告、学术论文和行业动态整理。\ 如需深入任何特定章节，可进一步展开。}
\end{center}

\end{document}

